% Velocity and Acceleration as Vectors — Pure Mathematics with Mechanics Upper Sixth — Cours signé OnBuch+
% Official MINESEC syllabus. Compiler avec : tectonic PMM10-velocity-and-acceleration-as-vectors.tex
\newcommand{\DOCMATIERE}{Pure Mathematics with Mechanics}
\newcommand{\DOCNIVEAU}{Upper Sixth}
\newcommand{\DOCMODULE}{Upper Sixth — Pure mathematics with mechanics}
\newcommand{\DOCLECON}{10}
\newcommand{\DOCTITRE}{Velocity and Acceleration as Vectors}
% OnBuch+ — préambule « cours signés » ENRICHI (compilé avec tectonic / XeLaTeX)
% Système visuel « Pop » d'OnBuch+ : fond crème (jamais blanc), bordures encre
% épaisses, ombres « dures » décalées, titres Archivo Black en capitales.
% Version MATHÉMATIQUES : graphes (pgfplots/tikz), boîtes pédagogiques
% (définition, propriété, méthode, attention, le savais-tu, exercice résolu,
% exercice, TP, bilan), page de garde et signature OnBuch+ ancrée en bas.
%
% Macros attendues dans main.tex AVANT \input{preamble} :
%   \DOCMATIERE  \DOCNIVEAU  \DOCMODULE  \DOCLECON (numéro)  \DOCTITRE
\documentclass[11pt]{article}

\usepackage{fontspec}
\usepackage[english]{babel}
\usepackage[a4paper,margin=2cm,top=3.4cm,bottom=2.8cm,headheight=1.4cm,footskip=1.3cm]{geometry}
\usepackage{amsmath,amssymb,amsfonts}
\usepackage{mathtools} % \xmapsto, \coloneqq... souvent utilisés par les modèles
\usepackage{cancel} % simplifications barrées (bilans de forces)
% \abs{z} (module, valeur absolue) et \norm{u} (norme), utilisés par les modèles
\providecommand{\abs}{}\let\abs\relax\DeclarePairedDelimiter{\abs}{\lvert}{\rvert}
\providecommand{\norm}{}\let\norm\relax\DeclarePairedDelimiter{\norm}{\lVert}{\rVert}
\usepackage[table]{xcolor} % [table] : fournit aussi \rowcolors (alternance de lignes)
\usepackage{pagecolor}
\usepackage{tikz}
\usetikzlibrary{arrows.meta,positioning,calc,shapes.geometric,shapes.misc,shapes.arrows,decorations.pathmorphing,decorations.pathreplacing,decorations.markings,patterns,shadows,mindmap,trees,fit,backgrounds,matrix,intersections,angles,quotes}
\usepackage{pgfplots}
\usepackage[version=4]{mhchem} % équations nucléaires \ce{^{235}_{92}U}
\usepackage[european,siunitx]{circuitikz} % schémas électriques
\pgfplotsset{compat=1.18}
\usepgfplotslibrary{fillbetween,groupplots} % aires entre courbes (calcul intégral)
\usepackage{enumitem}
\usepackage{fancyhdr}
% [most] charge toutes les bibliothèques tcolorbox (listings, minted, raster,
% documentation...) et gonfle inutilement la mémoire TeX sur un document long
% (~300 boîtes) : on ne charge que ce qui est réellement utilisé.
\usepackage[skins,breakable]{tcolorbox}
\usepackage{graphicx}
\usepackage{colortbl}
\usepackage{booktabs}
\usepackage{tabularx}
\usepackage{diagbox} % case d'en-tête diagonale des tableaux à double entrée
% Colonnes extensibles centrée / gauche / droite (C, L, R), souvent utilisées par les modèles
\newcolumntype{C}{>{\centering\arraybackslash}X}
\newcolumntype{L}{>{\raggedright\arraybackslash}X}
\newcolumntype{R}{>{\raggedleft\arraybackslash}X}
\usepackage{array}
\usepackage{multicol}
\usepackage{siunitx}
% parse-numbers=false : accepte aussi \SI{2,5 \times 10^{-3}}{...}
\sisetup{locale=UK,output-decimal-marker={.},group-minimum-digits=4,parse-numbers=false}
\DeclareSIUnit\ohm{\ensuremath{\symup{\Omega}}} % U+2126 absent de la police
\DeclareSIUnit\division{div} % réglages d'oscilloscope (ms/div, V/div)
\DeclareSIUnit\tr{tr} % tours (vitesse de rotation, ex. \SI{1500}{\tr\per\minute})
\usepackage{qrcode}
% \degree : commande fréquemment utilisée par les modèles pour un angle ou une
% température en dehors de \SI{}{} (non fournie par les paquets chargés).
\providecommand{\degree}{^{\circ}}
% \qed (fin de démonstration), utilisé par les modèles sans amsthm
\providecommand{\qed}{\hfill$\square$}
% \barre : double barre des valeurs interdites dans un tableau de variations (tkz-tab)
\providecommand{\barre}{\ensuremath{\|}}
% environnement proof (amsthm non chargé) utilisé par les modèles
\ifdefined\proof\else\newenvironment{proof}[1][Proof]{\par\noindent\textit{#1.}\ }{\qed\par}\fi
\usepackage{lastpage}
\usepackage{hyperref}
\hypersetup{hidelinks}

% ============================================================
% Palette « Pop » OnBuch+ — mêmes hex que la classe P (plus_ui.dart)
% ============================================================
\definecolor{poporange}{HTML}{F59321}
\definecolor{popdark}{HTML}{E07A0C}
\definecolor{popcream}{HTML}{FFF6E8}
\definecolor{popink}{HTML}{1C1714}
\definecolor{poppurple}{HTML}{7A5AE0}
\definecolor{popgreen}{HTML}{2E9E6A}
\definecolor{poppink}{HTML}{E0457B}
\definecolor{popblue}{HTML}{2F7DE1}
\definecolor{popgold}{HTML}{C98A12}
\definecolor{popmuted}{HTML}{8A7F76}
\definecolor{popline}{HTML}{EADFD2}
\definecolor{popgray}{HTML}{8A7F76}
\colorlet{popgrey}{popgray}
% Alias tolérants (le modèle nomme parfois les couleurs différemment)
\colorlet{popwhite}{white}
\colorlet{popblack}{popink}
\colorlet{popred}{poppink}
\colorlet{popyellow}{popgold}
% Teintes claires pour fonds de tableaux / zones
\colorlet{popblueL}{popblue!10!white}
\colorlet{poporangeL}{poporange!14!white}
\colorlet{popgreenL}{popgreen!12!white}
\colorlet{poppurpleL}{poppurple!12!white}
\colorlet{poppinkL}{poppink!12!white}
\colorlet{popgoldL}{popgold!14!white}

\pagecolor{popcream}

% ============================================================
% Typographie
% ============================================================
\defaultfontfeatures{Ligatures=TeX}
\setmainfont{PlusJakartaSans-Medium.ttf}[Path=../../../fonts/,BoldFont=PlusJakartaSans-Bold.ttf]
% Maths en sans-serif (Fira Math) avec chiffres et lettres droites de Plus
% Jakarta, pour que les formules se fondent dans le texte.
\usepackage[math-style=ISO,bold-style=ISO]{unicode-math}
\setmathfont{FiraMath-Regular.otf}[Path=../../../fonts/]
\setmathfont{PlusJakartaSans-Medium.ttf}[Path=../../../fonts/,range={up/num,up/{Latin,latin},\mathup}]
% (icomma retiré : décimales avec point en anglais)
% Caractères mathématiques Unicode absents des polices de texte (Plus Jakarta,
% Archivo Black) : rendus par leur équivalent mathématique, en texte comme en
% formule — sinon ils s'affichent en carré vide (ex. « Arithmétique dans ℤ »).
\usepackage{newunicodechar}
\newunicodechar{ℝ}{\ensuremath{\mathbb{R}}}
\newunicodechar{ℤ}{\ensuremath{\mathbb{Z}}}
\newunicodechar{ℂ}{\ensuremath{\mathbb{C}}}
\newunicodechar{ℕ}{\ensuremath{\mathbb{N}}}
\newunicodechar{ℚ}{\ensuremath{\mathbb{Q}}}
\newunicodechar{𝔻}{\ensuremath{\mathbb{D}}}
\newunicodechar{α}{\ensuremath{\alpha}}
\newunicodechar{β}{\ensuremath{\beta}}
\newunicodechar{γ}{\ensuremath{\gamma}}
\newunicodechar{δ}{\ensuremath{\delta}}
\newunicodechar{ε}{\ensuremath{\varepsilon}}
\newunicodechar{θ}{\ensuremath{\theta}}
\newunicodechar{λ}{\ensuremath{\lambda}}
\newunicodechar{μ}{\ensuremath{\mu}}
\newunicodechar{σ}{\ensuremath{\sigma}}
\newunicodechar{τ}{\ensuremath{\tau}}
\newunicodechar{φ}{\ensuremath{\varphi}}
\newunicodechar{ψ}{\ensuremath{\psi}}
\newunicodechar{ω}{\ensuremath{\omega}}
\newunicodechar{Σ}{\ensuremath{\Sigma}}
% majuscules produites par \MakeUppercase dans les titres (α-aminés -> Α)
\newunicodechar{Α}{\ensuremath{\alpha}}
\newunicodechar{Β}{\ensuremath{\beta}}
\newunicodechar{Γ}{\ensuremath{\Gamma}}
\newunicodechar{Δ}{\ensuremath{\Delta}}
\newunicodechar{Ε}{\ensuremath{\varepsilon}}
\newunicodechar{Θ}{\ensuremath{\Theta}}
\newunicodechar{Λ}{\ensuremath{\Lambda}}
\newunicodechar{Μ}{\ensuremath{\mu}}
\newunicodechar{Τ}{\ensuremath{\tau}}
\newunicodechar{Φ}{\ensuremath{\Phi}}
\newunicodechar{Ψ}{\ensuremath{\Psi}}
\newunicodechar{Ω}{\ensuremath{\Omega}}
\newunicodechar{↦}{\ensuremath{\mapsto}}
\newunicodechar{∈}{\ensuremath{\in}}
\newunicodechar{∉}{\ensuremath{\notin}}
\newunicodechar{∧}{\ensuremath{\wedge}}
\newunicodechar{⇔}{\ensuremath{\Leftrightarrow}}
\newunicodechar{⟺}{\ensuremath{\Longleftrightarrow}}
\newunicodechar{⇒}{\ensuremath{\Rightarrow}}
\newunicodechar{⟹}{\ensuremath{\Longrightarrow}}
\newunicodechar{∘}{\ensuremath{\circ}}
\newunicodechar{∪}{\ensuremath{\cup}}
\newunicodechar{∩}{\ensuremath{\cap}}
\newunicodechar{≡}{\ensuremath{\equiv}}
\newunicodechar{⊂}{\ensuremath{\subset}}
\newunicodechar{⊥}{\ensuremath{\perp}}
\newunicodechar{∃}{\ensuremath{\exists}}
\newunicodechar{∀}{\ensuremath{\forall}}
\newunicodechar{∛}{\ensuremath{\sqrt[3]{\,}}}
\newunicodechar{∜}{\ensuremath{\sqrt[4]{\,}}}
\newunicodechar{⃗}{\ensuremath{{}^{\to}}}
\newunicodechar{ⁿ}{\ifmmode^{n}\else\textsuperscript{\ensuremath{n}}\fi}
\newunicodechar{ˣ}{\ifmmode^{x}\else\textsuperscript{\ensuremath{x}}\fi}
\newunicodechar{ᵇ}{\ifmmode^{b}\else\textsuperscript{\ensuremath{b}}\fi}
\newunicodechar{ʳ}{\ifmmode^{r}\else\textsuperscript{\ensuremath{r}}\fi}
\newunicodechar{ᵘ}{\ifmmode^{u}\else\textsuperscript{\ensuremath{u}}\fi}
\newunicodechar{ʸ}{\ifmmode^{y}\else\textsuperscript{\ensuremath{y}}\fi}
\newunicodechar{ᵃ}{\ifmmode^{a}\else\textsuperscript{\ensuremath{a}}\fi}
\newunicodechar{ᵉ}{\ifmmode^{e}\else\textsuperscript{\ensuremath{e}}\fi}
\newunicodechar{ᵏ}{\ifmmode^{k}\else\textsuperscript{\ensuremath{k}}\fi}
\newunicodechar{ᵖ}{\ifmmode^{p}\else\textsuperscript{\ensuremath{p}}\fi}
\newunicodechar{ᴺ}{\ifmmode^{N}\else\textsuperscript{\ensuremath{N}}\fi}
\newunicodechar{ᶜ}{\ifmmode^{c}\else\textsuperscript{\ensuremath{c}}\fi}
\newunicodechar{ʰ}{\ifmmode^{h}\else\textsuperscript{\ensuremath{h}}\fi}
\newunicodechar{ᵠ}{\ifmmode^{q}\else\textsuperscript{\ensuremath{q}}\fi}
\newunicodechar{ₙ}{\ifmmode_{n}\else\textsubscript{\ensuremath{n}}\fi}
\newunicodechar{ₐ}{\ifmmode_{a}\else\textsubscript{\ensuremath{a}}\fi}
\newunicodechar{ᵢ}{\ifmmode_{i}\else\textsubscript{\ensuremath{i}}\fi}
\newunicodechar{ᵣ}{\ifmmode_{r}\else\textsubscript{\ensuremath{r}}\fi}
\newunicodechar{ₖ}{\ifmmode_{k}\else\textsubscript{\ensuremath{k}}\fi}
\newunicodechar{ᵦ}{\ifmmode_{\beta}\else\textsubscript{\ensuremath{\beta}}\fi}
\newunicodechar{ₓ}{\ifmmode_{x}\else\textsubscript{\ensuremath{x}}\fi}
\newfontfamily\popdisplay{ArchivoBlack-Regular.ttf}[Path=../../../fonts/]
\newfontfamily\poplogo{SpaceGrotesk-Bold.ttf}[Path=../../../fonts/]
\newfontfamily\popsemi{PlusJakartaSans-SemiBold.ttf}[Path=../../../fonts/]
\newfontfamily\popxbold{PlusJakartaSans-ExtraBold.ttf}[Path=../../../fonts/]
\newfontfamily\popxboldls{PlusJakartaSans-ExtraBold.ttf}[Path=../../../fonts/,LetterSpace=3.0]

\setlist[enumerate]{leftmargin=1.7em,itemsep=5pt,topsep=4pt}
\setlist[itemize]{leftmargin=1.7em,itemsep=4pt,topsep=4pt,label={\color{poporange}\textbullet}}
\setlength{\parindent}{0pt}
\setlength{\parskip}{5pt}
\renewcommand{\arraystretch}{1.25}

% pgfplots : style Pop par défaut
\pgfplotsset{
  every axis/.append style={axis line style={popink,line width=1pt},tick style={popink},
    grid style={popline,line width=0.5pt},label style={font=\small},tick label style={font=\footnotesize}},
  popcurve/.style={poporange,line width=1.8pt,no markers,smooth},
}
% Même style disponible dans un \draw TikZ simple (hors pgfplots)
\tikzset{popcurve/.style={poporange,line width=1.8pt}}
\tikzset{popgrid/.style={popline,very thin}}
\colorlet{popcurve}{poporange} % le nom du style est parfois utilisé comme couleur

% ============================================================
% Pill / squiggle
% ============================================================
\newcommand{\poppill}[3]{%
  \tikz[baseline=-0.55ex]\node[draw=popink,line width=1.2pt,rounded corners=2.6pt,%
    fill=#2,inner xsep=6.5pt,inner ysep=3pt]{%
    \popxboldls\fontsize{7}{7}\selectfont\color{#3}\MakeUppercase{#1}};%
}
\newcommand{\popsquiggle}[1]{%
  \tikz[baseline=-0.4ex]{
    \draw[#1,line width=2.6pt,line cap=round]
      plot [smooth] coordinates {(0,0.09) (0.34,0.01) (0.68,0.21) (1.02,0.01) (1.36,0.10)};
  }%
}
% Mot-clé surligné (marqueur orange sous le texte)
\newcommand{\cle}[1]{{\popxbold\color{popdark}#1}}

% ============================================================
% En-tête / pied
% ============================================================
\pagestyle{fancy}
\fancyhf{}
\renewcommand{\headrulewidth}{0pt}
\renewcommand{\footrulewidth}{0pt}
\lhead{\raisebox{-0.15ex}{\poplogo\fontsize{13}{13}\selectfont\color{popink}OnBuch\color{poporange}+}}
\rhead{\poppill{\DOCMATIERE\ \textbullet\ \DOCNIVEAU\ \textbullet\ Chapter \DOCLECON}{white}{popink}}
\lfoot{\popxbold\fontsize{7.6}{7.6}\selectfont\color{popmuted}\MakeUppercase{Signed course OnBuch+}\ \textbullet\ {\popsemi\DOCTITRE}}
\rfoot{\tikz[baseline=-0.6ex]\node[rounded corners=8.5pt,fill=popink,minimum height=17pt,minimum width=17pt,inner xsep=5pt,inner sep=0]{\popxbold\fontsize{8}{8}\selectfont\color{popcream}\thepage\,/\,\pageref*{LastPage}};}

% ============================================================
% Titres
% ============================================================
\newcommand{\chaptitre}[2]{%
  \begin{tcolorbox}[enhanced,colback=poporange,colframe=popink,boxrule=2.6pt,arc=11pt,
      drop shadow=popink,left=15pt,right=15pt,top=12pt,bottom=11pt,before skip=3pt,after skip=18pt]
    \poppill{#2}{popink}{popcream}\par\vspace{9pt}
    {\raggedright\hyphenpenalty=10000\exhyphenpenalty=10000\popdisplay\fontsize{22}{24}\selectfont\color{popcream}\MakeUppercase{#1}\par}
  \end{tcolorbox}
}
% Section numérotée : pastille orange avec le numéro + titre Archivo
\newcounter{popsec}
\newcommand{\coursec}[1]{%
  \stepcounter{popsec}\setcounter{popsub}{0}%
  \par\vspace{16pt}\noindent
  \begin{minipage}{\linewidth}
  \tikz[baseline=-0.7ex]\node[draw=popink,line width=1.6pt,fill=poporange,rounded corners=4pt,minimum size=22pt,inner sep=0]{\popdisplay\fontsize{12}{12}\selectfont\color{popcream}\thepopsec};%
  \hspace{8pt}{\popdisplay\fontsize{15}{17}\selectfont\color{popink}\MakeUppercase{#1}}\par
  \vspace{4pt}\noindent{\color{poporange}\rule{36pt}{3.6pt}}
  \end{minipage}\par\nopagebreak\vspace{8pt}
}
\newcounter{popsub}
\newcommand{\courssub}[1]{%
  \stepcounter{popsub}%
  \par\vspace{9pt}\noindent{\popxbold\fontsize{11.5}{13}\selectfont\color{popink}{\color{poporange}\thepopsec.\thepopsub}\hspace{6pt}#1}\par\nopagebreak\vspace{4pt}
}

% ============================================================
% Boîtes pédagogiques — même recette : fond blanc, bordure épaisse colorée,
% ombre dure encre, badge plein. Titre optionnel [#1].
% ============================================================
\tcbset{popbase/.style={breakable,enhanced,colback=white,boxrule=2.3pt,arc=9pt,
  shadow={3pt}{-3pt}{0pt}{popink},left=14pt,right=14pt,top=11pt,bottom=11pt,before skip=11pt,after skip=15pt,
  fontupper=\popsemi\fontsize{10.6}{14.5}\selectfont\color{popink},
  fontlower=\popsemi\fontsize{10.6}{14.5}\selectfont\color{popink}}}
% Coche dessinée (le glyphe ✓ n'existe pas dans les polices du document)
\newcommand{\popcheck}[1]{\tikz[baseline=-0.5ex]\draw[#1,line width=1.6pt,line cap=round,line join=round](-0.13,0.0)--(-0.04,-0.09)--(0.13,0.1);}
\providecommand{\coche}{\popcheck{popgreen}}
\providecommand{\resultat}[1]{\par\smallskip\noindent\fcolorbox{popgreen}{popgreenL}{\parbox{\dimexpr\linewidth-2\fboxsep-2\fboxrule\relax}{\textbf{Result.} #1}}\par\smallskip}
\newcommand{\popboxhead}[3]{\poppill{#1}{#2}{white}\ifx\relax#3\relax\else\hspace{6pt}{\popxbold\fontsize{10.6}{12}\selectfont\color{popink}#3}\fi\par\vspace{7pt}}

\newtcolorbox{aretenir}[1][]{popbase,colframe=popgreen,before upper={\popboxhead{Key points}{popgreen}{#1}}}
\newtcolorbox{exemplebox}[1][]{popbase,colframe=poporange,before upper={\popboxhead{Exemple}{poporange}{#1}}}
\newtcolorbox{definition}[1][]{popbase,colframe=popblue,before upper={\popboxhead{Definition}{popblue}{#1}}}
\newtcolorbox{propriete}[1][]{popbase,colframe=poppurple,before upper={\popboxhead{Property}{poppurple}{#1}}}
\newtcolorbox{methode}[1][]{popbase,colframe=popgold,before upper={\popboxhead{Method}{popgold}{#1}}}
\newtcolorbox{attention}[1][]{popbase,colframe=poppink,before upper={\popboxhead{Warning}{poppink}{#1}}}
\newtcolorbox{experience}[1][]{popbase,colframe=popblue,colback=popblueL,before upper={\popboxhead{Experiment}{popblue}{#1}}}
% « Le savais-tu ? » : carte violette pleine (contexte, culture, Cameroun)
\newtcolorbox{savaistu}[1][]{popbase,colback=poppurple,colframe=popink,
  fontupper=\popsemi\fontsize{10.4}{14.2}\selectfont\color{white},
  before upper={\poppill{Did you know?}{white}{poppurple}\ifx\relax#1\relax\else\hspace{6pt}{\popxbold\color{white}#1}\fi\par\vspace{7pt}}}
% Exercice résolu : énoncé en haut, \tcblower puis la solution sur fond crème
\newcounter{popexo}\newcounter{popexr}
\newtcolorbox{exoresolu}[1][]{popbase,colframe=poporange,colbacklower=popcream,
  segmentation style={popink,line width=1.2pt,dashed},
  before upper={\stepcounter{popexr}\popboxhead{Worked example \thepopexr}{poporange}{#1}},
  before lower={\poppill{Solution}{popink}{popcream}\par\vspace{6pt}}}
% Exercice (non résolu en place) — la correction est dans la partie Corrigés
\newtcolorbox{exercice}[1][]{popbase,colframe=popink,
  before upper={\stepcounter{popexo}\popboxhead{Exercise \thepopexo}{popink}{#1}}}
% Corrigé
\newtcolorbox{corrige}[1][]{popbase,colframe=popgreen,colback=popgreenL,
  before upper={\popboxhead{Solution}{popgreen}{#1}}}
% Objectifs / prérequis / bilan
\newtcolorbox{objectifs}{popbase,colframe=popink,colback=poporangeL,
  before upper={\popboxhead{Chapter objectives}{popink}{}}}
\newtcolorbox{prerequis}{popbase,colframe=popink,colback=popblueL,
  before upper={\popboxhead{Prerequisites}{popblue}{}}}
\newtcolorbox{bilan}{popbase,colframe=popink,colback=white,boxrule=2.8pt,
  before upper={\popboxhead{Fiche bilan}{poporange}{L'essentiel à connaître pour l'examen}}}
% Figure encadrée avec légende
\newtcolorbox{popfigure}[1][]{enhanced,breakable=false,colback=white,colframe=popline,boxrule=1.4pt,arc=8pt,
  left=10pt,right=10pt,top=10pt,bottom=8pt,before skip=10pt,after skip=12pt,halign=center,
  lower separated=false,halign lower=center,
  fontlower=\popsemi\fontsize{9}{11.5}\selectfont\color{popmuted},
  #1}
\newcommand{\legende}[1]{\par\vspace{6pt}{\popsemi\fontsize{9}{11.5}\selectfont\color{popmuted}#1}}

% ============================================================
% Signature OnBuch+
% ============================================================
\newcommand{\poplogoinline}[1]{{\poplogo\fontsize{#1}{#1}\selectfont\color{popink}OnBuch\color{poporange}+}}
\newcommand{\signatureonbuch}{%
  \begin{tcolorbox}[enhanced,colback=popink,colframe=popink,boxrule=2.2pt,arc=10pt,
      drop shadow=poporange,left=14pt,right=14pt,top=10pt,bottom=10pt,before skip=0pt,after skip=0pt]
    \tikz[baseline=-0.7ex]\node[circle,fill=popgreen,draw=popcream,line width=1.3pt,minimum size=22pt,inner sep=0]{\popcheck{white}};%
    \hspace{9pt}\begin{minipage}[c]{0.62\linewidth}
      {\popxbold\fontsize{12}{13}\selectfont\color{popcream}Signed\ \textbullet\ {\poplogo OnBuch\color{poporange}+}}\par\vspace{2pt}
      {\raggedright\popsemi\fontsize{8}{10.5}\selectfont\color{popline}\MakeUppercase{OnBuch+ teaching team}\\\MakeUppercase{Follows the official MINESEC syllabus}\par}
    \end{minipage}\hfill
    {\poplogo\fontsize{22}{22}\selectfont\color{popcream}OnBuch\color{poporange}+}
  \end{tcolorbox}%
}

% ============================================================
% Page de garde
% #1 titre · #2 matière · #3 niveau · #4 module · #5 n° leçon · #6 durée ·
% #7 description / objectifs (texte ou itemize)
% ============================================================
\newcommand{\pagegarde}[7]{%
  \thispagestyle{empty}
  \newgeometry{a4paper,margin=1.7cm,top=1.6cm,bottom=1.4cm}%
  \begin{tikzpicture}[remember picture,overlay]
    \fill[poporange!18] ([xshift=-3.2cm,yshift=-3.6cm]current page.north east) circle (4.6cm);
    \fill[poppurple!14] ([xshift=2.4cm,yshift=7.6cm]current page.south west) circle (3.8cm);
  \end{tikzpicture}%
  {\poplogo\fontsize{30}{30}\selectfont\color{popink}OnBuch\color{poporange}+}\hfill
  \raisebox{0.6ex}{\poppill{Signed course \textbullet\ Premium}{popink}{popcream}}\par
  \vspace{1pt}\popsquiggle{poppurple}\par
  \vspace{8pt}
  \begin{tcolorbox}[enhanced,colback=poporange,colframe=popink,boxrule=2.8pt,arc=13pt,
      drop shadow=popink,left=18pt,right=18pt,top=12pt,bottom=13pt,after skip=10pt]
    \poppill{#2 \textbullet\ #3}{popink}{popcream}\hspace{5pt}\poppill{#4}{white}{popink}\par\vspace{8pt}
    {\popdisplay\fontsize{14}{14}\selectfont\color{popink}CHAPTER #5}\par\vspace{4pt}
    {\raggedright\hyphenpenalty=10000\exhyphenpenalty=10000\popdisplay\fontsize{25}{27}\selectfont\color{popcream}\MakeUppercase{#1}\par}
    \vspace{8pt}
    {\popxbold\fontsize{9.5}{9.5}\selectfont\color{popink}\MakeUppercase{Suggested time: #6}}
  \end{tcolorbox}
  \begin{tcolorbox}[enhanced,colback=white,colframe=popink,boxrule=2.2pt,arc=10pt,
      drop shadow=popink,left=16pt,right=16pt,top=10pt,bottom=10pt,after skip=8pt]
    \poppill{In this chapter}{poppurple}{white}\par\vspace{5pt}
    \popsemi\fontsize{9.3}{12.6}\selectfont\color{popink}#7
  \end{tcolorbox}
  \noindent\begin{minipage}[t]{0.487\linewidth}
    \begin{tcolorbox}[enhanced,colback=popgreen,colframe=popink,boxrule=2.2pt,arc=10pt,
        drop shadow=popink,left=11pt,right=11pt,top=10pt,bottom=10pt,height=3.1cm,valign=center]
      \tcbox[colback=white,colframe=popink,boxrule=1.3pt,arc=5pt,boxsep=0pt,nobeforeafter,
        left=3pt,right=3pt,top=3pt,bottom=3pt,valign=center]{\qrcode[height=1.9cm]{https://play.google.com/store/apps/details?id=com.luvvix.onbuch&hl=en}}%
      \hspace{8pt}\begin{minipage}[c]{0.5\linewidth}
        {\popdisplay\fontsize{11}{12.5}\selectfont\color{white}\MakeUppercase{Download OnBuch}}\par\vspace{4pt}
        {\raggedright\popsemi\fontsize{8.4}{11}\selectfont\color{white}Free \textbullet\ Google Play\\AI tutor, past papers, results\par}
      \end{minipage}
    \end{tcolorbox}
  \end{minipage}\hfill
  \begin{minipage}[t]{0.487\linewidth}
    \begin{tcolorbox}[enhanced,colback=poppurple,colframe=popink,boxrule=2.2pt,arc=10pt,
        drop shadow=popink,left=11pt,right=11pt,top=10pt,bottom=10pt,height=3.1cm,valign=center]
      \tikz[baseline=-0.7ex]\node[circle,fill=white,draw=popink,line width=1.3pt,minimum size=1.6cm,inner sep=0]{\popdisplay\fontsize{24}{24}\selectfont\color{poppurple}+};%
      \hspace{8pt}\begin{minipage}[c]{0.6\linewidth}
        {\popdisplay\fontsize{11}{12.5}\selectfont\color{white}\MakeUppercase{Go OnBuch+}}\par\vspace{4pt}
        {\raggedright\popsemi\fontsize{8.4}{11}\selectfont\color{white}All signed courses, unlimited AI tutor, PDF sheets — OnBuch+ tab in the app\par}
      \end{minipage}
    \end{tcolorbox}
  \end{minipage}
  \vspace{8pt}
  \popcontenu
  \vfill
  \signatureonbuch
  \restoregeometry
  \newpage
}

% Rangée « Ce cours contient » (page de garde)
\newcommand{\poptile}[3]{% #1 couleur #2 gros texte #3 légende
  \begin{tcolorbox}[enhanced,colback=#1,colframe=popink,boxrule=2pt,arc=9pt,shadow={2.5pt}{-2.5pt}{0pt}{popink},
      left=6pt,right=6pt,top=9pt,bottom=9pt,halign=center,valign=center,height=1.75cm,width=\linewidth,nobeforeafter]
    {\popdisplay\fontsize{12.5}{13}\selectfont\color{white}\MakeUppercase{#2}}\par\vspace{3pt}
    {\centering\popsemi\fontsize{7.8}{9.5}\selectfont\color{white}#3\par}
  \end{tcolorbox}}
\newcommand{\popcontenu}{%
  \par\noindent{\popxboldls\fontsize{7.5}{7.5}\selectfont\color{popmuted}THIS SIGNED COURSE CONTAINS}\par\vspace{7pt}
  \noindent\begin{minipage}[t]{0.235\linewidth}\poptile{popblue}{Course}{complete, illustrated, step by step}\end{minipage}\hfill
  \begin{minipage}[t]{0.235\linewidth}\poptile{poporange}{Methods}{and worked examples}\end{minipage}\hfill
  \begin{minipage}[t]{0.235\linewidth}\poptile{poppink}{Exercises}{exam style + solutions}\end{minipage}\hfill
  \begin{minipage}[t]{0.235\linewidth}\poptile{popgreen}{Summary sheet}{the essentials to remember}\end{minipage}\par}

% Sommaire
\newtcolorbox{sommaire}{enhanced,colback=white,colframe=popink,boxrule=2.2pt,arc=10pt,
  drop shadow=popink,left=18pt,right=18pt,top=13pt,bottom=13pt,before skip=6pt,after skip=20pt,
  before upper={\poppill{Contents}{poppurple}{white}\par\vspace{9pt}\popsemi\fontsize{10.6}{14.5}\selectfont\color{popink}}}

% Signature de fin de cours — poussée tout en bas de la dernière page
\newcommand{\findecours}{%
  \par\vfill\nopagebreak
  \begin{center}\popsquiggle{poporange}\end{center}\vspace{4pt}
  \signatureonbuch
}
\AtBeginDocument{\def\llbracket{\mathopen{[\mkern-3.5mu[}}\def\rrbracket{\mathclose{]\mkern-3.5mu]}}} % intervalles d'entiers (glyphe ⟦ absent de la police)
% Glyphes absents de Fira Math : redéfinis à partir de glyphes disponibles
\AtBeginDocument{%
  \def\setminus{\mathbin{\backslash}}\let\smallsetminus\setminus
  \def\vdots{\mathord{\vbox{\baselineskip3.2pt\lineskiplimit0pt\kern4pt\hbox{.}\hbox{.}\hbox{.}}}}%
  \def\ddots{\mathinner{\mkern1mu\raise6.5pt\hbox{.}\mkern2mu\raise3.6pt\hbox{.}\mkern2mu\raise0.7pt\hbox{.}\mkern1mu}}%
  \def\triangle{\mathord{\scalebox{0.9}{$\symup{\Delta}$}}}%
  \def\nabla{\mathord{\raisebox{\depth}{\scalebox{1}[-1]{$\symup{\Delta}$}}}}%
  \def\checkmark{\popcheck{popdark}}%
  \def\star{\mathbin{\ast}}\def\bigstar{\mathord{\ast}}%
  \def\diamond{\mathbin{\scalebox{0.6}{$\Diamond$}}}%
  \def\bigcup{\mathop{\mathchoice{\raisebox{-0.3em}{\scalebox{1.7}{$\cup$}}}{\scalebox{1.25}{$\cup$}}{\cup}{\cup}}\displaylimits}%
  \def\bigcap{\mathop{\mathchoice{\raisebox{-0.3em}{\scalebox{1.7}{$\cap$}}}{\scalebox{1.25}{$\cap$}}{\cap}{\cap}}\displaylimits}%
  \def\lesseqqgtr{\mathrel{\lessgtr}}\def\bigoplus{\mathop{\oplus}\displaylimits}%
}
\providecommand{\ssi}{if and only if} % abréviation fréquente
\AtBeginDocument{\providecommand{\wideparen}{\overparen}} % arc de cercle

%% FIGFIX-BEGIN v5 — correctifs globaux des figures (docs/onbuch-generation/tools/figfix)
\usepackage{adjustbox}\usepackage{etoolbox}\tcbuselibrary{raster}
% toute figure TikZ/pgfplots plus large que la ligne est réduite à la largeur disponible
\BeforeBeginEnvironment{tikzpicture}{\begin{adjustbox}{max width=\linewidth}}
\AfterEndEnvironment{tikzpicture}{\end{adjustbox}}
% légendes pgfplots lisibles par-dessus les courbes
\pgfplotsset{every axis/.append style={legend style={fill=white,fill opacity=0.92,text opacity=1,draw=black!15,inner sep=3pt}}}
% étiquettes d'arêtes (syntaxe edge["..."]) avec fond blanc
\tikzset{every edge quotes/.append style={fill=white,inner sep=1.2pt}}
%% FIGFIX-END
\begin{document}

\pagegarde{Velocity and Acceleration as Vectors}{Pure Mathematics with Mechanics}{Upper Sixth}{Upper Sixth — Pure mathematics with mechanics}{10}{10 periods}{This chapter extends kinematics to two dimensions by treating velocity and acceleration as vectors. Learners master relative velocity, true velocity in river and wind problems, interception and collision courses, and relative position, with full mathematical treatment and GCE-style applications.
\vspace{4pt}
\begin{multicols}{2}\small
\begin{itemize}
\item By the end of this chapter I can represent velocities and accelerations as vectors using i, j components and bearings.
\item By the end of this chapter I can define and compute the relative velocity of one body with respect to another using v\_A/B = v\_A − v\_B.
\item By the end of this chapter I can resolve non-parallel velocities and forces to find resultant velocity magnitude and direction.
\item By the end of this chapter I can solve river-crossing and wind problems, distinguishing true (ground) velocity from heading velocity.
\item By the end of this chapter I can determine the course to steer for quickest crossing, shortest crossing, or arrival at a given point.
\item By the end of this chapter I can model interception and collision problems and find the time and position of meeting.
\end{itemize}\end{multicols}}

\begin{sommaire}
\begin{multicols}{2}
\begin{enumerate}
\item Velocity and Acceleration as Vectors; Concept of Relative Velocity
\item Non-Parallel Forces, Resultant Velocity and True Velocity in River or Wind Problems
\item Applications of Relative and True Velocity
\item Interception and Collision: Concept and Mathematical Treatment
\item Applications of Interception and Collision
\item Relative Position and Relative Motion: Concept, Treatment and Applications
\item Integration activity
\item Methods and worked examples
\item Exercises
\item Solutions to the exercises
\item Summary sheet
\end{enumerate}
\end{multicols}
\end{sommaire}

\begin{objectifs}
\begin{itemize}
\item By the end of this chapter I can represent velocities and accelerations as vectors using i, j components and bearings.
\item By the end of this chapter I can define and compute the relative velocity of one body with respect to another using v\_A/B = v\_A − v\_B.
\item By the end of this chapter I can resolve non-parallel velocities and forces to find resultant velocity magnitude and direction.
\item By the end of this chapter I can solve river-crossing and wind problems, distinguishing true (ground) velocity from heading velocity.
\item By the end of this chapter I can determine the course to steer for quickest crossing, shortest crossing, or arrival at a given point.
\item By the end of this chapter I can model interception and collision problems and find the time and position of meeting.
\item By the end of this chapter I can find the closest distance of approach between two moving bodies and the time at which it occurs.
\item By the end of this chapter I can interpret and compute relative position vectors and use them to analyse relative motion.
\item By the end of this chapter I can solve GCE-style applied problems combining relative velocity, interception and relative position.
\end{itemize}
\end{objectifs}

\begin{prerequis}
\begin{itemize}
\item Vectors in component form: addition, subtraction, magnitude, unit vectors i and j (Lower Sixth Pure Mathematics).
\item Trigonometry: sine and cosine rules, bearings, resolving a vector into perpendicular components.
\item Kinematics in one dimension: constant velocity and constant acceleration formulae (Lower Sixth Mechanics).
\item Resultant of forces by the parallelogram and triangle laws (Lower Sixth Mechanics).
\item Solving simultaneous linear equations and simple quadratic equations.
\end{itemize}
\end{prerequis}

\newpage

\coursec{Velocity and Acceleration as Vectors; Concept of Relative Velocity}

Every rainy season, the ferry crossing the Wouri estuary at Douala points its bow noticeably upstream, yet passengers arrive exactly at the landing stage on the opposite bank. Meanwhile, a Wouri pilot boat racing to intercept an incoming cargo ship must aim not at where the ship \emph{is}, but at where it \emph{will be}. On the highway between Yaoundé and Bafoussam, a driver overtaking a truck feels the truck ``crawl backwards'' even though both vehicles move forward. All these everyday observations share one mathematical engine: \cle{velocity treated as a vector}. This chapter gives you the tools to explain, predict and compute such motions — a favourite source of GCE Advanced Level mechanics questions.

The central problem of this first section is simple to state: \textbf{how do we describe motion when direction matters, and how does one moving object appear to another moving observer?} We begin by revising the vector description of motion, then we build the key new idea — \cle{relative velocity}.

\courssub{Position, Velocity and Acceleration as Vectors}

\textbf{Intuition.} Suppose a motorcycle taxi (\emph{bendskin}) moves through Douala. Saying ``its speed is $40\ \mathrm{km\,h^{-1}}$'' does not tell you where it is going: it could be heading north towards Mboppi or east towards the Wouri bridge. To describe the motion completely we need a \emph{magnitude} and a \emph{direction} — that is, a \cle{vector}.

Fix an origin $O$ and unit vectors $\vec{i}$ (pointing East) and $\vec{j}$ (pointing North). The position of a particle $P$ at time $t$ is described by its \cle{position vector}
\[
\vec{r}(t) = x(t)\,\vec{i} + y(t)\,\vec{j},
\]
where $x(t)$ and $y(t)$ are the coordinates of $P$ at time $t$. As $P$ moves, $\vec{r}$ changes with time, and the natural way to measure that change is differentiation.

\begin{definition}[Velocity and acceleration as vectors]
Let a particle have position vector $\vec{r}(t) = x(t)\,\vec{i} + y(t)\,\vec{j}$ at time $t$. Its \cle{velocity} is the vector
\[
\vec{v} = \frac{d\vec{r}}{dt} = \dot{x}\,\vec{i} + \dot{y}\,\vec{j},
\]
and its \cle{acceleration} is the vector
\[
\vec{a} = \frac{d\vec{v}}{dt} = \ddot{x}\,\vec{i} + \ddot{y}\,\vec{j}.
\]
The \cle{speed} is the magnitude of the velocity: $|\vec{v}| = \sqrt{\dot{x}^2 + \dot{y}^2}$. Velocity is measured in $\mathrm{m\,s^{-1}}$ and acceleration in $\mathrm{m\,s^{-2}}$.
\end{definition}

Because $\vec{i}$ and $\vec{j}$ are constant vectors, differentiating $\vec{r}$ simply means differentiating each component. This is what makes the component form so powerful: a two-dimensional problem becomes two one-dimensional problems.

\begin{exemplebox}[From position to velocity and acceleration]
A particle moves so that $\vec{r}(t) = (3t^2 + 1)\vec{i} + (4t - t^2)\vec{j}$ (metres). Then
\[
\vec{v} = \frac{d\vec{r}}{dt} = 6t\,\vec{i} + (4 - 2t)\,\vec{j}, \qquad
\vec{a} = \frac{d\vec{v}}{dt} = 6\,\vec{i} - 2\,\vec{j}.
\]
At $t = 1$: $\vec{v} = 6\vec{i} + 2\vec{j}$, so the speed is $|\vec{v}| = \sqrt{36 + 4} = \sqrt{40} \approx 6.32\ \mathrm{m\,s^{-1}}$. The acceleration is constant: $\vec{a} = 6\vec{i} - 2\vec{j}\ \mathrm{m\,s^{-2}}$.
\end{exemplebox}

\begin{attention}[Speed is not velocity]
The \emph{speed} $|\vec{v}|$ is a scalar; the \emph{velocity} $\vec{v}$ is a vector. Two vehicles can have the same speed and completely different velocities. In GCE answers, never write a velocity as a bare number unless the motion is one-dimensional and the direction is stated.
\end{attention}

\courssub{From Speed and Bearing to Components — and Back}

In navigation problems (rivers, wind, aircraft, ships), a velocity is usually given as a \textbf{speed together with a bearing}. Recall that a \cle{bearing} is measured \emph{clockwise from North}. To compute with such a velocity we must express it in the $\vec{i}, \vec{j}$ form, with $\vec{i}$ East and $\vec{j}$ North.

If a particle has speed $v$ on a bearing $\theta$, draw the velocity vector from the origin. The angle $\theta$ is measured from the North ($\vec{j}$) axis towards the East ($\vec{i}$) axis. Resolving:
\[
\vec{v} = v\sin\theta\,\vec{i} + v\cos\theta\,\vec{j}.
\]
Note the ``swap'' compared with angles measured from the $x$-axis: the $\vec{i}$-component carries $\sin\theta$ and the $\vec{j}$-component carries $\cos\theta$, because $\theta$ is measured from the \emph{North} axis.

\begin{methode}[Converting between (speed, bearing) and components]
\begin{enumerate}
\item \textbf{(Speed, bearing) $\to$ components.} Given speed $v$ and bearing $\theta$: write $\vec{v} = v\sin\theta\,\vec{i} + v\cos\theta\,\vec{j}$. Check the signs against the quadrant of the bearing.
\item \textbf{Components $\to$ speed.} Given $\vec{v} = v_1\vec{i} + v_2\vec{j}$: the speed is $v = \sqrt{v_1^2 + v_2^2}$.
\item \textbf{Components $\to$ bearing.} Compute the acute angle $\alpha = \tan^{-1}\left|\dfrac{v_1}{v_2}\right|$ from the North axis, then use the signs of $v_1$ (East $+$/West $-$) and $v_2$ (North $+$/South $-$) to place the vector in the correct quadrant and deduce the bearing $\theta$.
\item \textbf{Check.} Verify that $v\sin\theta$ and $v\cos\theta$ recover the components you started with.
\end{enumerate}
\end{methode}

\begin{popfigure}
\begin{center}
\begin{tikzpicture}[scale=1.1, >=stealth]
\draw[->, thick, popmuted] (-0.3,0) -- (4.2,0) node[below left] {East $\vec{i}$};
\draw[->, thick, popmuted] (0,-0.3) -- (0,3.6) node[below left] {North $\vec{j}$};
\draw[->, very thick, popblue] (0,0) -- (3.464,2) node[above right] {$\vec{v}$};
\draw[->, thick, poporange] (0,0) -- (3.464,0) node[midway, below] {$v\sin 60^\circ$};
\draw[->, thick, popgreen] (3.464,0) -- (3.464,2) node[midway, right] {$v\cos 60^\circ$};
\draw[popdark] (0,1.1) arc (90:30:1.1) node[midway, above left=1pt and -6pt] {$60^\circ$};
\draw[popmuted, dashed] (0,2) -- (3.464,2);
\node[popdark] at (1.7,-0.9) {Bearing $060^\circ$, speed $v$};
\end{tikzpicture}
\end{center}
\legende{Converting a velocity given on a bearing of $060^\circ$ into $\vec{i}, \vec{j}$ components: the bearing is measured clockwise from North, so the East component is $v\sin 60^\circ$ and the North component is $v\cos 60^\circ$.}
\end{popfigure}

\begin{exoresolu}[Speed and bearing to components, and back]
\textbf{(a)} A boat sails at $10\ \mathrm{km\,h^{-1}}$ on a bearing of $060^\circ$. Express its velocity in the form $v_1\vec{i} + v_2\vec{j}$. \quad \textbf{(b)} A wind has velocity $\vec{w} = -6\vec{i} + 8\vec{j}\ \mathrm{km\,h^{-1}}$. Find the wind speed and the bearing on which it blows.
\tcblower
\textbf{(a)} With $v = 10$ and $\theta = 60^\circ$:
\[
\vec{v} = 10\sin 60^\circ\,\vec{i} + 10\cos 60^\circ\,\vec{j} = 5\sqrt{3}\,\vec{i} + 5\,\vec{j} \approx 8.66\,\vec{i} + 5\,\vec{j}\ \mathrm{km\,h^{-1}}.
\]
Both components are positive, consistent with a bearing in the North-East quadrant.

\textbf{(b)} Speed: $|\vec{w}| = \sqrt{(-6)^2 + 8^2} = \sqrt{100} = 10\ \mathrm{km\,h^{-1}}$. The acute angle from the North axis is $\alpha = \tan^{-1}\left|\dfrac{-6}{8}\right| = \tan^{-1}(0.75) \approx 36.9^\circ$. Since the $\vec{i}$-component is negative (West) and the $\vec{j}$-component positive (North), the vector points North-West, so the bearing is
\[
\theta = 360^\circ - 36.9^\circ \approx 323.1^\circ.
\]
Check: $10\sin 323.1^\circ \approx -6$ and $10\cos 323.1^\circ \approx 8$. \checkmark
\end{exoresolu}

\courssub{The Concept of Relative Velocity}

\textbf{Intuition.} You are in a car on the Douala--Yaoundé highway travelling at $90\ \mathrm{km\,h^{-1}}$, and you overtake a truck moving at $70\ \mathrm{km\,h^{-1}}$ in the same direction. Through your window, the truck does not appear to rush forwards at $70\ \mathrm{km\,h^{-1}}$; it appears to drift \emph{backwards} past you at $20\ \mathrm{km\,h^{-1}}$. Your brain has automatically subtracted the two velocities. That subtraction — performed on \emph{vectors}, not just on numbers — is exactly the notion of relative velocity.

\begin{definition}[Relative velocity]
Let $A$ and $B$ be two particles moving with velocities $\vec{v}_A$ and $\vec{v}_B$ (measured relative to the ground). The \cle{velocity of $A$ relative to $B$}, also called the \cle{relative velocity of $A$ with respect to $B$}, is
\[
\vec{v}_{A/B} = \vec{v}_A - \vec{v}_B.
\]
It is the velocity that an observer moving with $B$ would attribute to $A$.
\end{definition}

\textbf{Geometric interpretation.} Since $\vec{v}_{A/B} = \vec{v}_A - \vec{v}_B = \vec{v}_A + (-\vec{v}_B)$, we construct it by the triangle law: draw $\vec{v}_A$, then from its tip draw $-\vec{v}_B$ (the vector $\vec{v}_B$ reversed); the closing side from the origin to the final tip is $\vec{v}_{A/B}$. Equivalently, if $\vec{v}_A$ and $\vec{v}_B$ are drawn from the \emph{same} origin, then $\vec{v}_{A/B}$ is the vector from the \emph{tip of $\vec{v}_B$ to the tip of $\vec{v}_A$}.

\begin{popfigure}
\begin{center}
\begin{tikzpicture}[scale=0.55, >=stealth]
\draw[thick, popmuted] (0,-1.2) -- (14,-1.2);
\draw[thick, popmuted] (0,-3.2) -- (14,-3.2);
\foreach \x in {1,3,5,7,9,11,13} \draw[popmuted, dashed] (\x,-2.2) -- (\x+0.7,-2.2);
\draw[fill=popblueL] (1.5,-1.2) rectangle (3.5,-0.4);
\node[popblue] at (2.5,-0.05) {$A$};
\draw[->, very thick, popblue] (3.5,-0.8) -- (8.5,-0.8) node[midway, above] {$\vec{v}_A$};
\draw[fill=popgreenL] (1.5,-3.2) rectangle (3.5,-2.4);
\node[popgreen] at (2.5,-3.55) {$B$};
\draw[->, very thick, popgreen] (3.5,-2.8) -- (6.5,-2.8) node[midway, above] {$\vec{v}_B$};
\node[popdark, align=center] at (7,-4.6) {Same direction: $|\vec{v}_{A/B}| = v_A - v_B$};
\end{tikzpicture}
\end{center}
\legende{Car $A$ overtaking truck $B$ on a straight road. Because the velocities are parallel, the relative velocity of $A$ with respect to $B$ has magnitude $v_A - v_B$: $A$ appears to pull ahead slowly.}
\end{popfigure}

\begin{popfigure}
\begin{center}
\begin{tikzpicture}[scale=0.9, >=stealth]
\draw[->, thick, popmuted] (-0.3,0) -- (6.5,0) node[below left] {$\vec{i}$};
\draw[->, thick, popmuted] (0,-2.6) -- (0,4.2) node[below left] {$\vec{j}$};
\draw[->, very thick, popblue] (0,0) -- (5,3) node[above right] {$\vec{v}_A$};
\draw[->, very thick, popgreen] (0,0) -- (4,-1.5) node[below right] {$\vec{v}_B$};
\draw[->, very thick, poporange] (4,-1.5) -- (5,3) node[midway, right] {$\vec{v}_{A/B} = \vec{v}_A - \vec{v}_B$};
\draw[popmuted, dashed] (5,0) -- (5,3);
\draw[popmuted, dashed] (0,3) -- (5,3);
\node[popblue] at (5.35,1.5) {$v_{Ay}$};
\node[popblue] at (2.5,3.25) {$v_{Ax}$};
\node[popdark, align=center, text width=4.5cm] at (2.2,-2.3) {From tip of $\vec{v}_B$ to tip of $\vec{v}_A$};
\end{tikzpicture}
\end{center}
\legende{Vector triangle for relative velocity: with $\vec{v}_A$ and $\vec{v}_B$ drawn from the same origin, $\vec{v}_{A/B}$ runs from the tip of $\vec{v}_B$ to the tip of $\vec{v}_A$. Components subtract: $\vec{v}_{A/B} = (v_{Ax} - v_{Bx})\vec{i} + (v_{Ay} - v_{By})\vec{j}$.}
\end{popfigure}

In components the subtraction is immediate: if $\vec{v}_A = a_1\vec{i} + a_2\vec{j}$ and $\vec{v}_B = b_1\vec{i} + b_2\vec{j}$, then
\[
\vec{v}_{A/B} = (a_1 - b_1)\,\vec{i} + (a_2 - b_2)\,\vec{j}.
\]

\begin{methode}[Computing relative velocity $\vec{v}_{A/B}$]
\begin{enumerate}
\item Write both velocities in component form: $\vec{v}_A = a_1\vec{i} + a_2\vec{j}$, $\vec{v}_B = b_1\vec{i} + b_2\vec{j}$.
\item Subtract component-wise: $\vec{v}_{A/B} = (a_1 - b_1)\vec{i} + (a_2 - b_2)\vec{j}$.
\item If required, find the magnitude: $|\vec{v}_{A/B}| = \sqrt{(a_1 - b_1)^2 + (a_2 - b_2)^2}$.
\item If required, find the bearing: compute $\alpha = \tan^{-1}\left|\dfrac{a_1 - b_1}{a_2 - b_2}\right|$ and determine the quadrant from the signs of the components.
\end{enumerate}
\end{methode}

\begin{exoresolu}[Computing a relative velocity: magnitude and bearing]
A passenger ferry $A$ crosses the Wouri estuary with velocity $\vec{v}_A = 3\vec{i} + 4\vec{j}\ \mathrm{m\,s^{-1}}$, while a pirogue $B$ moves with velocity $\vec{v}_B = 5\vec{i} + \vec{j}\ \mathrm{m\,s^{-1}}$. Find the velocity of the ferry relative to the pirogue, giving its magnitude and bearing.
\tcblower
\[
\vec{v}_{A/B} = \vec{v}_A - \vec{v}_B = (3 - 5)\vec{i} + (4 - 1)\vec{j} = -2\vec{i} + 3\vec{j}\ \mathrm{m\,s^{-1}}.
\]
Magnitude: $|\vec{v}_{A/B}| = \sqrt{(-2)^2 + 3^2} = \sqrt{13} \approx 3.61\ \mathrm{m\,s^{-1}}$.

Direction: the acute angle from the North axis is $\alpha = \tan^{-1}\left|\dfrac{-2}{3}\right| \approx 33.7^\circ$. The $\vec{i}$-component is negative (West) and the $\vec{j}$-component positive (North), so the vector points North-West and the bearing is
\[
360^\circ - 33.7^\circ \approx 326.3^\circ.
\]
Interpretation: to someone sitting in the pirogue, the ferry appears to move at about $3.61\ \mathrm{m\,s^{-1}}$ on a bearing of $326^\circ$.
\end{exoresolu}

\courssub{Velocity of B Relative to A — and the One-Dimensional Special Case}

\begin{propriete}[Anti-symmetry of relative velocity]
For any two particles $A$ and $B$,
\[
\vec{v}_{B/A} = \vec{v}_B - \vec{v}_A = -\bigl(\vec{v}_A - \vec{v}_B\bigr) = -\vec{v}_{A/B}.
\]
The two relative velocities have the \emph{same magnitude} but \emph{opposite directions}. (Proof: immediate from the definition — not examinable, but you must use it fluently.)
\end{propriete}

This matches experience perfectly: if you overtake the truck, you see it drifting backwards at $20\ \mathrm{km\,h^{-1}}$, while the truck driver sees \emph{you} pulling forwards at $20\ \mathrm{km\,h^{-1}}$. Same magnitude, opposite direction.

\begin{attention}[Do not confuse $\vec{v}_{A/B}$ with $\vec{v}_{B/A}$]
A very common GCE error is to subtract the velocities in the wrong order. Remember: \emph{``velocity of $A$ relative to $B$''} means \emph{subtract the velocity of the observer $B$}: $\vec{v}_{A/B} = \vec{v}_A - \vec{v}_B$. If your answer comes out in the opposite direction to what the situation suggests, you have probably computed $\vec{v}_{B/A}$. A useful consistency check: $\vec{v}_{A/B} + \vec{v}_{B/A} = \vec{0}$.
\end{attention}

\begin{attention}[Subtract vectors, not speeds]
Relative velocity is a \emph{vector} subtraction. Writing $|\vec{v}_{A/B}| = |\vec{v}_A| - |\vec{v}_B|$ is \textbf{wrong} in general — it only works when the two velocities are parallel and in the same sense. For example, if $\vec{v}_A = 3\vec{i} + 4\vec{j}$ (speed $5$) and $\vec{v}_B = 5\vec{i} + \vec{j}$ (speed $\sqrt{26} \approx 5.1$), then $|\vec{v}_{A/B}| = \sqrt{13} \approx 3.61$, which is neither $5 - 5.1$ nor $5.1 - 5$. Always subtract the \emph{components}, then take the magnitude.
\end{attention}

\textbf{GCE extra: relative velocity in one dimension.} When both motions lie on the same straight line, choose a positive sense and represent velocities by signed scalars. Two cases occur constantly in examinations:

\begin{itemize}
\item \textbf{Same direction.} Two vehicles moving the same way at speeds $v_A > v_B$: the faster one gains on the slower at $v_A - v_B$. Overtaking feels slow because the \emph{relative} speed is small.
\item \textbf{Opposite directions.} Two trains passing each other at speeds $v_A$ and $v_B$: taking the direction of $A$ as positive, $\vec{v}_B$ is negative, so the relative speed is $v_A - (-v_B) = v_A + v_B$. This is why an oncoming train flashes past your window so violently: the relative speed is the \emph{sum} of the two speeds.
\end{itemize}

\begin{exemplebox}[Trains passing on the Douala--Ngaoundéré line]
Train $A$ travels North at $60\ \mathrm{km\,h^{-1}}$ and train $B$ travels South at $80\ \mathrm{km\,h^{-1}}$ on the parallel track. Taking North as positive, $v_A = 60$ and $v_B = -80\ \mathrm{km\,h^{-1}}$. The velocity of $B$ relative to $A$ is
\[
v_{B/A} = v_B - v_A = -80 - 60 = -140\ \mathrm{km\,h^{-1}},
\]
i.e.\ a passenger in $A$ sees $B$ sweep past southwards at $140\ \mathrm{km\,h^{-1}}$ — the sum of the speeds, because the motions are in opposite directions.
\end{exemplebox}

\begin{aretenir}[Key points of this section]
\begin{itemize}
\item Position, velocity and acceleration are vectors: $\vec{v} = \dfrac{d\vec{r}}{dt} = \dot{x}\vec{i} + \dot{y}\vec{j}$, $\vec{a} = \dfrac{d\vec{v}}{dt} = \ddot{x}\vec{i} + \ddot{y}\vec{j}$; speed $= |\vec{v}|$.
\item Speed $v$ on bearing $\theta$ converts to $\vec{v} = v\sin\theta\,\vec{i} + v\cos\theta\,\vec{j}$; conversely $v = \sqrt{v_1^2 + v_2^2}$ and the bearing comes from $\tan^{-1}|v_1/v_2|$ placed in the correct quadrant.
\item Relative velocity: $\vec{v}_{A/B} = \vec{v}_A - \vec{v}_B$ — the velocity of $A$ as seen by an observer moving with $B$. Geometrically it is the vector from the tip of $\vec{v}_B$ to the tip of $\vec{v}_A$.
\item $\vec{v}_{B/A} = -\vec{v}_{A/B}$: same magnitude, opposite direction.
\item In one dimension: subtract speeds for the same direction, add speeds for opposite directions.
\item Never subtract \emph{speeds} when the velocities are not parallel: subtract \emph{components}.
\end{itemize}
\end{aretenir}

\begin{exercice}[Practice]
\begin{enumerate}
\item A particle has position vector $\vec{r}(t) = (2t^3 - 5t)\vec{i} + (t^2 + 3)\vec{j}$ metres. Find its velocity and acceleration at $t = 2$, and its speed at that instant.
\item A helicopter flies at $120\ \mathrm{km\,h^{-1}}$ on a bearing of $150^\circ$. Express its velocity in $\vec{i}, \vec{j}$ form.
\item Car $A$ has velocity $18\vec{i} + 6\vec{j}\ \mathrm{m\,s^{-1}}$ and car $B$ has velocity $10\vec{i} + 14\vec{j}\ \mathrm{m\,s^{-1}}$. Find $\vec{v}_{A/B}$ and $\vec{v}_{B/A}$, giving the magnitude and bearing of each, and verify that they are opposite vectors.
\item Two buses leave Bafoussam on the same straight road in the same direction, one at $80\ \mathrm{km\,h^{-1}}$ and one at $65\ \mathrm{km\,h^{-1}}$. At what rate does the distance between them increase? What would the answer be if they travelled in opposite directions?
\item A ship sails with velocity $\vec{v}_S = 12\vec{i} + 5\vec{j}\ \mathrm{km\,h^{-1}}$. A current flows with velocity $\vec{v}_C = -3\vec{i} + 4\vec{j}\ \mathrm{km\,h^{-1}}$. Find the velocity of the ship relative to the current, and its magnitude and bearing.
\end{enumerate}
\end{exercice}

\begin{corrige}[Exercise 1]
\textbf{1.} $\vec{v} = (6t^2 - 5)\vec{i} + 2t\vec{j}$, so at $t = 2$: $\vec{v} = 19\vec{i} + 4\vec{j}$, speed $= \sqrt{361 + 16} = \sqrt{377} \approx 19.4\ \mathrm{m\,s^{-1}}$; $\vec{a} = 12t\,\vec{i} + 2\vec{j} = 24\vec{i} + 2\vec{j}\ \mathrm{m\,s^{-2}}$.

\textbf{2.} $\vec{v} = 120\sin 150^\circ\,\vec{i} + 120\cos 150^\circ\,\vec{j} = 60\vec{i} - 60\sqrt{3}\,\vec{j} \approx 60\vec{i} - 103.9\vec{j}\ \mathrm{km\,h^{-1}}$ (South-East quadrant, consistent with bearing $150^\circ$).

\textbf{3.} $\vec{v}_{A/B} = (18 - 10)\vec{i} + (6 - 14)\vec{j} = 8\vec{i} - 8\vec{j}$; magnitude $\sqrt{128} = 8\sqrt{2} \approx 11.3\ \mathrm{m\,s^{-1}}$; angle from North $\tan^{-1}(8/8) = 45^\circ$, South-East quadrant, bearing $135^\circ$. $\vec{v}_{B/A} = -8\vec{i} + 8\vec{j}$; same magnitude $8\sqrt{2}$, bearing $315^\circ$. Check: $\vec{v}_{A/B} + \vec{v}_{B/A} = \vec{0}$. \checkmark

\textbf{4.} Same direction: relative speed $= 80 - 65 = 15\ \mathrm{km\,h^{-1}}$, so the gap grows at $15\ \mathrm{km\,h^{-1}}$. Opposite directions: relative speed $= 80 + 65 = 145\ \mathrm{km\,h^{-1}}$.

\textbf{5.} $\vec{v}_{S/C} = \vec{v}_S - \vec{v}_C = (12 - (-3))\vec{i} + (5 - 4)\vec{j} = 15\vec{i} + \vec{j}\ \mathrm{km\,h^{-1}}$. Magnitude $= \sqrt{225 + 1} = \sqrt{226} \approx 15.03\ \mathrm{km\,h^{-1}}$. Angle from North: $\alpha = \tan^{-1}(15/1) \approx 86.2^\circ$, North-East quadrant, bearing $\approx 086^\circ$.
\end{corrige}

\coursec{Non-Parallel Forces, Resultant Velocity and True Velocity in River or Wind Problems}

In the previous section we established that velocity and acceleration are vectors, and we introduced the concept of \cle{relative velocity}: the velocity of one body as observed from another moving body. We now apply this vector machinery to one of the most concrete and examination-relevant situations in mechanics: a body moving \emph{through a medium that is itself moving}. A pirogue crossing the Wouri estuary, a ferry on the Sanaga, an aircraft flying from Douala to Yaound\'e through a crosswind — in each case the body's velocity relative to the ground is the \emph{vector sum} of two non-parallel velocities. Mastering this composition is the heart of this section.

\courssub{From Non-Parallel Forces to Non-Parallel Velocities}

You already know how to combine non-parallel \emph{forces} acting on a particle. The \cle{resultant} $\vec{R}$ of two forces $\vec{F}_1$ and $\vec{F}_2$ is obtained by the \cle{parallelogram law} (equivalently, the \cle{triangle law}): place the vectors head-to-tail; the resultant runs from the tail of the first to the head of the second. Its magnitude follows from the cosine rule,
\[
R^2 = F_1^2 + F_2^2 + 2F_1F_2\cos\alpha,
\]
where $\alpha$ is the angle between $\vec{F}_1$ and $\vec{F}_2$ placed tail-to-tail, and its direction is found via the sine rule or by resolving into perpendicular components.

The key observation is that \textbf{velocity obeys exactly the same vector algebra}. Nothing in the parallelogram law depends on the physical nature of the vectors — it is a geometric property of vectors themselves.

\begin{propriete}[Composition of Velocities (Relative Velocity)]
If a body $B$ has velocity $\vec{v}_{B/M}$ relative to a medium $M$, and the medium $M$ has velocity $\vec{v}_{M/G}$ relative to the ground $G$, then the velocity of the body relative to the ground (the \cle{true velocity} or \cle{resultant velocity}) is the vector sum
\[
\vec{v}_{B/G} = \vec{v}_{B/M} + \vec{v}_{M/G}.
\]
This is found by the parallelogram or triangle law, exactly as for forces. Resolving along perpendicular axes $Ox$ and $Oy$:
\[
v_x = v_{B/M,x} + v_{M/G,x}, \qquad v_y = v_{B/M,y} + v_{M/G,y},
\]
\[
|\vec{v}_{B/G}| = \sqrt{v_x^2 + v_y^2}, \qquad \tan\phi = \frac{v_y}{v_x} \quad (\phi \text{ measured from } Ox).
\]
\end{propriete}

\begin{definition}[True (Resultant) Velocity]
When a body moves through a moving medium (water current, wind), its \cle{true velocity} (also called \emph{resultant velocity} or \emph{velocity over ground}) is
\[
\vec{v}_{\text{true}} = \vec{v}_{\text{body}} + \vec{v}_{\text{medium}},
\]
where $\vec{v}_{\text{body}}$ (often denoted $\vec{v}_{B/M}$) is the velocity the body would have in a still medium — its \emph{heading} velocity — and $\vec{v}_{\text{medium}}$ ($\vec{v}_{M/G}$) is the velocity of the medium relative to the ground.
\end{definition}

\begin{attention}[Still-water speed is not ground speed]
A very common mistake is to quote the boat's still-water speed as its speed over the ground. When a current acts, the speed over ground is $|\vec{v}_{\text{boat}} + \vec{v}_{\text{current}}|$, which is generally \emph{different} from $|\vec{v}_{\text{boat}}|$. \textbf{Always draw the velocity triangle first.}
\end{attention}

\begin{savaistu}[Navigation on the Wouri]
The Wouri estuary, where Douala stands, experiences strong tidal currents that can reach $2\text{--}3\ \mathrm{m\,s^{-1}}$. Local pilots (often called \emph{bar pilots}) must calculate the true velocity of ships entering the port by combining the ship's engine speed (relative to water) with the tidal stream velocity (relative to the quay). A misjudgement of the vector sum can ground a vessel on the mudbanks — a real-life application of the triangle law you are learning.
\end{savaistu}

\courssub{River Problems: Heading, Current and True Velocity}

Fix the vocabulary carefully; GCE examiners test it directly. We adopt the standard relative velocity notation: $\vec{v}_{b/w}$ = velocity of boat relative to water, $\vec{v}_{w/g}$ = velocity of water (current) relative to ground, $\vec{v}_{b/g}$ = velocity of boat relative to ground (true velocity).

\begin{definition}[River Terminology]
For a boat crossing a river:
\begin{itemize}
\item the \cle{velocity of the boat in still water}, $\vec{v}_{b/w}$, is the velocity the engine and rudder produce relative to the water; its direction is the \emph{heading} of the boat;
\item the \cle{velocity of the current}, $\vec{v}_{w/g}$, is the velocity of the water relative to the banks, directed parallel to the banks;
\item the \cle{true velocity} is $\vec{v}_{b/g} = \vec{v}_{b/w} + \vec{v}_{w/g}$: the velocity of the boat as observed from the bank.
\end{itemize}
\end{definition}

Set up coordinates with $Ox$ pointing across the river (perpendicular to the banks) and $Oy$ pointing downstream. The river has width $d$, the boat's speed in still water is $u = |\vec{v}_{b/w}|$, and the current speed is $c = |\vec{v}_{w/g}|$.

\textbf{Case 1 — The quickest crossing.}

Suppose the pilot wants to reach the far bank in the \emph{shortest possible time}, and does not care where he lands. The displacement across the river is fixed at $d$; the crossing time is
\[
t = \frac{d}{\text{component of true velocity across the river}}.
\]
The current flows \emph{along} the river, so it has zero component across the river: it cannot help or hinder the crossing. The across-river component of the true velocity is exactly the across-river component of $\vec{v}_{b/w}$, namely $u\sin\beta$ where $\beta$ is the angle between the heading and the banks. This is maximised when $\beta = 90^\circ$, i.e. when the boat is steered \emph{perpendicular to the banks}.

\begin{methode}[Quickest Crossing of a River]
\begin{enumerate}
\item Steer the boat \textbf{perpendicular to the banks} (heading straight across).
\item The current does not affect the crossing time: $t_{\min} = \dfrac{d}{u}$.
\item During the crossing the current carries the boat downstream a distance $s = c\,t_{\min} = \dfrac{cd}{u}$.
\item The true speed over ground is $\sqrt{u^2 + c^2}$, directed at angle $\arctan(c/u)$ downstream of the straight-across line.
\end{enumerate}
\end{methode}

\begin{popfigure}
\begin{center}
\begin{tikzpicture}[scale=1.0, >=stealth]
% banks
\fill[popblueL] (0,0) rectangle (7,3);
\draw[thick, popdark] (0,0) -- (7,0);
\draw[thick, popdark] (0,3) -- (7,3);
\node[popdark] at (6.3,-0.25) {\scriptsize near bank};
\node[popdark] at (6.3,3.25) {\scriptsize far bank};
% current vector
\draw[->, very thick, popblue] (0.4,1.5) -- (2.0,1.5) node[midway, below, popblue] {\scriptsize current $\vec{v}_{w/g}$};
% boat vectors
\coordinate (A) at (1.0,0);
\draw[->, very thick, poporange] (A) -- ++(0,2.4) node[midway, left, poporange] {\scriptsize heading $\vec{v}_{b/w}$};
\draw[->, very thick, popgreen] (A) -- ++(1.6,0) node[midway, below, popgreen] {\scriptsize $\vec{v}_{w/g}$};
\draw[->, very thick, poppurple] (A) -- ++(1.6,2.4) node[midway, right, poppurple] {\scriptsize true $\vec{v}_{b/g}$};
% dashed construction
\draw[dashed, gray] (1.0,2.4) -- (2.6,2.4);
\draw[dashed, gray] (2.6,0) -- (2.6,2.4);
% width label
\draw[<->, popdark] (6.5,0) -- (6.5,3) node[midway, right] {\scriptsize width $d$};
% landing point
\fill[poppurple] (2.6,3) circle (2pt) node[above right, poppurple] {\scriptsize landing point};
\end{tikzpicture}
\end{center}
\legende{Quickest crossing: the boat is steered straight across; the current only causes downstream drift.}
\end{popfigure}

\textbf{Case 2 — Crossing to the point directly opposite (shortest path).}

Now the pilot insists on landing at the point $B$ directly opposite the starting point $A$. The true velocity $\vec{v}_{b/g}$ must then point \emph{straight across} the river, so the downstream component of the heading must exactly cancel the current. If the boat is steered upstream at angle $\theta$ to the straight-across direction (i.e. $\theta$ is the angle between $\vec{v}_{b/w}$ and the perpendicular to the banks), then
\[
u\sin\theta = c \qquad\Longrightarrow\qquad \sin\theta = \frac{c}{u}.
\]
The true speed across is then $u\cos\theta = \sqrt{u^2 - c^2}$, and the crossing time is
\[
t = \frac{d}{\sqrt{u^2 - c^2}}.
\]

\begin{methode}[Crossing to the Point Directly Opposite]
\begin{enumerate}
\item Check that $u > c$ (boat faster than current); otherwise the crossing is impossible.
\item Steer upstream at angle $\theta$ to the straight-across line, where $\sin\theta = \dfrac{c}{u}$.
\item The true velocity is perpendicular to the banks, of magnitude $\sqrt{u^2 - c^2}$.
\item Crossing time: $t = \dfrac{d}{\sqrt{u^2 - c^2}}$.
\end{enumerate}
\end{methode}

\begin{attention}[When direct crossing is impossible]
If $c \geq u$, the equation $\sin\theta = c/u$ has no solution: even pointing the boat fully upstream, the current matches or exceeds the boat's upstream component, and the boat can never make headway directly across. It is then impossible to land at the point directly opposite. \textbf{Never write $\theta = \arcsin(c/u)$ without first checking $c < u$.}
\end{attention}

\begin{popfigure}
\begin{center}
\begin{tikzpicture}[scale=1.0, >=stealth]
\fill[popblueL] (0,0) rectangle (7,3);
\draw[thick, popdark] (0,0) -- (7,0);
\draw[thick, popdark] (0,3) -- (7,3);
\draw[->, very thick, popblue] (5.2,1.5) -- (6.6,1.5) node[midway, below, popblue] {\scriptsize current $\vec{v}_{w/g}$};
\coordinate (A) at (3.5,0);
\coordinate (B) at (3.5,3);
\fill[popdark] (A) circle (2pt) node[below left, popdark] {\scriptsize $A$};
\fill[popdark] (B) circle (2pt) node[above left, popdark] {\scriptsize $B$};
% heading upstream
\draw[->, very thick, poporange] (A) -- ++(-1.15,2.4) node[midway, left, poporange] {\scriptsize heading $\vec{v}_{b/w}$};
% current
\draw[->, very thick, popgreen] (A) -- ++(1.15,0) node[midway, below, popgreen] {\scriptsize $\vec{v}_{w/g}$};
% true velocity straight across
\draw[->, very thick, poppurple] (A) -- ++(0,2.4) node[midway, right, poppurple] {\scriptsize true $\vec{v}_{b/g}$};
% angle theta
\draw[popdark] (3.5,1.0) arc (90:115.6:1.0);
\node[popdark] at (3.15,1.35) {\scriptsize $\theta$};
\draw[dashed, gray] (A) -- (B);
\end{tikzpicture}
\end{center}
\legende{Landing directly opposite: the boat is steered upstream at angle $\theta$ with $\sin\theta = c/u$, so the upstream component of the heading cancels the current.}
\end{popfigure}

\textbf{Case 3 — Shortest time versus shortest path.}

The two strategies answer two different optimisation questions, and a good candidate keeps them distinct:

\begin{center}
\begin{tabularx}{\linewidth}{|l|X|X|}
\hline
\rowcolor{popblueL} & \textbf{Quickest crossing} & \textbf{Directly opposite (shortest path)} \\
\hline
Heading & Straight across the river & Upstream at $\theta$, $\sin\theta = c/u$ \\
\hline
Time & $d/u$ (minimum possible) & $d/\sqrt{u^2-c^2}$ (longer) \\
\hline
Path over ground & Slanted downstream, length $d\sqrt{u^2+c^2}/u$ & Straight across, length $d$ (minimum) \\
\hline
Condition & Always possible & Possible only if $u > c$ \\
\hline
\end{tabularx}
\end{center}

Note that the quickest crossing always exists, while the shortest-path crossing requires $u > c$; and the shortest path takes \emph{longer} than the quickest crossing, because part of the boat's effort is spent fighting the current.

\begin{experience}[Paper Boat in a Tray]
Fill a long rectangular tray with water. Create a steady "current" by tilting the tray slightly or using a small pump at one end. Float a light paper boat powered by a rubber band or a small fan (constant thrust). 
\begin{enumerate}
\item Mark a starting point on one long side. Let the boat go with heading perpendicular to the banks. Measure crossing time and downstream drift.
\item Now adjust the heading upstream until the boat lands exactly opposite. Measure the angle and the crossing time.
\item Compare with the formulas $t_{\min}=d/u$ and $t=d/\sqrt{u^2-c^2}$.
\end{enumerate}
This simple experiment makes the vector addition $\vec{v}_{b/g}=\vec{v}_{b/w}+\vec{v}_{w/g}$ tangible.
\end{experience}

\begin{exoresolu}[Crossing the Sanaga]
A motorised pirogue can travel at $u = 5\ \mathrm{m\,s^{-1}}$ in still water. It crosses a stretch of the Sanaga river of width $d = 200\ \mathrm{m}$ where the current flows at $c = 3\ \mathrm{m\,s^{-1}}$.
\begin{enumerate}
\item Find the shortest possible crossing time and the downstream drift in that case.
\item Find the angle at which the boat must be steered to land directly opposite, and the crossing time then.
\end{enumerate}
\tcblower
\textbf{1. Quickest crossing.} Steer perpendicular to the banks. The current does not affect the across-river motion:
\[
t_{\min} = \frac{d}{u} = \frac{200}{5} = 40\ \mathrm{s}.
\]
Downstream drift: $s = c\,t_{\min} = 3 \times 40 = 120\ \mathrm{m}$. The true speed is $\sqrt{5^2+3^2} = \sqrt{34} \approx 5.83\ \mathrm{m\,s^{-1}}$.

\textbf{2. Directly opposite.} Since $u = 5 > c = 3$, a direct crossing is possible. Steer upstream at angle $\theta$ with
\[
\sin\theta = \frac{c}{u} = \frac{3}{5} = 0.6 \quad\Longrightarrow\quad \theta \approx 36.9^\circ.
\]
True speed across: $\sqrt{u^2 - c^2} = \sqrt{25-9} = 4\ \mathrm{m\,s^{-1}}$, so
\[
t = \frac{200}{4} = 50\ \mathrm{s}.
\]
As expected, the direct crossing ($50$ s) takes longer than the quickest crossing ($40$ s).
\end{exoresolu}

\courssub{Wind Problems: Aircraft, Track and Drift Angle}

Exactly the same mathematics governs an aircraft flying through moving air. We use the notation: $\vec{v}_{a/w}$ = velocity of aircraft relative to wind (heading + airspeed), $\vec{v}_{w/g}$ = wind velocity (velocity of air relative to ground), $\vec{v}_{a/g}$ = velocity of aircraft relative to ground (track velocity).

\begin{definition}[Wind Terminology]
For an aircraft:
\begin{itemize}
\item the \cle{velocity in still air} $\vec{v}_{a/w}$ (heading and airspeed) is the velocity relative to the air;
\item the \cle{wind velocity} $\vec{v}_{w/g}$ is the velocity of the air relative to the ground;
\item the \cle{true velocity} (or \emph{track velocity}) is $\vec{v}_{a/g} = \vec{v}_{a/w} + \vec{v}_{w/g}$: the velocity over the ground;
\item the \cle{drift angle} is the angle between the heading $\vec{v}_{a/w}$ and the track $\vec{v}_{a/g}$.
\end{itemize}
\end{definition}

Two standard inverse problems arise. 
\textbf{(a) Course to steer:} given the desired track $\vec{v}_{a/g}$ and the wind $\vec{v}_{w/g}$, find the required heading $\vec{v}_{a/w}$. 
Using $\vec{v}_{a/w} = \vec{v}_{a/g} - \vec{v}_{w/g}$: draw the track vector first; from its head, draw the wind vector \emph{reversed} (i.e. $-\vec{v}_{w/g}$); the vector from the tail of the track to the head of the reversed wind is the required heading.
\textbf{(b) Finding the wind:} given the heading $\vec{v}_{a/w}$ and the observed track $\vec{v}_{a/g}$, the wind is $\vec{v}_{w/g} = \vec{v}_{a/g} - \vec{v}_{a/w}$.

\begin{popfigure}
\begin{center}
\begin{tikzpicture}[scale=1.1, >=stealth]
\coordinate (O) at (0,0);
% heading
\draw[->, very thick, poporange] (O) -- (3.4,1.2) node[midway, above left, poporange] {\scriptsize heading $\vec{v}_{a/w}$};
% wind
\draw[->, very thick, popblue] (3.4,1.2) -- (4.2,2.6) node[midway, right, popblue] {\scriptsize wind $\vec{v}_{w/g}$};
% true velocity
\draw[->, very thick, poppurple] (O) -- (4.2,2.6) node[midway, above right, poppurple] {\scriptsize track $\vec{v}_{a/g}$};
% drift angle
\draw[popdark] (1.6,0.565) arc (19.4:31.7:1.7);
\node[popdark] at (2.15,0.85) {\scriptsize drift};
\end{tikzpicture}
\end{center}
\legende{Velocity triangle for an aircraft: the track (true velocity) is the sum of the heading and wind vectors; the drift angle lies between heading and track.}
\end{popfigure}

\begin{exoresolu}[Wind over the Littoral]
A small aircraft heads due \textbf{east} with airspeed $200\ \mathrm{km\,h^{-1}}$. A wind blows from the south at $40\ \mathrm{km\,h^{-1}}$ (i.e. the wind velocity points due \textbf{north}). Find the true velocity of the aircraft and the drift angle.
\tcblower
Take $Ox$ east, $Oy$ north. Then $\vec{v}_{a/w} = (200, 0)$ and $\vec{v}_{w/g} = (0, 40)$, so
\[
\vec{v}_{a/g} = (200, 40), \qquad |\vec{v}_{a/g}| = \sqrt{200^2 + 40^2} = \sqrt{41600} \approx 204.0\ \mathrm{km\,h^{-1}}.
\]
The drift angle $\delta$ satisfies $\tan\delta = \dfrac{40}{200} = 0.2$, giving $\delta \approx 11.3^\circ$ north of east. Although the nose points due east, the aircraft actually travels along a track $11.3^\circ$ north of east at about $204\ \mathrm{km\,h^{-1}}$.
\end{exoresolu}

\begin{exercice}[Ferry on the Wouri]
A ferry crosses a $600\ \mathrm{m}$ wide stretch of the Wouri river where the current flows at $1.5\ \mathrm{m\,s^{-1}}$. The ferry's speed in still water is $2.5\ \mathrm{m\,s^{-1}}$.
\begin{enumerate}
\item Calculate the minimum crossing time and the distance carried downstream.
\item Determine whether the ferry can reach the point directly opposite; if so, find the steering angle and the crossing time.
\item If the current doubled to $3\ \mathrm{m\,s^{-1}}$, could the ferry still land directly opposite? Justify.
\end{enumerate}
\end{exercice}

\begin{corrige}[Exercise: Ferry on the Wouri]
\textbf{1.} Quickest crossing: steer straight across, $t_{\min} = 600/2.5 = 240\ \mathrm{s}$. Drift: $s = 1.5 \times 240 = 360\ \mathrm{m}$.

\textbf{2.} Since $2.5 > 1.5$, direct crossing is possible: $\sin\theta = 1.5/2.5 = 0.6$, so $\theta \approx 36.9^\circ$ upstream. True speed $\sqrt{2.5^2 - 1.5^2} = 2\ \mathrm{m\,s^{-1}}$; time $= 600/2 = 300\ \mathrm{s}$.

\textbf{3.} With $c = 3 > u = 2.5$, $\sin\theta = 3/2.5 = 1.2 > 1$ has no solution: the ferry \textbf{cannot} land directly opposite; it would be swept downstream no matter how it is steered.
\end{corrige}

\begin{aretenir}[Key Points]
\begin{itemize}
\item True velocity: $\vec{v}_{b/g} = \vec{v}_{b/w} + \vec{v}_{w/g}$ (river) or $\vec{v}_{a/g} = \vec{v}_{a/w} + \vec{v}_{w/g}$ (air) — the same parallelogram law as for forces.
\item Quickest river crossing: steer perpendicular to the banks; $t_{\min} = d/u$; drift $= cd/u$.
\item Directly opposite (shortest path): steer upstream with $\sin\theta = c/u$; possible only if $u > c$; time $= d/\sqrt{u^2 - c^2}$.
\item Aircraft: track $=$ heading $+$ wind; drift angle lies between heading and track; wind $=$ track $-$ heading.
\item Always draw the velocity triangle, label vectors clearly ($\vec{v}_{b/w}$, $\vec{v}_{w/g}$, $\vec{v}_{b/g}$), and resolve into perpendicular components.
\end{itemize}
\end{aretenir}

\coursec{Applications of Relative and True Velocity}

In the previous section we learnt how to combine non-parallel velocities: a boat steered across a flowing river, an aircraft flying through a cross-wind, and the \emph{true} (resultant) velocity obtained by vector addition. In this section we put those tools to work on the situations the GCE examiner loves most: the \cle{apparent wind} felt by a moving cyclist, the \cle{apparent direction of rain}, river crossings with \cle{return journeys}, and problems where the current or wind is \emph{unknown} and must be deduced from observations. Everything reduces to one disciplined routine, which we state first and then practise until it becomes automatic.

\courssub{A Four-Step Problem-Solving Strategy}

Every problem in this chapter is a vector problem wearing a disguise. The disguise may be a cyclist, a canoe on the Wouri river, or rain on a taxi windscreen — underneath, it is always the same algebra.

\begin{methode}[The Four-Step Strategy]
\begin{enumerate}
\item \textbf{Choose axes.} Fix unit vectors $\vec{i}$ (East) and $\vec{j}$ (North). If the problem uses bearings, remember that a bearing $\theta$ is measured \emph{clockwise from North}, so a velocity of magnitude $v$ on bearing $\theta$ is
\[ \vec{v} = v\sin\theta\,\vec{i} + v\cos\theta\,\vec{j}. \]
\item \textbf{Write every velocity in components.} Translate each sentence of the question into a vector $a\vec{i}+b\vec{j}$. Never subtract magnitudes or directions directly.
\item \textbf{Form the required vector.} Resultant (true) velocity: add. Relative/apparent velocity: subtract, $\vec{v}_{A/B}=\vec{v}_A-\vec{v}_B$.
\item \textbf{Extract magnitude and direction.} Magnitude: $|\vec{v}|=\sqrt{a^{2}+b^{2}}$. Direction: convert back to a bearing or an angle using $\tan\theta = \dfrac{\text{east component}}{\text{north component}}$, placing the angle in the correct quadrant from the signs of the components.
\end{enumerate}
\end{methode}

\begin{attention}[Bearing versus Component Angle]
A bearing is measured \textbf{clockwise from North}, while the mathematical angle is measured anticlockwise from the $x$-axis. Mixing the two is the single most frequent source of error at the GCE. The safe recipe: bearing $\theta$ gives components $(v\sin\theta,\,v\cos\theta)$. Check: bearing $000^{\circ}$ (North) gives $(0,v)$ — pure $\vec{j}$. Bearing $090^{\circ}$ (East) gives $(v,0)$ — pure $\vec{i}$. 
\end{attention}

\courssub{Apparent Velocity: Wind and Rain as Seen by a Moving Observer}

\textbf{Intuition.} Stand still on a calm day and the air feels still. Now ride a bicycle at \SI{20}{km.h^{-1}} on that same calm day: you feel a wind of \SI{20}{km.h^{-1}} blowing straight into your face. That wind does not exist — it is the \emph{apparent} wind created by your own motion. If a real wind is also blowing, what you feel is the combination of the true wind and the wind due to your motion. Formally, the velocity of the air \emph{relative to you} is the velocity of the air minus your velocity.

\begin{definition}[Apparent Velocity of Wind or Rain]
Let $\vec{v}_{W}$ be the true velocity of the wind (or rain) relative to the ground, and $\vec{v}_{O}$ the velocity of the observer relative to the ground. The \cle{apparent velocity} of the wind (or rain) as experienced by the observer is the relative velocity
\[ \vec{v}_{W/O} \;=\; \vec{v}_{W} - \vec{v}_{O}. \]
The observer feels the wind coming \emph{from} the direction \emph{opposite} to $\vec{v}_{W/O}$.
\end{definition}

\begin{popfigure}
\begin{center}
\begin{tikzpicture}[scale=0.9,>=stealth]
\draw[->,thick,popblue] (0,0) -- (3,0) node[midway,below]{$\vec{v}_{O}$ (cyclist)};
\draw[->,thick,popgreen] (0,0) -- (-1.5,2.2) node[midway,above left]{$\vec{v}_{W}$ (true wind)};
\draw[->,thick,poporange] (0,0) -- (-4.5,2.2) node[midway,above]{$\vec{v}_{W/O}$ (apparent)};
\draw[->,dashed,popmuted] (3,0) -- (-4.5,2.2) node[midway,below right]{$\vec{v}_{W}-\vec{v}_{O}$};
\fill (0,0) circle (2pt) node[below left]{observer};
\end{tikzpicture}
\end{center}
\legende{The apparent wind $\vec{v}_{W/O}=\vec{v}_{W}-\vec{v}_{O}$ closes the triangle from the tip of $\vec{v}_{O}$ to the tip of $\vec{v}_{W}$.}
\end{popfigure}

\begin{exoresolu}[Apparent Wind on a Cyclist]
A cyclist rides due East at \SI{6}{m.s^{-1}}. The true wind blows from the South-West at \SI{5\sqrt{2}}{m.s^{-1}} (i.e.\ towards the North-East). Determine the apparent wind velocity experienced by the cyclist, and state the direction from which the wind appears to blow.
\tcblower
\textbf{Step 1 — axes.} $\vec{i}$ East, $\vec{j}$ North.

\textbf{Step 2 — components.} Cyclist: $\vec{v}_{O}=6\vec{i}$. Wind blows \emph{towards} NE, i.e.\ on bearing $045^{\circ}$:
\[ \vec{v}_{W}=5\sqrt{2}\sin 45^{\circ}\,\vec{i}+5\sqrt{2}\cos 45^{\circ}\,\vec{j}=5\vec{i}+5\vec{j}. \]

\textbf{Step 3 — relative vector.}
\[ \vec{v}_{W/O}=\vec{v}_{W}-\vec{v}_{O}=(5-6)\vec{i}+5\vec{j}=-\vec{i}+5\vec{j}. \]

\textbf{Step 4 — magnitude and direction.}
\[ |\vec{v}_{W/O}|=\sqrt{1+25}=\sqrt{26}\approx \SI{5.1}{m.s^{-1}}. \]
The apparent wind vector points West of North with $\tan\alpha=\dfrac{1}{5}$, so $\alpha\approx 11.3^{\circ}$; the apparent wind blows \emph{towards} bearing $360^{\circ}-11.3^{\circ}=348.7^{\circ}$, hence it appears to come \emph{from} the opposite direction, bearing $168.7^{\circ}$ — nearly from the South-South-East, slightly ahead of the cyclist's right shoulder.
\end{exoresolu}

\begin{attention}[Sign Errors When Moving Towards the Wind]
When the observer moves \emph{towards} the wind, the apparent speed is the \emph{sum} of the magnitudes; when moving \emph{away}, it is the difference. Students who subtract magnitudes blindly get this wrong half the time. The cure is infallible: \textbf{always write both vectors in components first}, then subtract component by component. The signs take care of themselves.
\end{attention}

\courssub{Rain Problems}

Rain problems are identical in structure to wind problems, with one simplification: true rain falls vertically, so $\vec{v}_{R}=-v\,\vec{j}$ where $v$ is the (unknown) speed of fall. A moving observer sees the rain slanting; the angle of slant reveals the true speed.

\begin{popfigure}
\begin{center}
\begin{tikzpicture}[scale=1.0,>=stealth]
\draw[->,thick,popblue] (0,0) -- (2.5,0) node[midway,below]{$\vec{v}_{O}$ (observer)};
\draw[->,thick,popgreen] (4,3) -- (4,0.6) node[midway,right]{true rain $\vec{v}_{R}$};
\draw[->,thick,poporange] (1.5,3) -- (4,0.6) node[midway,above left]{apparent rain};
\draw[dashed,popmuted] (4,3) -- (1.5,3);
\draw (4,2.2) arc[start angle=90,end angle=124, radius=1.6];
\node at (3.35,2.75){$\theta$};
\fill (4,0.6) circle (2pt) node[below]{observer's eye};
\end{tikzpicture}
\end{center}
\legende{To the moving observer the rain appears to fall along $\vec{v}_{R}-\vec{v}_{O}$, tilted through an angle $\theta$ from the vertical, with $\tan\theta=\dfrac{|\vec{v}_{O}|}{|\vec{v}_{R}|}$.}
\end{popfigure}

\begin{propriete}[Apparent Angle of Rain]
If rain falls vertically with speed $u$ and an observer moves horizontally with speed $v$, the rain appears to fall at an angle $\theta$ to the vertical where
\[ \tan\theta=\frac{v}{u}. \]
Conversely, measuring $\theta$ for two different speeds of the observer determines both $u$ and the true direction of the rain. (The derivation is immediate from the components of $\vec{v}_{R/O}$; it is a standard bookwork result and may be quoted.)
\end{propriete}

\begin{exoresolu}[Finding the True Velocity of Rain]
To a man walking due East at \SI{4}{km.h^{-1}}, rain appears to fall vertically. When he doubles his speed to \SI{8}{km.h^{-1}} (still due East), the rain appears to fall at $30^{\circ}$ to the vertical, hitting him from ahead. Find the true velocity of the rain.
\tcblower
\textbf{Step 1–2.} Let the true rain velocity be $\vec{v}_{R}=a\vec{i}+b\vec{j}$ with $b<0$.

At speed $4$: $\vec{v}_{R/O}=(a-4)\vec{i}+b\vec{j}$. This is vertical, so its $\vec{i}$-component vanishes: $a-4=0$, hence $a=4$.

At speed $8$: $\vec{v}_{R/O}=(4-8)\vec{i}+b\vec{j}=-4\vec{i}+b\vec{j}$. The angle to the vertical is $30^{\circ}$, so
\[ \tan 30^{\circ}=\frac{|\text{horizontal}|}{|\text{vertical}|}=\frac{4}{|b|}
\quad\Longrightarrow\quad |b|=\frac{4}{\tan 30^{\circ}}=4\sqrt{3}. \]

\textbf{Step 3–4.} $\vec{v}_{R}=4\vec{i}-4\sqrt{3}\,\vec{j}$, giving
\[ |\vec{v}_{R}|=\sqrt{16+48}=\SI{8}{km.h^{-1}},\qquad \tan\varphi=\frac{4}{4\sqrt{3}}=\frac{1}{\sqrt{3}}\ \Rightarrow\ \varphi=30^{\circ}. \]
The rain truly falls at \SI{8}{km.h^{-1}} at $30^{\circ}$ to the vertical, blowing towards the East (from the West). Note the elegance: the first observation told us the rain's horizontal drift matches the walker at \SI{4}{km.h^{-1}}.
\end{exoresolu}

\courssub{River Problems: Return Journeys and Average Speed}

A canoe crossing a river with current speed $c$, steered at speed $u$ relative to the water, has been treated in Section 2. The GCE frequently extends this to a \cle{round trip}: downstream then back upstream (or across and back). The key warning is that the average speed is \emph{not} $u$.

\begin{propriete}[Round Trip Along the Stream]
For a journey of distance $d$ each way, with boat speed $u$ in still water and current $c$ ($u>c$):
\[ t_{\text{down}}=\frac{d}{u+c},\qquad t_{\text{up}}=\frac{d}{u-c},\qquad
T=\frac{d}{u+c}+\frac{d}{u-c}=\frac{2ud}{u^{2}-c^{2}}. \]
The average speed for the round trip is
\[ \bar{v}=\frac{2d}{T}=\frac{u^{2}-c^{2}}{u}=u\left(1-\frac{c^{2}}{u^{2}}\right)<u. \]
\end{propriete}

\textbf{Why this matters:} the current \emph{always} reduces the average speed on a round trip, even though it helps on the way down. The time lost fighting the current upstream exceeds the time gained drifting with it, because the slow leg lasts longer. The factor $1-\frac{c^{2}}{u^{2}}$ quantifies the loss.

\begin{exoresolu}[Round Trip on the River]
A motorboat travels at \SI{10}{km.h^{-1}} in still water. It makes a trip of \SI{12}{km} downstream and back on a river flowing at \SI{2}{km.h^{-1}}. Find the total time and the average speed for the round trip, and compare with the time in still water.
\tcblower
Downstream speed $=10+2=12$ km/h, so $t_{\text{down}}=\dfrac{12}{12}=1$ h.

Upstream speed $=10-2=8$ km/h, so $t_{\text{up}}=\dfrac{12}{8}=1.5$ h.

Total time $T=2.5$ h for $24$ km, giving $\bar{v}=\dfrac{24}{2.5}=\SI{9.6}{km.h^{-1}}$.

Check with the formula: $\bar{v}=10\left(1-\dfrac{4}{100}\right)=\SI{9.6}{km.h^{-1}}$. \checkmark

In still water the trip would take $\dfrac{24}{10}=2.4$ h. The current costs the boatman $0.1$ h $=6$ minutes, and the average speed drops by $4\%$ — exactly $\left(\frac{c}{u}\right)^{2}=\left(\frac{2}{10}\right)^{2}$.
\end{exoresolu}

\courssub{Deducing an Unknown Current or Wind}

Sometimes the current or wind is the \emph{unknown}: we observe the true (resultant) velocity and the heading, and must work backwards. The strategy is unchanged — components first — but now we solve equations instead of evaluating them.

\begin{exoresolu}[Unknown Current from an Observed Track]
A boat is steered due North at \SI{8}{km.h^{-1}} in still water. It is actually observed to travel on a bearing of $030^{\circ}$ at \SI{10}{km.h^{-1}}. Find the velocity of the current.
\tcblower
\textbf{Step 2 — components.} Boat relative to water: $\vec{v}_{B/W}=8\vec{j}$. Observed true velocity:
\[ \vec{v}_{B}=10\sin 30^{\circ}\,\vec{i}+10\cos 30^{\circ}\,\vec{j}=5\vec{i}+5\sqrt{3}\,\vec{j}. \]
\textbf{Step 3 — solve.} Since $\vec{v}_{B}=\vec{v}_{B/W}+\vec{v}_{C}$,
\[ \vec{v}_{C}=\vec{v}_{B}-\vec{v}_{B/W}=5\vec{i}+(5\sqrt{3}-8)\vec{j}\approx 5\vec{i}+0.66\vec{j}. \]
\textbf{Step 4.} $|\vec{v}_{C}|=\sqrt{25+(5\sqrt{3}-8)^{2}}=\sqrt{25+0.436}\approx \SI{5.04}{km.h^{-1}}$, with $\tan\theta=\dfrac{5}{0.66}$, giving a bearing of about $082.5^{\circ}$. The current flows almost due East at about \SI{5.0}{km.h^{-1}}.
\end{exoresolu}

\courssub{Mixed GCE-Style Practice}

\begin{exemplebox}[Exam Technique: Read, Convert, Compute, Conclude]
GCE questions on this topic are worth 6–10 marks and are marked methodically: one mark for correct components, one for the correct relative/resultant vector, one for the magnitude, one for the direction \emph{with correct quadrant}. Even if your arithmetic slips, a clearly set-up vector equation earns most of the marks. Always finish with a sentence: ``The wind appears to blow from bearing \ldots at \ldots m/s.''
\end{exemplebox}

\begin{exercice}[Guided Practice]
\begin{enumerate}
\item A cyclist travels due South at \SI{7}{m.s^{-1}}. The true wind blows from the East at \SI{4}{m.s^{-1}}. Find the apparent wind velocity and the bearing from which it appears to come.
\item To a man running due West at \SI{3}{m.s^{-1}}, rain appears to fall at $45^{\circ}$ to the vertical. Find the true speed of the (vertically falling) rain.
\item A boat with still-water speed \SI{15}{km.h^{-1}} travels \SI{20}{km} upstream and returns, on a river flowing at \SI{3}{km.h^{-1}}. Find the total time and the average speed for the round trip.
\item An aircraft heads due East with airspeed \SI{400}{km.h^{-1}}. Its true track is bearing $080^{\circ}$ at \SI{420}{km.h^{-1}}. Find the wind velocity.
\end{enumerate}
\end{exercice}

\begin{corrige}[Exercise 1]
\textbf{1.} $\vec{v}_{O}=-7\vec{j}$; wind blows towards West: $\vec{v}_{W}=-4\vec{i}$. Apparent: $\vec{v}_{W/O}=-4\vec{i}+7\vec{j}$; magnitude $\sqrt{65}\approx \SI{8.06}{m.s^{-1}}$; it blows towards bearing $\arctan(4/7)\approx 29.7^{\circ}$ West of North, i.e.\ bearing $330.3^{\circ}$, so it appears to come \emph{from} bearing $150.3^{\circ}$.

\textbf{2.} $\vec{v}_{R}=-u\vec{j}$, $\vec{v}_{O}=-3\vec{i}$; $\vec{v}_{R/O}=3\vec{i}-u\vec{j}$; $\tan 45^{\circ}=\frac{3}{u}=1$, so $u=\SI{3}{m.s^{-1}}$.

\textbf{3.} $t=\dfrac{20}{12}+\dfrac{20}{18}=\dfrac{5}{3}+\dfrac{10}{9}=\dfrac{25}{9}\approx \SI{2.78}{h}$; $\bar{v}=\dfrac{40}{25/9}=\dfrac{72}{5}=\SI{14.4}{km.h^{-1}} =15\left(1-\frac{9}{225}\right)$. \checkmark

\textbf{4.} $\vec{v}_{P}=420\sin 80^{\circ}\vec{i}+420\cos 80^{\circ}\vec{j}\approx 413.6\vec{i}+72.9\vec{j}$; $\vec{v}_{P/A}=400\vec{i}$; wind $\vec{v}_{W}=\vec{v}_{P}-\vec{v}_{P/A}\approx 13.6\vec{i}+72.9\vec{j}$; $|\vec{v}_{W}|\approx \SI{74.2}{km.h^{-1}}$ on bearing $\arctan(13.6/72.9)\approx 010.6^{\circ}$.
\end{corrige}

\begin{aretenir}[Key Points of This Section]
\begin{itemize}
\item One strategy solves everything: axes $\to$ components $\to$ add (resultant) or subtract (relative/apparent) $\to$ magnitude and bearing.
\item Apparent wind or rain: $\vec{v}_{\text{apparent}}=\vec{v}_{\text{true}}-\vec{v}_{\text{observer}}$; the wind appears to come \emph{from} the direction opposite to this vector.
\item Vertical rain seen at angle $\theta$ by an observer moving at speed $v$: $\tan\theta=v/u$.
\item On a round trip the current always lowers the average speed: $\bar{v}=u\left(1-\frac{c^{2}}{u^{2}}\right)$.
\item Unknown current or wind: subtract the known velocity from the observed true velocity, in components.
\end{itemize}
\end{aretenir}

\begin{savaistu}[The Michelson–Morley Analogy]
The formula $\bar{v}=u\left(1-\frac{c^{2}}{u^{2}}\right)$ for a round trip in a current is exactly the calculation Michelson and Morley used in their famous experiment, treating light as a boat and the hypothetical ``aether wind'' as the current. The expected tiny slowing was never detected — a null result that helped open the road to Einstein's relativity. Your river-crossing algebra is, quite literally, the mathematics of one of the most important experiments in physics.
\end{savaistu}

\coursec{Interception and Collision: Concept and Mathematical Treatment}

In the previous section we used relative velocity to answer questions such as ``how fast does ship $B$ appear to move as seen from ship $A$?'' We now turn that idea into a powerful weapon. Suppose a coastguard launch leaves Limbe to intercept a fishing vessel, or a goalkeeper at the Amadou Ahidjo stadium must decide where to run to catch a through ball. In each case one body must \emph{steer} so as to arrive at exactly the same point as a moving target at exactly the same instant. This is the problem of \cle{interception}, and when neither body wishes the meeting, we call it a \cle{collision}. The whole theory rests on one simple idea: \emph{work in the reference frame of the target}.

\courssub{The Concept of Interception and Collision}

Imagine two ships $P$ and $Q$ moving on the open sea with constant velocities $\vec{v}_P$ and $\vec{v}_Q$. Ship $P$ is the \cle{interceptor} and ship $Q$ is the \cle{target}. We say that $P$ \emph{intercepts} $Q$ if there exists a time $t > 0$ at which both ships occupy the same point.

\begin{definition}[Interception and collision]
Two bodies $P$ and $Q$, with position vectors $\vec{r}_P(t)$ and $\vec{r}_Q(t)$ measured from a fixed origin $O$, \textbf{intercept} (or \textbf{collide}) at time $t$ if and only if
\[ \vec{r}_P(t) = \vec{r}_Q(t) \qquad \text{for some } t > 0. \]
Equivalently, the relative position $\vec{r}_{P/Q}(t) = \vec{r}_P(t) - \vec{r}_Q(t)$ vanishes: $\vec{r}_{P/Q}(t) = \vec{0}$. The word \emph{interception} is used when the meeting is intended; \emph{collision} when it is not. Mathematically they are the same event.
\end{definition}

The key intuition is this: \textbf{stand on the target}. From the target's frame, the target itself is stationary, and the interceptor moves with the relative velocity $\vec{v}_{P/Q} = \vec{v}_P - \vec{v}_Q$. For the interceptor to reach the target, it must travel in this frame along the straight line joining the two initial positions. This gives the fundamental condition.

\begin{propriete}[Condition for interception]
Let $P$ and $Q$ move with constant velocities, and let $\vec{r}_{P/Q}(0) = \vec{r}_P(0) - \vec{r}_Q(0)$ be the initial position of $P$ relative to $Q$. Interception occurs if and only if the relative velocity $\vec{v}_{P/Q}$ is directed along the line joining the initial positions, pointing from $P$ towards $Q$:
\[ \vec{v}_{P/Q} = -\lambda\, \vec{r}_{P/Q}(0) \quad \text{for some } \lambda > 0, \]
i.e.\ $\vec{v}_{P/Q}$ is parallel to $\overrightarrow{QP}(0)$ (the vector from $Q$ to $P$) — $P$ closes in on $Q$. The interception time is then $t = \dfrac{|\vec{r}_{P/Q}(0)|}{|\vec{v}_{P/Q}|}$.
\end{propriete}

\emph{Proof (instructive, not required for the GCE).} With constant velocities, $\vec{r}_{P/Q}(t) = \vec{r}_{P/Q}(0) + t\,\vec{v}_{P/Q}$. Interception requires $\vec{r}_{P/Q}(t) = \vec{0}$, i.e.\ $t\,\vec{v}_{P/Q} = -\vec{r}_{P/Q}(0)$. Since $t > 0$ is a scalar, $\vec{v}_{P/Q}$ must be a negative scalar multiple of $\vec{r}_{P/Q}(0)$, which is exactly the stated condition. Dividing magnitudes gives $t = |\vec{r}_{P/Q}(0)|/|\vec{v}_{P/Q}|$. \hfill$\blacksquare$

\begin{attention}[Aiming at the target's present position]
The most frequent error is to steer the interceptor \emph{directly at where the target is now}. By the time the interceptor arrives there, the target has moved on! You must aim at the \textbf{future meeting point} — equivalently, choose the course so that the \emph{relative} velocity points along the initial line joining the ships. In the target's frame, the interceptor's path must be a straight line onto the stationary target.
\end{attention}

\courssub{Setting up Interception with Components}

In examination questions positions and velocities are given in components with unit vectors $\vec{i}$ (east) and $\vec{j}$ (north). For constant velocities,
\[ \vec{r}_P(t) = \vec{r}_P(0) + t\,\vec{v}_P, \qquad \vec{r}_Q(t) = \vec{r}_Q(0) + t\,\vec{v}_Q. \]
The interception condition $\vec{r}_P(t) = \vec{r}_Q(t)$ is a \emph{vector} equation, so it yields \textbf{two simultaneous scalar equations} (one from the $\vec{i}$-components, one from the $\vec{j}$-components). Typically the unknowns are the time $t$ and one steering parameter (the direction of $\vec{v}_P$, or one component of it).

\begin{methode}[Component method for interception]
\begin{enumerate}
\item Write down $\vec{r}_P(0)$, $\vec{r}_Q(0)$, $\vec{v}_Q$ and $\vec{v}_P$ (with the unknown steering angle or component left as a letter).
\item Form $\vec{r}_P(t) = \vec{r}_P(0) + t\vec{v}_P$ and $\vec{r}_Q(t) = \vec{r}_Q(0) + t\vec{v}_Q$.
\item Equate the $\vec{i}$-components and the $\vec{j}$-components: two simultaneous equations.
\item Solve for $t$ and the steering parameter.
\item \textbf{Check that $t > 0$}; a negative time means interception already happened in the past — the chosen course is impossible.
\item Substitute $t$ back to find the position vector of the interception point.
\end{enumerate}
\end{methode}

\begin{attention}[Forgetting to check $t > 0$]
The algebra often produces two mathematical solutions, or a solution with $t < 0$. A negative time is physically meaningless for a problem starting at $t = 0$: it would mean the ships met \emph{before} the motion started. Always verify the sign of $t$ and state explicitly that the interception is possible.
\end{attention}

\courssub{The Course to Steer: Velocity Triangle Method}

When the interceptor's \emph{speed} is fixed but its \emph{direction} is free, the neatest construction is the \cle{velocity triangle}. We require
\[ \vec{v}_{P/Q} = \vec{v}_P - \vec{v}_Q \quad \text{to be parallel to } \overrightarrow{QP}(0). \]
Draw the target's velocity $\vec{v}_Q$; from its tip, the interceptor's velocity $\vec{v}_P$ (of known magnitude $u$) must be chosen so that the closing side $\vec{v}_{P/Q}$ lies along the initial line $QP$. Geometrically: draw a ray from the tip of $\vec{v}_Q$ in the direction from $Q$ to $P$; the interceptor's velocity vector must end on this ray, with $|\vec{v}_P| = u$. A circle of radius $u$ centred at the tail of $\vec{v}_P$ cuts the ray at the required point(s). There may be two, one, or zero intersections — corresponding to two possible courses, a unique grazing course, or impossibility.

\begin{popfigure}
\begin{center}
\begin{tikzpicture}[scale=0.9, >=stealth]
% initial positions
\coordinate (P) at (0,0);
\coordinate (Q) at (6,2);
\coordinate (I) at (3.6,3.4);
\fill[popblue] (P) circle (2.2pt) node[below left, popdark] {$P(0)$ interceptor};
\fill[poporange] (Q) circle (2.2pt) node[below right, popdark] {$Q(0)$ target};
\draw[dashed, popmuted] (P) -- (Q) node[midway, below, sloped, popmuted] {initial line $PQ$};
% target course
\draw[->, very thick, poporange] (Q) -- ++(1.6,2.4) node[right, popdark] {$\vec{v}_Q$};
\draw[dashed, poporange] (Q) -- (I);
% interceptor course
\draw[->, very thick, popblue] (P) -- (I) node[midway, above left, popdark] {$\vec{v}_P$ (course to steer)};
\fill[poppurple] (I) circle (2.5pt) node[above, popdark] {interception point};
% velocity triangle inset
\begin{scope}[shift={(8.2,-0.4)}, scale=0.85]
\coordinate (A) at (0,0);
\coordinate (B) at (1.6,2.4);
\coordinate (C) at (-1.1,2.9);
\draw[->, very thick, poporange] (A) -- (B) node[midway, right, popdark] {$\vec{v}_Q$};
\draw[->, very thick, popblue] (A) -- (C) node[midway, left, popdark] {$\vec{v}_P$};
\draw[->, very thick, poppurple] (B) -- (C) node[midway, above, popdark] {$\vec{v}_{P/Q}$};
\draw[dashed, popmuted] (C) -- ++(1.6,-1.2);
\node[popmuted] at (0.6,-0.5) {$\vec{v}_{P/Q} \parallel \overrightarrow{QP}$};
\end{scope}
\end{tikzpicture}
\legende{The interceptor $P$ steers a course such that the relative velocity $\vec{v}_{P/Q}$ lies along the initial line joining the ships; the velocity triangle (right) fixes the direction of $\vec{v}_P$.}
\end{center}
\end{popfigure}

\begin{exoresolu}[Course to steer for interception]
At noon, a patrol boat $P$ is at the origin and a trawler $Q$ is at the point with position vector $(20\vec{i} + 12\vec{j})\ \mathrm{km}$. The trawler steams with constant velocity $(6\vec{i} + 8\vec{j})\ \mathrm{km\,h^{-1}}$. The patrol boat can travel at $26\ \mathrm{km\,h^{-1}}$. Find the velocity the patrol boat must steer to intercept the trawler, and the time of interception.
\tcblower
\textbf{Solution.} Let the patrol boat's velocity be $\vec{v}_P = (a\vec{i} + b\vec{j})\ \mathrm{km\,h^{-1}}$ with $a^2 + b^2 = 26^2 = 676$.

Position vectors at time $t$ hours:
\[ \vec{r}_P(t) = t(a\vec{i} + b\vec{j}), \qquad \vec{r}_Q(t) = (20 + 6t)\vec{i} + (12 + 8t)\vec{j}. \]
Interception: $\vec{r}_P(t) = \vec{r}_Q(t)$, giving
\[ at = 20 + 6t \quad (1), \qquad bt = 12 + 8t \quad (2). \]
From (1): $(a-6)t = 20$; from (2): $(b-8)t = 12$. Hence
\[ \frac{a-6}{b-8} = \frac{20}{12} = \frac{5}{3} \;\Longrightarrow\; 3(a-6) = 5(b-8) \;\Longrightarrow\; 3a - 5b = -22. \]
So $a = \dfrac{5b - 22}{3}$. Substituting into $a^2 + b^2 = 676$:
\[ \left(\frac{5b-22}{3}\right)^2 + b^2 = 676 \;\Longrightarrow\; 25b^2 - 220b + 484 + 9b^2 = 6084 \]
\[ 34b^2 - 220b - 5600 = 0 \;\Longrightarrow\; 17b^2 - 110b - 2800 = 0. \]
The discriminant is $\Delta = 110^2 + 4 \times 17 \times 2800 = 12100 + 190400 = 202500 = 450^2$.
\[ b = \frac{110 \pm 450}{34}. \]
Taking the positive root (steering north-east towards the trawler): $b = \dfrac{560}{34} = \dfrac{280}{17} \approx 16.47$, then $a = \dfrac{5(280/17)-22}{3} = \dfrac{342}{17} \approx 20.12$. Check: $(342/17)^2 + (280/17)^2 = 676$. \checkmark

Then $t = \dfrac{20}{a-6} = \dfrac{20}{342/17 - 102/17} = \dfrac{20}{240/17} = \dfrac{17}{12} \approx 1.417\ \mathrm{h} > 0$. \checkmark

\textbf{Answer:} $\vec{v}_P = \left(\dfrac{342}{17}\vec{i} + \dfrac{280}{17}\vec{j}\right)\ \mathrm{km\,h^{-1}} \approx (20.1\vec{i} + 16.5\vec{j})\ \mathrm{km\,h^{-1}}$, interception after $\dfrac{17}{12}\ \mathrm{h}$ ($1\ \mathrm{h}\ 25\ \mathrm{min}$), at position $\vec{r} = \left(\dfrac{57}{2}\vec{i} + \dfrac{70}{3}\vec{j}\right)\ \mathrm{km} \approx (28.5\vec{i} + 23.3\vec{j})\ \mathrm{km}$.
\end{exoresolu}

\begin{popfigure}
\begin{center}
\begin{tikzpicture}[scale=0.16, >=stealth]
\draw[->, gray] (0,0) -- (34,0) node[right] {$\vec{i}$ (km)};
\draw[->, gray] (0,0) -- (0,30) node[above] {$\vec{j}$ (km)};
\coordinate (P0) at (0,0);
\coordinate (Q0) at (20,12);
\coordinate (M) at (28.5,23.333);
\fill[popblue] (P0) circle (0.5) node[below left, popdark] {$P$, $t=0$};
\fill[poporange] (Q0) circle (0.5) node[below right, popdark] {$Q$, $t=0$};
\draw[->, very thick, popblue] (P0) -- (M) node[midway, above left, popdark] {$\vec{r}_P(t)$};
\draw[->, very thick, poporange] (Q0) -- (M) node[midway, below right, popdark] {$\vec{r}_Q(t)$};
\fill[poppurple] (M) circle (0.6) node[above, popdark] {$t = 17/12$ h};
\draw[dashed, popmuted] (P0) -- (Q0);
\end{tikzpicture}
\legende{Position vectors of the patrol boat and the trawler meeting at the interception point; both arrive at the same instant $t = 17/12$ h.}
\end{center}
\end{popfigure}

\courssub{Collision of Two Bodies Moving with Constant Velocities}

When \emph{both} velocities are given and fixed, the question becomes: \emph{will they collide?} By the property above, we simply test whether $\vec{v}_{P/Q}$ is parallel to the initial relative position $\vec{r}_{P/Q}(0)$ and directed towards $Q$.

\begin{methode}[Testing for collision]
\begin{enumerate}
\item Compute $\vec{r}_{P/Q}(0) = \vec{r}_P(0) - \vec{r}_Q(0)$ and $\vec{v}_{P/Q} = \vec{v}_P - \vec{v}_Q$.
\item Check whether $\vec{v}_{P/Q} = -\lambda\,\vec{r}_{P/Q}(0)$ for some $\lambda > 0$ (compare ratios of components).
\item If yes: collision at $t = 1/\lambda$ (equivalently $t = |\vec{r}_{P/Q}(0)|/|\vec{v}_{P/Q}|$). If no: no collision — proceed to closest approach.
\end{enumerate}
\end{methode}

\begin{exemplebox}[Testing a collision course]
Ship $A$ starts at $(2\vec{i} + 14\vec{j})\ \mathrm{km}$ with velocity $(4\vec{i} - 2\vec{j})\ \mathrm{km\,h^{-1}}$; ship $B$ starts at $(14\vec{i} + 2\vec{j})\ \mathrm{km}$ with velocity $(-2\vec{i} + 4\vec{j})\ \mathrm{km\,h^{-1}}$. Then
\[ \vec{r}_{A/B}(0) = (-12\vec{i} + 12\vec{j})\ \mathrm{km}, \qquad \vec{v}_{A/B} = (6\vec{i} - 6\vec{j})\ \mathrm{km\,h^{-1}} = -\tfrac{1}{2}\,\vec{r}_{A/B}(0). \]
Since $\vec{v}_{A/B}$ is a \emph{negative} scalar multiple of $\vec{r}_{A/B}(0)$, the ships collide, at $t = 2\ \mathrm{h}$, at the point $\vec{r}_A(2) = (10\vec{i} + 10\vec{j})\ \mathrm{km}$.
\end{exemplebox}

\courssub{Closest Distance of Approach}

If the collision test fails, the bodies pass each other at some \cle{closest distance of approach}. The separation is
\[ d(t) = \left|\vec{r}_{P/Q}(t)\right| = \left|\vec{r}_{P/Q}(0) + t\,\vec{v}_{P/Q}\right|. \]
Since $d(t) \geq 0$, minimising $d(t)$ is equivalent to minimising $D(t) = d(t)^2$, which is a \emph{quadratic in $t$} — so no heavy calculus is needed: complete the square, or differentiate.

\begin{methode}[Finding the closest approach]
\begin{enumerate}
\item Write $\vec{r}_{P/Q}(t) = (x_0 + v_x t)\vec{i} + (y_0 + v_y t)\vec{j}$.
\item Form $D(t) = (x_0 + v_x t)^2 + (y_0 + v_y t)^2$.
\item Either differentiate: $D'(t) = 0$ gives $t_{\min} = -\dfrac{x_0 v_x + y_0 v_y}{v_x^2 + v_y^2} = -\dfrac{\vec{r}_{P/Q}(0) \cdot \vec{v}_{P/Q}}{|\vec{v}_{P/Q}|^2}$; or complete the square.
\item Check $t_{\min} > 0$ (otherwise the closest approach occurred before $t = 0$).
\item Compute $d_{\min} = \sqrt{D(t_{\min})}$.
\end{enumerate}
\end{methode}

\begin{propriete}[Closest approach — geometric meaning (GCE insight)]
At the instant of closest approach, $\vec{r}_{P/Q}(t_{\min}) \cdot \vec{v}_{P/Q} = 0$: the relative position is \textbf{perpendicular} to the relative velocity. Geometrically, in the target's frame the interceptor travels along a straight line; the closest distance of approach is the \emph{perpendicular distance from the target to that line}.
\end{propriete}

\begin{exoresolu}[Closest distance of approach]
At $t = 0$ (hours), aircraft $P$ is at $(30\vec{i} + 10\vec{j})\ \mathrm{km}$ flying with velocity $(-200\vec{i} + 100\vec{j})\ \mathrm{km\,h^{-1}}$, and aircraft $Q$ is at the origin flying with velocity $(100\vec{i} + 150\vec{j})\ \mathrm{km\,h^{-1}}$. Find the time of closest approach and the minimum separation.
\tcblower
\textbf{Solution.} Relative motion of $P$ with respect to $Q$:
\[ \vec{r}_{P/Q}(0) = 30\vec{i} + 10\vec{j}, \qquad \vec{v}_{P/Q} = -300\vec{i} - 50\vec{j}. \]
So $\vec{r}_{P/Q}(t) = (30 - 300t)\vec{i} + (10 - 50t)\vec{j}$, and
\[ D(t) = (30 - 300t)^2 + (10 - 50t)^2 = 92500t^2 - 19000t + 1000. \]
Completing the square (or $D'(t) = 185000t - 19000 = 0$):
\[ t_{\min} = \frac{19000}{2 \times 92500} = \frac{19}{185} \approx 0.1027\ \mathrm{h} \approx 6.16\ \mathrm{min} > 0.\ \checkmark \]
Minimum separation:
\[ D(t_{\min}) = 1000 - \frac{19000^2}{4 \times 92500} = 1000 - \frac{361000000}{370000} = 1000 - \frac{36100}{37} = \frac{900}{37} \approx 24.324, \]
\[ d_{\min} = \sqrt{\frac{900}{37}} = \frac{30}{\sqrt{37}} \approx 4.93\ \mathrm{km}. \]
\textbf{Answer:} closest approach after $\dfrac{19}{185}\ \mathrm{h}$ (about $6.2$ minutes), with minimum separation $\dfrac{30}{\sqrt{37}}\ \mathrm{km} \approx 4.9\ \mathrm{km}$ — no collision.
\end{exoresolu}

\begin{popfigure}
\begin{center}
\begin{tikzpicture}
\begin{axis}[width=0.72\linewidth, height=5.2cm, xlabel={$t$ (h)}, ylabel={$d(t)$ (km)}, xmin=0, xmax=0.25, ymin=0, ymax=34, grid=both, grid style={popline, dashed}, axis lines=left, tick label style={font=\small}, label style={popdark}]
\addplot[very thick, popblue, domain=0:0.25, samples=100] {sqrt(max(0,(30-300*x)^2 + (10-50*x)^2))};
\addplot[only marks, mark=*, mark size=2pt, poporange] coordinates {(0.1027,4.93)};
\draw[dashed, poporange] (axis cs:0.1027,0) -- (axis cs:0.1027,4.93);
\node[popdark, anchor=west, font=\small] at (axis cs:0.115,6.5) {closest approach};
\node[popdark, anchor=west, font=\small] at (axis cs:0.115,3.5) {$(0.103,\ 4.93)$};
\end{axis}
\end{tikzpicture}
\legende{Separation $d(t) = |\vec{r}_P(t) - \vec{r}_Q(t)|$ against time: the curve bottoms out at the closest distance of approach.}
\end{center}
\end{popfigure}

\begin{aretenir}[Key points]
\begin{itemize}
\item Interception/collision: $\vec{r}_P(t) = \vec{r}_Q(t)$ for some $t > 0$ — two simultaneous component equations.
\item Interception is possible iff $\vec{v}_{P/Q}$ is directed along the initial line joining the bodies, from interceptor towards target.
\item Aim at the \emph{future} meeting point, never at the target's present position.
\item Always check $t > 0$.
\item No collision $\Rightarrow$ minimise $D(t) = |\vec{r}_{P/Q}(t)|^2$ by completing the square or differentiating; at closest approach $\vec{r}_{P/Q} \perp \vec{v}_{P/Q}$.
\end{itemize}
\end{aretenir}

\begin{exercice}[Interception in the estuary]
At $t = 0$, a customs launch is at the origin near the Wouri estuary and can travel at $30\ \mathrm{km\,h^{-1}}$. A boat is at $(12\vec{i} + 5\vec{j})\ \mathrm{km}$ moving with velocity $(8\vec{i} + 6\vec{j})\ \mathrm{km\,h^{-1}}$. Find the velocity the launch must steer to intercept the boat, and the time and position of interception.
\end{exercice}

\begin{exercice}[Closest approach of two ships]
Ship $A$ is at $(40\vec{i})\ \mathrm{km}$ with velocity $(-10\vec{i} + 5\vec{j})\ \mathrm{km\,h^{-1}}$; ship $B$ is at the origin with velocity $(6\vec{i} + 8\vec{j})\ \mathrm{km\,h^{-1}}$. Show that the ships do not collide, and find the time of closest approach and the minimum distance between them.
\end{exercice}

\begin{corrige}[Exercise 1]
Let $\vec{v}_L = (a\vec{i} + b\vec{j})$ with $a^2+b^2 = 900$. Interception: $at = 12 + 8t$, $bt = 5 + 6t$. Then $\dfrac{a-8}{b-6} = \dfrac{12}{5}$, so $5a - 12b = -32$, $a = \dfrac{12b-32}{5}$. Substituting: $(12b-32)^2 + 25b^2 = 22500$, i.e.\ $169b^2 - 768b - 21476 = 0$. The discriminant is $\Delta = 768^2 + 4 \times 169 \times 21476 = 15\,107\,600$, $\sqrt{\Delta} \approx 3886.85$. Taking the positive root: $b = \dfrac{768 + 3886.85}{338} \approx 13.77$, $a = \dfrac{12 \times 13.77 - 32}{5} \approx 26.65$. Check: $26.65^2 + 13.77^2 \approx 900$. \checkmark\ Then $t = \dfrac{12}{a-8} \approx \dfrac{12}{18.65} \approx 0.643\ \mathrm{h} \approx 38.6\ \mathrm{min} > 0$. \checkmark\ Interception point: $\vec{r} \approx (17.1\vec{i} + 8.9\vec{j})\ \mathrm{km}$.
\end{corrige}

\begin{corrige}[Exercise 2]
$\vec{r}_{A/B}(0) = 40\vec{i}$, $\vec{v}_{A/B} = -16\vec{i} - 3\vec{j}$. This is not a scalar multiple of $40\vec{i}$ (the $\vec{j}$-component is non-zero), so no collision. $D(t) = (40-16t)^2 + (3t)^2 = 265t^2 - 1280t + 1600$. $D'(t) = 530t - 1280 = 0 \Rightarrow t_{\min} = \dfrac{128}{53} \approx 2.415\ \mathrm{h} > 0$. \checkmark\ $D(t_{\min}) = 1600 - \dfrac{1280^2}{4 \times 265} = 1600 - \dfrac{1\,638\,400}{1060} = \dfrac{57600}{1060} = \dfrac{2880}{53} \approx 54.34$, so $d_{\min} = \sqrt{2880/53} \approx 7.37\ \mathrm{km}$.
\end{corrige}

\begin{savaistu}[Collision avoidance at sea and in the air]
Modern ships use radar systems (ARPA) that continuously compute exactly the quantities of this section — relative velocity, time and distance of closest approach (called CPA, ``closest point of approach'') — and sound an alarm when the CPA falls below a safety margin. Aircraft collision-avoidance systems (TCAS) do the same in three dimensions. The mathematics protecting travellers over the Atlantic and in the Douala and Yaoundé flight corridors is precisely the relative-motion theory you have just mastered.
\end{savaistu}

\coursec{Applications of Interception and Collision}

In the previous section we built the mathematical machinery of interception: two bodies $P$ and $Q$ with position vectors $\vec{r}_P(t)$ and $\vec{r}_Q(t)$ meet at time $t$ if and only if $\vec{r}_P(t)=\vec{r}_Q(t)$, or equivalently if the \cle{relative position} $\vec{r}_{P/Q}(t)=\vec{r}_P(t)-\vec{r}_Q(t)$ vanishes. We now put this machinery to work on the families of problems that the GCE Advanced Level examiner loves most: maritime interception off our own coasts, collision-avoidance at sea, interceptions against the clock, delayed pursuits, and (as an extension) pursuits in which one body accelerates briefly before settling into uniform motion. Throughout, the golden rule is unchanged: \emph{reduce the two-body problem to a one-body problem by working with relative position and relative velocity}.

\courssub{Maritime Interception: Choosing Speed or Course}

\begin{experience}[A patrol off the Kribi coast]
A patrol boat of the Cameroonian navy, stationed off the deep-sea port of Kribi, spots a suspect vessel steaming on a fixed course at constant speed. The patrol boat is faster. Two questions arise immediately: \emph{which course} should the patrol boat steer to intercept, and \emph{what is the earliest possible time} of interception? These are the two standard GCE questions, and both are answered by the same vector equation.
\end{experience}

\begin{methode}[Setting up an interception problem]
\begin{enumerate}
\item Fix one Cartesian frame $(O;\vec{i},\vec{j})$ for the whole problem, with $\vec{i}$ pointing East and $\vec{j}$ pointing North. Choose a single time origin $t=0$.
\item Write the position vectors $\vec{r}_P(t)=\vec{r}_P(0)+t\vec{v}_P$ and $\vec{r}_Q(t)=\vec{r}_Q(0)+t\vec{v}_Q$ of the interceptor $P$ and the target $Q$.
\item Impose the interception condition $\vec{r}_P(t)=\vec{r}_Q(t)$, i.e.
\[ t\,\vec{v}_P = \vec{r}_Q(0)-\vec{r}_P(0)+t\,\vec{v}_Q . \]
\item If the course of $P$ is required, solve the two component equations for $t$ and the components of $\vec{v}_P$ (with $|\vec{v}_P|$ given). If the \emph{earliest} interception is required, use the method below.
\end{enumerate}
\end{methode}

\begin{methode}[Earliest interception]
The relative position at time $t$ is $\vec{r}_{Q/P}(t)=\vec{r}_{Q/P}(0)+t(\vec{v}_Q-\vec{v}_P)$. Interception occurs when this vanishes, i.e. when $\vec{v}_P-\vec{v}_Q$ points \emph{along the line of sight} from $P$ to $Q$. To make $t$ as small as possible:
\begin{enumerate}
\item Write the closing velocity $\vec{v}_P-\vec{v}_Q$ and require it to be parallel to $\vec{r}_{Q/P}(0)$ (the line of sight).
\item The interception time is $t=\dfrac{|\vec{r}_{Q/P}(0)|}{\text{component of }(\vec{v}_P-\vec{v}_Q)\text{ along the line of sight}}$.
\item Hence $t$ is minimised by \textbf{maximising the component of the interceptor's velocity along the line of sight}: steer $P$ so that as much of its speed as possible closes the gap directly.
\end{enumerate}
\end{methode}

\begin{exoresolu}[The Kribi patrol boat]
At noon, a patrol boat $P$ is at the point $A$. A vessel $Q$ is $20\ \mathrm{km}$ due East of $P$ and is steaming due North at a steady $15\ \mathrm{km\,h^{-1}}$. The patrol boat can make $25\ \mathrm{km\,h^{-1}}$. Find the course the patrol boat must steer to intercept the vessel, and the time of interception.
\tcblower
\textbf{Set-up.} Take $A$ as origin, $\vec{i}$ East, $\vec{j}$ North, $t=0$ at noon, distances in km, times in hours.
\[ \vec{r}_P(0)=\vec{0},\qquad \vec{r}_Q(0)=20\vec{i},\qquad \vec{v}_Q=15\vec{j},\qquad \vec{v}_P=25(\cos\theta\,\vec{i}+\sin\theta\,\vec{j}), \]
where $\theta$ is measured from East. \hfill\textbf{[M1: vector set-up]}

\textbf{Interception condition.} $\vec{r}_P(t)=\vec{r}_Q(t)$ gives
\[ 25t\cos\theta = 20, \qquad 25t\sin\theta = 15t. \]
\textbf{[M1]}
From the second equation, $\sin\theta=\dfrac{15}{25}=\dfrac{3}{5}$, so $\theta\approx 36.9^{\circ}$ and $\cos\theta=\dfrac{4}{5}$. \hfill\textbf{[A1]}
Then $25t\cdot\tfrac45=20$ gives $t=1\ \mathrm{h}$. \hfill\textbf{[A1]}

\textbf{Answer.} The patrol boat steers on a bearing of $90^{\circ}-36.9^{\circ}=053.1^{\circ}$ (i.e. N $36.9^{\circ}$ E) and intercepts the vessel at $1$ p.m., $25$ km from $A$. \hfill\textbf{[B1]}
\end{exoresolu}

\courssub{Collision or Not? Closest Distance of Approach}

Very often the two vessels are \emph{not} trying to meet — quite the opposite. Given their courses and speeds, the captain wants to know: \emph{will we collide, and if not, how close will we pass, and when?}

\begin{propriete}[Collision test]
Let $P$ and $Q$ move uniformly with $\vec{r}_{P/Q}(t)=\vec{r}_{P/Q}(0)+t\,\vec{v}_{P/Q}$, where $\vec{v}_{P/Q}=\vec{v}_P-\vec{v}_Q$ is the relative velocity. Then
\[ \text{collision} \iff \vec{r}_{P/Q}(0)\ \text{is parallel to}\ \vec{v}_{P/Q}\ \text{with opposite sense}, \]
i.e. $\vec{r}_{P/Q}(0)=\lambda\,\vec{v}_{P/Q}$ for some $\lambda<0$; the collision time is $t=-\lambda>0$. Geometrically: in $Q$'s frame, $P$ moves in a straight line, and collision occurs exactly when that line passes through $Q$.
\end{propriete}

\begin{propriete}[Closest approach]
If the collision test fails, the separation $d(t)=|\vec{r}_{P/Q}(0)+t\vec{v}_{P/Q}|$ is minimised when
\[ t_{\min}=-\frac{\vec{r}_{P/Q}(0)\cdot\vec{v}_{P/Q}}{|\vec{v}_{P/Q}|^{2}}, \qquad
d_{\min}=\frac{\bigl|\vec{r}_{P/Q}(0)\times\vec{v}_{P/Q}\bigr|}{|\vec{v}_{P/Q}|}, \]
provided $t_{\min}>0$ (otherwise the ships are already separating). The proof (differentiate $d^{2}(t)$, or complete the square) is instructive but is \emph{not} required at GCE; quoting and applying the result, or minimising $d^2(t)$ directly, earns full marks.
\end{propriete}

\begin{attention}[Sign of $t_{\min}$]
A negative $t_{\min}$ means the closest approach happened \emph{before} $t=0$: the ships are moving apart and there is nothing to fear. Always check the sign before reporting a "closest approach in the future".
\end{attention}

\begin{exoresolu}[Two ships in the Gulf of Guinea]
At 08:00, ship $P$ is at the point with position vector $(2\vec{i}+18\vec{j})$ km and steams with velocity $(6\vec{i}-8\vec{j})\ \mathrm{km\,h^{-1}}$. Ship $Q$ is at $(-10\vec{i}+6\vec{j})$ km steaming with velocity $(4\vec{i}+2\vec{j})\ \mathrm{km\,h^{-1}}$. Determine whether the ships collide; if not, find their closest distance of approach and the time at which it occurs.
\tcblower
\textbf{Relative motion.} \hfill\textbf{[M1]}
\[ \vec{r}_{P/Q}(0)=(2-(-10))\vec{i}+(18-6)\vec{j}=12\vec{i}+12\vec{j},\qquad
\vec{v}_{P/Q}=(6-4)\vec{i}+(-8-2)\vec{j}=2\vec{i}-10\vec{j}. \]

\textbf{Collision test.} Is $12\vec{i}+12\vec{j}$ a negative multiple of $2\vec{i}-10\vec{j}$? We would need $12=2\lambda$ and $12=-10\lambda$, i.e. $\lambda=6$ and $\lambda=-1.2$ simultaneously — impossible. \textbf{No collision.} \hfill\textbf{[M1 A1]}

\textbf{Closest approach.}
\[ t_{\min}=-\frac{(12)(2)+(12)(-10)}{2^{2}+(-10)^{2}}=-\frac{24-120}{104}=\frac{96}{104}=\frac{12}{13}\ \mathrm{h}\approx 0.923\ \mathrm{h}. \]
\textbf{[M1 A1]}
Relative position at closest approach:
\[ \vec{r}_{P/Q}\!\left(\tfrac{12}{13}\right)=\left(12+\tfrac{24}{13}\right)\vec{i}+\left(12-\tfrac{120}{13}\right)\vec{j}=\tfrac{180}{13}\vec{i}+\tfrac{36}{13}\vec{j}, \]
\[ d_{\min}=\frac{1}{13}\sqrt{180^{2}+36^{2}}=\frac{36}{13}\sqrt{26}\approx 14.1\ \mathrm{km}. \]
\textbf{[M1 A1]}
\textbf{Answer.} No collision; the ships pass within about $14.1$ km of each other at approximately 08:55. \hfill\textbf{[B1]}
\end{exoresolu}

\begin{popfigure}
\begin{tikzpicture}[scale=0.42, >=stealth]
\draw[->, gray] (-2,0) -- (22,0) node[below left] {\small East};
\draw[->, gray] (0,-2) -- (0,20) node[below left] {\small North};
% P course
\fill[popblue] (2,18) circle (0.18) node[above left] {\small $P$ at 08:00};
\draw[->, very thick, popblue] (2,18) -- (8,10) node[midway, right] {\small $\vec{v}_P$};
\draw[dashed, popblue] (2,18) -- (12.6,3.8);
% Q course
\fill[poporange] (-10,6) circle (0.18) node[left] {\small $Q$ at 08:00};
\draw[->, very thick, poporange] (-10,6) -- (-6,8) node[midway, below] {\small $\vec{v}_Q$};
\draw[dashed, poporange] (-10,6) -- (2,12);
% closest approach points
\fill[popdark] (12.6,3.8) circle (0.15);
\fill[popdark] (2,12) circle (0.15);
\draw[<->, very thick, poppink] (12.6,3.8) -- (2,12) node[midway, below right] {\small $d_{\min}\approx14.1$ km};
\node[popdark] at (13.5,3.2) {\small $P$ at 08:55};
\node[popdark] at (2.6,12.8) {\small $Q$ at 08:55};
\end{tikzpicture}
\legende{The two courses (dashed) do cross on the map, but the ships reach the crossing region at different times: the real test is the relative-velocity test, not the intersection of the drawn tracks.}
\end{popfigure}

\begin{attention}[Crossing tracks is not a collision]
The most frequent error in this topic: seeing that the two \emph{drawn} course lines intersect and concluding that the ships collide. They collide only if they are at the same point \emph{at the same time} — which is exactly the relative-velocity test above.
\end{attention}

\courssub{Interception Against the Clock}

In a third family of problems the interception \emph{time} is imposed ("intercept in exactly 30 minutes") and we must find the required speed or bearing of the interceptor.

\begin{methode}[Given interception time]
\begin{enumerate}
\item Write the interception condition $\vec{r}_P(T)=\vec{r}_Q(T)$ with the \emph{given} $T$.
\item Solve for the unknown vector $\vec{v}_P=\vec{v}_Q+\dfrac{\vec{r}_Q(0)-\vec{r}_P(0)}{T}$.
\item Read off the required speed $|\vec{v}_P|$ and bearing from its components.
\end{enumerate}
\end{methode}

\begin{exoresolu}[Intercept in 40 minutes]
A coastguard launch at $O$ must intercept, in exactly $40$ minutes, a trawler which is currently $12\ \mathrm{km}$ away on a bearing of $030^{\circ}$ and moving due East at $9\ \mathrm{km\,h^{-1}}$. Find the velocity (speed and bearing) the launch must maintain.
\tcblower
Take $\vec{i}$ East, $\vec{j}$ North, $T=\tfrac{2}{3}\ \mathrm{h}$. The trawler's initial position is
\[ \vec{r}_Q(0)=12(\sin 30^{\circ}\,\vec{i}+\cos 30^{\circ}\,\vec{j})=6\vec{i}+6\sqrt{3}\,\vec{j}\ \mathrm{km}. \]
\textbf{[M1]}
Interception condition with $\vec{v}_Q=9\vec{i}$:
\[ \vec{v}_P=\vec{v}_Q+\frac{\vec{r}_Q(0)}{T}=9\vec{i}+\frac{3}{2}\left(6\vec{i}+6\sqrt{3}\,\vec{j}\right)=18\vec{i}+9\sqrt{3}\,\vec{j}\ \mathrm{km\,h^{-1}}. \]
\textbf{[M1 A1]}
Speed: $|\vec{v}_P|=\sqrt{18^{2}+(9\sqrt3)^{2}}=\sqrt{324+243}=\sqrt{567}=9\sqrt7\approx 23.8\ \mathrm{km\,h^{-1}}$. \hfill\textbf{[A1]}
Direction: $\tan\alpha=\dfrac{9\sqrt3}{18}=\dfrac{\sqrt3}{2}$, so $\alpha\approx 40.9^{\circ}$ North of East, i.e. a bearing of about $049.1^{\circ}$. \hfill\textbf{[A1]}
\end{exoresolu}

\courssub{Delayed Pursuit: One Body Starts Later}

\begin{attention}[One clock for everybody]
The commonest source of disaster in pursuit problems is mixing the two bodies' start times. \textbf{Always write every position vector using the same time origin.} If $P$ starts at $t=0$ and $Q$ starts $T_0$ later, then
\[ \vec{r}_Q(t)=\vec{r}_Q(T_0)+(t-T_0)\vec{v}_Q \quad (t\ge T_0), \]
and the interception condition $\vec{r}_P(t)=\vec{r}_Q(t)$ is solved for $t>T_0$. Never write $t\vec{v}_Q$ for a body that has only been moving for $t-T_0$ hours.
\end{attention}

\begin{popfigure}
\begin{tikzpicture}[scale=1.0, >=stealth]
\draw[->, thick] (0,0) -- (11,0) node[below left] {\small $t$ (h)};
\foreach \x/\lab in {0/0, 2/1, 4/2, 6/3, 8/4, 10/5} \draw (\x,0.08) -- (\x,-0.08) node[below] {\small \lab};
% P timeline
\draw[line width=2.2pt, popblue] (0,1.0) -- (10,1.0);
\fill[popblue] (0,1.0) circle (0.09); \node[left, popblue] at (0,1.0) {\small $P$ moving};
% Q timeline
\draw[line width=2.2pt, poporange] (4,2.0) -- (10,2.0);
\draw[dashed, poporange] (0,2.0) -- (4,2.0);
\fill[poporange] (4,2.0) circle (0.09); \node[left, poporange] at (0,2.0) {\small $Q$ moving};
\draw[<->, popdark, thick] (0,2.55) -- (4,2.55) node[midway, above] {\small delay $T_0=2$ h};
% interception
\draw[dashed, poppink, thick] (8,-0.3) -- (8,2.7);
\fill[poppink] (8,1.0) circle (0.11); \fill[poppink] (8,2.0) circle (0.11);
\node[poppink] at (8,-0.6) {\small interception $t=4$};
\end{tikzpicture}
\legende{Timeline for a delayed pursuit: $P$ moves from $t=0$; $Q$ only from $t=T_0=2$ h. At the interception instant $t=4$ h, $P$ has travelled for $4$ h but $Q$ only for $2$ h — the factor $(t-T_0)$ is essential.}
\end{popfigure}

\begin{exoresolu}[The delayed speedboat]
A smugglers' boat leaves a creek near Limbe at 06:00, heading due West at $12\ \mathrm{km\,h^{-1}}$. At 08:00 a patrol speedboat leaves the same creek and pursues on the same straight line at $30\ \mathrm{km\,h^{-1}}$. At what time does the patrol catch the boat, and how far from the creek?
\tcblower
Take $t=0$ at 06:00, the creek as origin, $\vec{i}$ pointing West.
\[ \vec{r}_S(t)=12t\,\vec{i},\qquad \vec{r}_P(t)=30(t-2)\,\vec{i}\ \ (t\ge 2). \]
\textbf{[M1: same time origin, offset $(t-2)$]}
Interception: $12t=30(t-2)$, so $18t=60$, $t=\dfrac{10}{3}\ \mathrm{h}=3\ \mathrm{h}\ 20\ \mathrm{min}$. \hfill\textbf{[M1 A1]}
Distance: $12\times\tfrac{10}{3}=40\ \mathrm{km}$. \hfill\textbf{[A1]}
\textbf{Answer.} The patrol catches the boat at 09:20, $40$ km from the creek. \hfill\textbf{[B1]}
\end{exoresolu}

\courssub{Extension: Interception with Piecewise Motion}

\begin{savaistu}[GCE extension]
Occasionally a harder question lets one body \emph{accelerate briefly} and then continue uniformly. The technique is to split the motion into phases: compute the position and velocity at the end of the acceleration phase, then restart a uniform-motion problem from that instant — exactly as in the delayed-pursuit method, with the clock reset at the phase change. This is beyond the routine GCE question but well within reach of a strong candidate.
\end{savaistu}

\begin{exoresolu}[Accelerate, then cruise]
A drone $Q$ is $600\ \mathrm{m}$ due North of a launcher and flies due East at a steady $20\ \mathrm{m\,s^{-1}}$. An interceptor $P$ starts from rest at the launcher, accelerates uniformly at $4\ \mathrm{m\,s^{-2}}$ due North for $5\ \mathrm{s}$, then continues at constant velocity. Show that $P$ cannot intercept $Q$, and find the closest distance of approach.
\tcblower
\textbf{Phase 1} ($0\le t\le 5$): $P$ reaches speed $v=4\times 5=20\ \mathrm{m\,s^{-1}}$ North and covers $\tfrac12(4)(5^{2})=50\ \mathrm{m}$. \hfill\textbf{[M1 A1]}

\textbf{Phase 2} ($t\ge 5$): reset the clock with $\tau=t-5$. Then
\[ \vec{r}_P=(50+20\tau)\vec{j},\qquad \vec{r}_Q=20(5+\tau)\vec{i}+600\vec{j}=(100+20\tau)\vec{i}+600\vec{j}. \]
\textbf{[M1]}
Relative position: $\vec{r}_{P/Q}=-(100+20\tau)\vec{i}+(20\tau-550)\vec{j}$.
The East component $-(100+20\tau)$ never vanishes for $\tau\ge 0$, so \textbf{interception is impossible}. \hfill\textbf{[A1]}

Minimise $d^{2}(\tau)=(100+20\tau)^{2}+(550-20\tau)^{2}$:
\[ \frac{d(d^{2})}{d\tau}=40(100+20\tau)-40(550-20\tau)=0 \implies 100+20\tau=550-20\tau \implies \tau=11.25\ \mathrm{s}. \]
\textbf{[M1 A1]}
Then $d_{\min}=\sqrt{(325)^{2}+(325)^{2}}=325\sqrt2\approx 460\ \mathrm{m}$, occurring at $t=5+11.25=16.25\ \mathrm{s}$. \hfill\textbf{[A1]}
\end{exoresolu}

\courssub{Graded GCE-Style Practice}

\begin{exercice}[Exercise 1 — course for interception]
A ferry leaves port $A$ and steams due North at $14\ \mathrm{km\,h^{-1}}$. A water taxi is $10\ \mathrm{km}$ due West of $A$ and can make $20\ \mathrm{km\,h^{-1}}$. Find the bearing the taxi must steer to intercept the ferry, and the time taken.
\end{exercice}

\begin{exercice}[Exercise 2 — collision test]
At 10:00, ship $A$ is at $(4\vec{i}+3\vec{j})$ km moving with velocity $(5\vec{i}+4\vec{j})\ \mathrm{km\,h^{-1}}$; ship $B$ is at $(-6\vec{i}+11\vec{j})$ km moving with velocity $(3\vec{i}+2\vec{j})\ \mathrm{km\,h^{-1}}$. Determine whether the ships collide; if they do, find the time and position of the collision.
\end{exercice}

\begin{exercice}[Exercise 3 — closest approach]
Ship $P$ is $30\ \mathrm{km}$ due North of ship $Q$. $P$ steams on a bearing of $120^{\circ}$ at $16\ \mathrm{km\,h^{-1}}$; $Q$ steams due North at $10\ \mathrm{km\,h^{-1}}$. Find the closest distance of approach and the time at which it occurs.
\end{exercice}

\begin{exercice}[Exercise 4 — delayed pursuit]
A cyclist leaves Buea at 07:00 riding towards Mutengene at $15\ \mathrm{km\,h^{-1}}$. At 07:30 a motorist leaves Buea on the same road at $60\ \mathrm{km\,h^{-1}}$. Find when and where the motorist overtakes the cyclist.
\end{exercice}

\begin{corrige}[Exercise 1]
With $A$ as origin, $\vec{r}_T(0)=-10\vec{i}$, $\vec{v}_F=14\vec{j}$, $\vec{v}_T=20(\cos\theta\,\vec{i}+\sin\theta\,\vec{j})$. Interception: $20t\cos\theta=10$, $20t\sin\theta=14t$, so $\sin\theta=\tfrac{7}{10}$, $\theta\approx44.4^{\circ}$ North of East, bearing $\approx045.6^{\circ}$; $\cos\theta=\tfrac{\sqrt{51}}{10}$ and $t=\dfrac{10}{2\sqrt{51}}\approx0.70\ \mathrm{h}\approx42$ min.
\end{corrige}

\begin{corrige}[Exercise 2]
$\vec{r}_{A/B}(0)=10\vec{i}-8\vec{j}$, $\vec{v}_{A/B}=2\vec{i}+2\vec{j}$. Is $10\vec{i}-8\vec{j}=\lambda(2\vec{i}+2\vec{j})$? $\lambda=5$ from the $\vec{i}$-component but $\lambda=-4$ from the $\vec{j}$-component: \emph{not} parallel, so the ships do \textbf{not} collide. (Closest approach: $t_{\min}=-\dfrac{20-16}{8}=-0.5$ h $<0$, so the ships are already separating; their closest approach was at 09:30.) \emph{Lesson: always carry out the test rather than assuming the conclusion.}
\end{corrige}

\begin{corrige}[Exercise 3]
$\vec{r}_{P/Q}(0)=30\vec{j}$, $\vec{v}_P=16(\sin120^{\circ}\vec{i}+\cos120^{\circ}\vec{j})=8\sqrt3\,\vec{i}-8\vec{j}$, $\vec{v}_Q=10\vec{j}$, so $\vec{v}_{P/Q}=8\sqrt3\,\vec{i}-18\vec{j}$, $|\vec{v}_{P/Q}|^{2}=192+324=516$.
$t_{\min}=-\dfrac{(30)(-18)}{516}=\dfrac{540}{516}\approx1.047\ \mathrm{h}\approx63$ min.
$d_{\min}=\dfrac{|30\times 8\sqrt3|}{\sqrt{516}}=\dfrac{240\sqrt3}{\sqrt{516}}\approx18.3\ \mathrm{km}$.
\end{corrige}

\begin{corrige}[Exercise 4]
Same time origin at 07:00: $r_C(t)=15t$, $r_M(t)=60(t-0.5)$. Overtaking: $15t=60t-30$, so $t=\tfrac{2}{3}\ \mathrm{h}=40$ min; distance $15\times\tfrac23=10$ km. The motorist overtakes the cyclist at 07:40, $10$ km from Buea.
\end{corrige}

\begin{aretenir}[Key points of the section]
\begin{itemize}
\item Every interception or collision problem reduces to $\vec{r}_{P/Q}(t)=\vec{r}_{P/Q}(0)+t\,\vec{v}_{P/Q}=\vec{0}$.
\item \textbf{Collision test:} collision iff $\vec{r}_{P/Q}(0)$ is parallel to $\vec{v}_{P/Q}$ in the opposite sense. Crossing tracks on a map is \emph{not} a collision.
\item \textbf{Closest approach:} $t_{\min}=-\dfrac{\vec{r}_{P/Q}(0)\cdot\vec{v}_{P/Q}}{|\vec{v}_{P/Q}|^{2}}$; check that $t_{\min}>0$.
\item \textbf{Earliest interception:} maximise the component of the interceptor's velocity along the line of sight.
\item \textbf{Given time $T$:} $\vec{v}_P=\vec{v}_Q+\dfrac{\vec{r}_Q(0)-\vec{r}_P(0)}{T}$.
\item \textbf{Delayed starts:} one clock for all bodies; use $(t-T_0)$ for the late starter.
\item \textbf{Piecewise motion:} compute position and velocity at each phase change and restart a uniform problem.
\end{itemize}
\end{aretenir}

\begin{attention}[Exam technique]
State each vector equation \emph{before} substituting numbers: "$\vec{r}_P(t)=\vec{r}_Q(t)$ gives \dots". GCE marking schemes award method marks for the correct vector set-up even when arithmetic slips later. Conversely, a correct numerical answer with no visible vector equation earns very little.
\end{attention}

\coursec{Relative Position and Relative Motion: Concept, Treatment and Applications}

In Section 5 we answered questions of the type ``will the two boats collide, and if not, how close do they get?''. In every case the decisive step was to stop watching the two bodies separately and to watch \emph{one body from the other}. This final section makes that idea precise and general: we define the \cle{relative position} vector, differentiate it to recover \cle{relative velocity}, and show that the whole of this chapter --- relative velocity, true velocity, interception, closest approach, safety ranges --- is contained in a single formula, $d(t)=\left|\vec{r}_0+\vec{v}_{\mathrm{rel}}\,t\right|$. We finish with a synthesis of the chapter and timed GCE-style practice.

\courssub{Relative Position: ``Where is $A$ as seen from $B$?''}

Imagine you are sitting in a taxi $B$ at the Mokolo junction in Yaound\'e, and a friend in another taxi $A$ phones you: ``Where are you?'' You do not answer with map coordinates; you answer ``I am 200 metres east and 150 metres north of you.'' That answer is a \emph{relative position}: the position of $A$ measured from $B$, not from a fixed origin.

\begin{definition}[Relative position and relative velocity]
Let $A$ and $B$ be two moving points whose position vectors relative to a fixed origin $O$ are $\vec{r}_A(t)$ and $\vec{r}_B(t)$. The \cle{relative position} of $A$ with respect to $B$ is
\[
\vec{r}_{A/B}(t)=\vec{r}_A(t)-\vec{r}_B(t).
\]
It is the position of $A$ \emph{as seen by an observer moving with $B$}. The \cle{relative velocity} of $A$ with respect to $B$ is its time derivative:
\[
\vec{v}_{A/B}(t)=\frac{d}{dt}\vec{r}_{A/B}(t)=\vec{v}_A(t)-\vec{v}_B(t),
\]
exactly the relative velocity of Lesson 102.
\end{definition}

Note the antisymmetry: $\vec{r}_{B/A}=-\vec{r}_{A/B}$. If $A$ is 200 m east of $B$, then $B$ is 200 m west of $A$. Also note that $\left|\vec{r}_{A/B}\right|=\left|\vec{r}_{B/A}\right|$: the \emph{distance} between the two bodies does not depend on who is observing.

\begin{popfigure}
\begin{center}
\begin{tikzpicture}[scale=0.9, >=stealth]
% axes
\draw[->, gray] (0,0) -- (8.6,0) node[below left] {$x$};
\draw[->, gray] (0,0) -- (0,4.6) node[below left] {$y$};
\node[below left] at (0,0) {$O$};
% time t1
\coordinate (A1) at (5.2,3.6);
\coordinate (B1) at (2.0,1.4);
\draw[->, thick, popblue] (0,0) -- (A1) node[midway, above left, popblue] {$\vec{r}_A(t_1)$};
\draw[->, thick, poporange] (0,0) -- (B1) node[midway, below left, poporange] {$\vec{r}_B(t_1)$};
\draw[->, very thick, poppurple] (B1) -- (A1) node[midway, above right, poppurple] {$\vec{r}_{A/B}(t_1)$};
\fill[popblue] (A1) circle (2pt) node[above] {$A(t_1)$};
\fill[poporange] (B1) circle (2pt) node[below right] {$B(t_1)$};
% time t2 (dashed, lighter)
\coordinate (A2) at (7.6,2.6);
\coordinate (B2) at (4.4,0.9);
\draw[->, thick, popblue, dashed] (0,0) -- (A2);
\draw[->, thick, poporange, dashed] (0,0) -- (B2);
\draw[->, very thick, poppurple, dashed] (B2) -- (A2) node[midway, below right, poppurple] {$\vec{r}_{A/B}(t_2)$};
\fill[popblue] (A2) circle (2pt) node[above] {$A(t_2)$};
\fill[poporange] (B2) circle (2pt) node[below] {$B(t_2)$};
\end{tikzpicture}
\legende{The vector triangle $\vec{r}_A=\vec{r}_B+\vec{r}_{A/B}$ at two different times. As $A$ and $B$ move, the relative position vector $\vec{r}_{A/B}$ changes in magnitude and direction.}
\end{center}
\end{popfigure}

The figure shows the fundamental triangle of this chapter: for any common origin $O$,
\[
\vec{r}_A(t)=\vec{r}_B(t)+\vec{r}_{A/B}(t),
\]
so $\vec{r}_{A/B}$ is simply the third side of the triangle, the arrow that goes \emph{from $B$ to $A$}.

\courssub{Uniform Relative Motion and the Separation $d(t)$}

Suppose now that both bodies move with \emph{constant} velocities $\vec{v}_A$ and $\vec{v}_B$ (the standard GCE hypothesis). Then
\[
\vec{r}_A(t)=\vec{r}_A(0)+\vec{v}_A t, \qquad \vec{r}_B(t)=\vec{r}_B(0)+\vec{v}_B t.
\]

\begin{propriete}[Uniform relative motion --- proof examinable and instructive]
If $\vec{v}_A$ and $\vec{v}_B$ are constant, then the relative motion is uniform:
\[
\vec{r}_{A/B}(t)=\vec{r}_0+\vec{v}_{\mathrm{rel}}\,t,
\qquad\text{where }\ \vec{r}_0=\vec{r}_A(0)-\vec{r}_B(0),\ \ \vec{v}_{\mathrm{rel}}=\vec{v}_A-\vec{v}_B,
\]
and the \cle{distance of separation} is
\[
d(t)=\left|\vec{r}_{A/B}(t)\right|=\left|\vec{r}_0+\vec{v}_{\mathrm{rel}}\,t\right|.
\]
In particular, seen from $B$, the body $A$ moves along a \emph{straight line} with constant velocity $\vec{v}_{\mathrm{rel}}$.
\end{propriete}

\textbf{Proof.} Subtracting the two expressions for $\vec{r}_A(t)$ and $\vec{r}_B(t)$:
\[
\vec{r}_{A/B}(t)=\vec{r}_A(t)-\vec{r}_B(t)=\bigl(\vec{r}_A(0)-\vec{r}_B(0)\bigr)+\bigl(\vec{v}_A-\vec{v}_B\bigr)t=\vec{r}_0+\vec{v}_{\mathrm{rel}}\,t.
\]
This is a linear function of $t$, i.e.\ the equation of uniform (straight-line, constant-velocity) motion. Taking magnitudes gives $d(t)$. \hfill$\blacksquare$

This single property \emph{unifies} Lessons 107--112:
\begin{itemize}
\item \textbf{Collision (interception)} occurs iff $d(t)=0$ for some $t\geq 0$, i.e.\ iff $\vec{r}_0+\vec{v}_{\mathrm{rel}}\,t=\vec{0}$, i.e.\ iff $\vec{r}_0$ is a negative scalar multiple of $\vec{v}_{\mathrm{rel}}$: the relative velocity points exactly along the initial line of sight.
\item \textbf{Closest approach} is the minimum of $d(t)$, found (Section 4) at $t^*=-\dfrac{\vec{r}_0\cdot\vec{v}_{\mathrm{rel}}}{|\vec{v}_{\mathrm{rel}}|^2}$, giving $d_{\min}^2=|\vec{r}_0|^2-\dfrac{(\vec{r}_0\cdot\vec{v}_{\mathrm{rel}})^2}{|\vec{v}_{\mathrm{rel}}|^2}$.
\end{itemize}

\begin{attention}[Differentiate the vector, not its magnitude]
A very common error: a candidate writes $d(t)=|\vec{r}_{A/B}(t)|$ and then ``differentiates $d$'' hoping to get the relative speed. In general
\[
\frac{d}{dt}\bigl|\vec{r}_{A/B}\bigr|\ \neq\ \left|\frac{d}{dt}\vec{r}_{A/B}\right|=|\vec{v}_{\mathrm{rel}}|.
\]
The rate of change of the \emph{distance} is only the \emph{component} of $\vec{v}_{\mathrm{rel}}$ along the line joining the bodies. Always differentiate the \textbf{vector} $\vec{r}_{A/B}$ first; take magnitudes only at the end, when you need a distance.
\end{attention}

\courssub{Time Within Range: the Quadratic Inequality $d(t)\leq R$}

Many GCE questions add a practical twist: two ships must stay within radio range $R$, two aircraft must not come closer than a safety distance, a patrol boat can detect a canoe only within \SI{5}{km}. In all these problems we must find the \emph{time interval} during which $d(t)\leq R$.

\begin{methode}[Solving $d(t)\leq R$]
\begin{enumerate}
\item Write $\vec{r}_0=\vec{r}_A(0)-\vec{r}_B(0)$ and $\vec{v}_{\mathrm{rel}}=\vec{v}_A-\vec{v}_B$ in components.
\item Form $d^2(t)=|\vec{r}_0+\vec{v}_{\mathrm{rel}}\,t|^2$: this is a quadratic $at^2+bt+c$ with $a=|\vec{v}_{\mathrm{rel}}|^2$, $b=2\,\vec{r}_0\cdot\vec{v}_{\mathrm{rel}}$, $c=|\vec{r}_0|^2$.
\item Solve the inequality $d^2(t)\leq R^2$, i.e.\ $at^2+bt+(c-R^2)\leq 0$ (squaring is legal since both sides are non-negative).
\item If the discriminant is positive with roots $t_1<t_2$, then since $a>0$ the solution is $t_1\leq t\leq t_2$ (intersected with $t\geq 0$). The bodies are within range for a duration $t_2-t_1$.
\item If the discriminant is negative, $d(t)>R$ always: the bodies never come within range.
\end{enumerate}
\end{methode}

\begin{exoresolu}[Radio contact between two boats on the Wouri]
At noon, boat $A$ is at the origin and moves with velocity $(6\vec{i}+8\vec{j})\ \mathrm{km\,h^{-1}}$, while boat $B$ is at the point $(10\vec{i}+0\vec{j})$ km and moves with velocity $(-2\vec{i}+8\vec{j})\ \mathrm{km\,h^{-1}}$. Their radios have a range of $R=10$ km. Find the time interval during which the boats can communicate, and the duration of contact.
\tcblower
\textbf{Step 1: relative position.} Take $\vec{r}_{A/B}$:
\[
\vec{r}_0=\vec{r}_A(0)-\vec{r}_B(0)=-10\vec{i},\qquad
\vec{v}_{\mathrm{rel}}=\vec{v}_A-\vec{v}_B=(6-(-2))\vec{i}+(8-8)\vec{j}=8\vec{i}.
\]
Hence $\vec{r}_{A/B}(t)=(-10+8t)\vec{i}$ (with $t$ in hours after noon).

\textbf{Step 2: separation.}
\[
d(t)=|{-10}+8t|,\qquad d^2(t)=(8t-10)^2.
\]

\textbf{Step 3: inequality.} Contact requires $d^2(t)\leq 100$:
\[
(8t-10)^2\leq 100\iff |8t-10|\leq 10\iff 0\leq 8t\leq 20\iff 0\leq t\leq 2.5.
\]
\textbf{Conclusion.} The boats are within radio range from noon until 2:30 p.m.; contact lasts $\mathbf{2.5}$ hours. (Here the relative motion is one-dimensional, so the ``quadratic'' factorises at once --- the general template of the \emph{M\'ethode} box still applies.)
\end{exoresolu}

\begin{exoresolu}[A genuine quadratic: safety distance at a road junction]
Two vehicles approach a crossroads along perpendicular roads. Taking the junction as origin, at $t=0$ (in seconds) car $A$ is at $(-60\vec{i})$ m moving with $\vec{v}_A=12\vec{i}\ \mathrm{m\,s^{-1}}$, and car $B$ is at $(-80\vec{j})$ m moving with $\vec{v}_B=16\vec{j}\ \mathrm{m\,s^{-1}}$. Road-safety rules require the cars never to be less than $R=20$ m apart. Are the rules violated, and if so, during which time interval?
\tcblower
\textbf{Step 1.} $\vec{r}_0=\vec{r}_A(0)-\vec{r}_B(0)=-60\vec{i}+80\vec{j}$ and
\[
\vec{v}_{\mathrm{rel}}=12\vec{i}-16\vec{j},\qquad
\vec{r}_{A/B}(t)=(-60+12t)\vec{i}+(80-16t)\vec{j}.
\]
\textbf{Step 2.} $d^2(t)=(12t-60)^2+(80-16t)^2=144t^2-1440t+3600+256t^2-2560t+6400$:
\[
d^2(t)=400t^2-4000t+10000=400\bigl(t^2-10t+25\bigr)=400(t-5)^2.
\]
So $d(t)=20|t-5|$: the minimum separation is $d_{\min}=0$ at $t=5$ s --- the cars actually collide at the junction! (Check: at $t=5$, $A$ is at the origin and $B$ is at the origin.)

\textbf{Step 3.} The safety rule is violated when $d(t)\leq 20$:
\[
400(t-5)^2\leq 400\iff (t-5)^2\leq 1\iff 4\leq t\leq 6.
\]
\textbf{Conclusion.} The rule is violated (indeed catastrophically) between $t=4$ s and $t=6$ s; the cars are within 20 m of each other for 2 seconds, colliding at $t=5$ s. This is exactly the kind of computation an engineer uses when timing traffic lights at a Douala intersection.
\end{exoresolu}

\begin{popfigure}
\begin{center}
\begin{tikzpicture}
\begin{axis}[
  width=11cm, height=6.5cm,
  xlabel={$t$ (s)}, ylabel={$d(t)$ (m)},
  xmin=0, xmax=10, ymin=0, ymax=110,
  xtick={0,1,2,3,4,5,6,7,8,9,10},
  ytick={0,20,40,60,80,100},
  grid=both, grid style={popline, very thin},
  axis line style={popdark},
]
\addplot[domain=0:10, samples=100, very thick, popblue] {20*abs(x-5)};
\addplot[domain=0:10, samples=2, thick, poporange, dashed] {20};
\addplot[only marks, mark=*, poppurple] coordinates {(4,20) (6,20)};
\addplot[domain=4:6, samples=2, very thick, popgreen] {20*abs(x-5)};
\node[poporange, anchor=south west] at (axis cs:0.2,21) {$d=R=20$ m};
\node[popgreen, anchor=north] at (axis cs:5,2) {within range: $4\le t\le 6$};
\node[popblue, anchor=west] at (axis cs:6.3,55) {$d(t)=20|t-5|$};
\end{axis}
\end{tikzpicture}
\legende{Separation $d(t)$ for the junction example. The horizontal line $d=R$ cuts the curve at $t=4$ and $t=6$: between these times (green portion) the vehicles are within the critical distance.}
\end{center}
\end{popfigure}

\begin{attention}[Squaring and signs]
Two traps when solving $d(t)\leq R$:
\begin{itemize}
\item Squaring $d(t)\leq R$ is valid only because both sides are non-negative; never square an inequality whose sides may be negative.
\item The inequality $at^2+bt+c\leq 0$ with $a>0$ holds \emph{between} the roots, not outside them. Reversing the interval is one of the most frequent GCE errors in this chapter.
\end{itemize}
\end{attention}

\courssub{Applications: Vehicles, Boats and Aircraft}

The same template --- one relative vector, one quadratic --- solves all the standard situations:

\begin{itemize}
\item \textbf{Road junctions.} Two vehicles on perpendicular or oblique roads: $\vec{r}_0$ and $\vec{v}_{\mathrm{rel}}$ come straight from the geometry; the safety distance gives $R$.
\item \textbf{Boats on the Wouri.} With a current, remember (Lesson 104) that each boat's \emph{true} velocity $\vec{v}=\vec{v}_{\text{boat/water}}+\vec{v}_{\text{water}}$ must be used before forming $\vec{v}_{\mathrm{rel}}$. If both boats are in the \emph{same} current, the current cancels in $\vec{v}_A-\vec{v}_B$: a beautiful simplification worth noticing in the exam.
\item \textbf{Aircraft on crossing routes.} Air-traffic control requires a minimum horizontal separation (a few nautical miles). With constant velocities, checking compliance is exactly $d_{\min}\geq R$, computed from $d^2(t)$ without ever solving for $t$ explicitly: evaluate $d_{\min}^2=|\vec{r}_0|^2-\dfrac{(\vec{r}_0\cdot\vec{v}_{\mathrm{rel}})^2}{|\vec{v}_{\mathrm{rel}}|^2}$ and compare with $R^2$.
\end{itemize}

\begin{exemplebox}[Aircraft on crossing routes]
Aircraft $A$ flies at $(300\vec{i})\ \mathrm{km\,h^{-1}}$ and aircraft $B$ at $(240\vec{j})\ \mathrm{km\,h^{-1}}$ (same altitude). At a certain instant $A$ is 50 km west and 40 km south of $B$. Regulations require a minimum separation of 13 km. Is the regulation respected?
\[
\vec{r}_0=\vec{r}_A-\vec{r}_B=-50\vec{i}-40\vec{j},\qquad
\vec{v}_{\mathrm{rel}}=300\vec{i}-240\vec{j}.
\]
$\vec{r}_0\cdot\vec{v}_{\mathrm{rel}}=-15000+9600=-5400$ (negative: separation decreasing). $|\vec{v}_{\mathrm{rel}}|^2=90000+57600=147600$.
\[
d_{\min}^2=|\vec{r}_0|^2-\frac{(\vec{r}_0\cdot\vec{v}_{\mathrm{rel}})^2}{|\vec{v}_{\mathrm{rel}}|^2}
=4100-\frac{5400^2}{147600}=4100-197.6\approx 3902.4,
\]
so $d_{\min}\approx 62.5$ km $\gg 13$ km: the regulation is comfortably respected.
\end{exemplebox}

\courssub{Chapter Synthesis}

\begin{aretenir}[The whole chapter in one table]
\begin{center}
\begin{tabularx}{\linewidth}{|l|X|X|}
\hline
\rowcolor{popblueL} \textbf{Notion} & \textbf{Question it answers} & \textbf{Key formula} \\
\hline
Relative velocity & How does $A$ move as seen from $B$? & $\vec{v}_{A/B}=\vec{v}_A-\vec{v}_B$ \\
\hline
True velocity (river, wind) & What is the actual velocity over the ground? & $\vec{v}_{\text{true}}=\vec{v}_{\text{body/medium}}+\vec{v}_{\text{medium}}$ \\
\hline
Resultant velocity & Combined effect of non-parallel influences & Vector sum; magnitude by cosine rule \\
\hline
Interception / collision & Do the bodies meet? When and where? & $\vec{r}_A(t)=\vec{r}_B(t)$, i.e.\ $\vec{r}_0+\vec{v}_{\mathrm{rel}}t=\vec{0}$ \\
\hline
Closest approach & How near do they get? & $t^*=-\dfrac{\vec{r}_0\cdot\vec{v}_{\mathrm{rel}}}{|\vec{v}_{\mathrm{rel}}|^2}$ \\
\hline
Relative position / motion & Where is $A$ as seen from $B$ at time $t$? & $\vec{r}_{A/B}(t)=\vec{r}_0+\vec{v}_{\mathrm{rel}}t$;\quad $d(t)=|\vec{r}_{A/B}(t)|$ \\
\hline
Time within range & When is $d(t)\leq R$? & Solve $|\vec{r}_0+\vec{v}_{\mathrm{rel}}t|^2\leq R^2$ (quadratic in $t$) \\
\hline
\end{tabularx}
\end{center}
Every problem in this chapter reduces to: \textbf{(1)} choose an origin and axes, \textbf{(2)} form $\vec{r}_0$ and $\vec{v}_{\mathrm{rel}}$, \textbf{(3)} answer the question using the row of the table that matches it.
\end{aretenir}

\begin{attention}[Exam tip: the standard template]
Markers' reports on GCE mechanics papers show that most questions on this chapter reduce to forming \emph{one relative vector and one quadratic}. Practise this template until it is automatic:
\begin{enumerate}
\item Diagram with origin, axes, and the two position vectors at $t=0$.
\item $\vec{r}_0$ and $\vec{v}_{\mathrm{rel}}$ in components (watch the signs: relative velocity is a \emph{difference}).
\item $d^2(t)$ expanded as $at^2+bt+c$.
\item Read off what is asked: roots (collision or range), vertex (closest approach), or discriminant (does it happen at all?).
\end{enumerate}
A clean diagram and correct signs are worth more than clever algebra.
\end{attention}

\courssub{Revision and Timed GCE Practice}

Work the following exercises under examination conditions: allow about 25 minutes each, with full working shown.

\begin{exercice}[Exercise 1 --- Relative position and closest approach]
At $t=0$, ship $P$ is at the point $20\vec{i}$ km and sails with velocity $(-4\vec{i}+3\vec{j})\ \mathrm{km\,h^{-1}}$; ship $Q$ is at the origin and sails with velocity $(5\vec{i}+3\vec{j})\ \mathrm{km\,h^{-1}}$.
\begin{enumerate}
\item Find $\vec{r}_{P/Q}(t)$ and show that the separation satisfies $d^2(t)=81t^2-360t+400$.
\item Find the time of closest approach and the minimum distance between the ships.
\item The ships can communicate by radio within 8 km. Find the time interval during which communication is possible.
\end{enumerate}
\end{exercice}

\begin{corrige}[Exercise 1]
\textbf{(a)} $\vec{r}_0=20\vec{i}$, $\vec{v}_{\mathrm{rel}}=\vec{v}_P-\vec{v}_Q=-9\vec{i}$, so $\vec{r}_{P/Q}(t)=(20-9t)\vec{i}$. Then $d^2(t)=(20-9t)^2=81t^2-360t+400$.

\textbf{(b)} Minimum at $t^*=\dfrac{360}{2\times 81}=\dfrac{20}{9}\approx 2.22$ h; $d_{\min}=\left|20-9\times\dfrac{20}{9}\right|=0$ km: the ships in fact collide (the relative velocity points exactly along the initial line of sight).

\textbf{(c)} $(20-9t)^2\leq 64\iff |20-9t|\leq 8\iff \dfrac{12}{9}\leq t\leq\dfrac{28}{9}$, i.e.\ from $t=1\ \mathrm{h}\ 20\ \mathrm{min}$ to $t=3\ \mathrm{h}\ 7\ \mathrm{min}$ (approximately); contact lasts $\dfrac{16}{9}$ h $\approx 1$ h $47$ min.
\end{corrige}

\begin{exercice}[Exercise 2 --- GCE style: boats and a safety zone]
Two canoes cross the Wouri. Relative to the ground, canoe $A$ has velocity $(4\vec{i}+2\vec{j})\ \mathrm{m\,s^{-1}}$ and starts at the origin; canoe $B$ has velocity $(\vec{i}+6\vec{j})\ \mathrm{m\,s^{-1}}$ and starts at $90\vec{i}$ m.
\begin{enumerate}
\item Show that $\vec{r}_{A/B}(t)=(-90+3t)\vec{i}+(-4t)\vec{j}$.
\item Find the minimum distance between the canoes and the time at which it occurs.
\item A floating marker is safe only if both canoes stay more than 30 m apart. Do the canoes ever come within 30 m of each other? Justify.
\end{enumerate}
\end{exercice}

\begin{corrige}[Exercise 2]
\textbf{(a)} $\vec{r}_0=-90\vec{i}$, $\vec{v}_{\mathrm{rel}}=3\vec{i}-4\vec{j}$, hence $\vec{r}_{A/B}(t)=(-90+3t)\vec{i}-4t\vec{j}$.

\textbf{(b)} $d^2(t)=(3t-90)^2+16t^2=25t^2-540t+8100$. Minimum at $t^*=\dfrac{540}{50}=10.8$ s; $d_{\min}^2=8100-\dfrac{540^2}{100}=8100-2916=5184$, so $d_{\min}=72$ m.

\textbf{(c)} Since $d_{\min}=72$ m $>30$ m, the canoes never come within 30 m of each other: the discriminant of $25t^2-540t+8100=900$ is negative, confirming there is no solution.
\end{corrige}

\begin{savaistu}[Why pilots and captains think in relative motion]
Air-traffic controllers and ship pilots do not track every craft against the map; they track \emph{relative} positions and \emph{closest points of approach} (CPA), exactly the $d_{\min}$ of this chapter. Modern collision-avoidance systems compute $t^*$ and $d_{\min}$ continuously and sound an alarm when $d_{\min}$ falls below a threshold --- your quadratic $at^2+bt+c$, evaluated thousands of times per second.
\end{savaistu}

With this section the chapter is complete. Relative velocity, true velocity, interception, closest approach and relative position are not ten separate tricks but one single idea --- \emph{change the observer} --- applied with one tool, the vector $\vec{r}_0+\vec{v}_{\mathrm{rel}}t$. Master the template, and this whole family of GCE questions becomes routine.

\newpage

\coursec{Integration activity}

\courssub{Aims of the activity}

This integration activity, \emph{Operation Wouri Patrol}, mobilises in a single authentic case study all the competences developed in this chapter. By the end of the activity, you should be able to:

\begin{itemize}
\item express velocities in component form $\vec{v}=v_x\vec{i}+v_y\vec{j}$ and interpret each component physically;
\item compute the \cle{true velocity} of a boat moving in a current by vector addition of the velocity in still water and the velocity of the current;
\item use the concept of \cle{relative velocity} to determine the course to steer for an \cle{interception}, and compute the interception time and position;
\item analyse a non-intercepting pursuit by computing the \cle{closest distance of approach} using relative position and relative velocity;
\item communicate scientific results clearly in the form of a structured mission report with vector diagrams.
\end{itemize}

The activity is done in groups of three or four. Allow about two hours: one hour of group work, then a plenary session where groups compare their answers and defend their methods.

\courssub{Situation}

The Wouri estuary at Douala is one of the busiest waterways in Central Africa. The Maritime Patrol Unit based near the Port of Douala operates fast patrol boats that can make a speed of $12\ \mathrm{m\,s^{-1}}$ in still water. On the day of the mission, the Wouri current flows steadily towards the sea at $2\ \mathrm{m\,s^{-1}}$.

Set up axes fixed to the ground with:
\begin{itemize}
\item unit vector $\vec{i}$ pointing \textbf{across} the estuary (from the Douala bank towards the Bonabéri side);
\item unit vector $\vec{j}$ pointing \textbf{downstream} (towards the sea).
\end{itemize}

The patrol boat $P$ is at the origin $O$ when, at time $t=0$, the radar operator sights a suspect boat $S$ exactly $3\ \text{km}$ downstream, i.e. at position
\[\vec{r}_S(0)=\begin{pmatrix}0\\3000\end{pmatrix}\ \text{m}.\]
Radar tracking shows that the suspect boat is moving with constant velocity
\[\vec{v}_S=(4\vec{i}+1\vec{j})\ \mathrm{m\,s^{-1}},\]
that is, it is crossing the estuary while drifting slowly downstream. The suspect boat keeps this velocity throughout. The current velocity is $\vec{v}_C=2\vec{j}\ \mathrm{m\,s^{-1}}$, and the patrol boat's speed relative to the water is $12\ \mathrm{m\,s^{-1}}$.

\begin{popfigure}
\begin{center}
\begin{tikzpicture}[scale=1.0,>=stealth]
% water
\fill[popblueL] (0,0) rectangle (7,4.6);
% axes
\draw[->,thick,popink] (0.4,0.4) -- (2.0,0.4) node[below left,popink]{\small $\vec{i}$ (across)};
\draw[->,thick,popink] (0.4,0.4) -- (0.4,2.0) node[below left,popink]{\small $\vec{j}$ (downstream)};
% patrol boat
\fill[popgreen] (0.4,0.4) circle (0.09);
\node[popgreen,below right] at (0.4,0.32) {\small Patrol $P$};
% suspect boat
\fill[poporange] (0.4,3.9) circle (0.09);
\node[poporange,right] at (0.5,3.9) {\small Suspect $S$};
% 3 km
\draw[<->,popmuted] (0.15,0.4) -- (0.15,3.9) node[midway,left,popmuted]{\small $3$ km};
% suspect velocity
\draw[->,very thick,poporange] (0.4,3.9) -- (2.0,4.15) node[right,poporange]{\small $\vec{v}_S=4\vec{i}+\vec{j}$};
% current
\draw[->,very thick,popblue] (5.6,1.2) -- (5.6,2.4) node[right,popblue]{\small $\vec{v}_C=2\vec{j}$};
% patrol speed label
\node[popgreen,align=center] at (4.6,0.6) {\small $12\ \mathrm{m\,s^{-1}}$ in still water};
\end{tikzpicture}
\end{center}
\legende{Situation map of Operation Wouri Patrol (not to scale).}
\end{popfigure}

The commander asks your team four operational questions: What is the patrol boat's true velocity over the ground for a given heading? What course must be steered to intercept the suspect boat, and when and where will interception occur? If instead the helmsman simply steers straight at the sighted position, how close does the patrol boat actually get? Your group must deliver a one-page \textbf{mission report}.

\courssub{Tasks}

\begin{enumerate}
\item \textbf{Component form.} Write down, in $\vec{i},\vec{j}$ form: the velocity $\vec{v}_S$ of the suspect boat, the velocity $\vec{v}_C$ of the current, and the position vector $\vec{r}_S(0)$ of the suspect boat at $t=0$. State the position vector $\vec{r}_S(t)$ of the suspect boat at time $t$.
\item \textbf{True velocity for a chosen heading.} The patrol boat first steers a heading making an angle $\theta$ with $\vec{i}$ (measured towards $\vec{j}$), so its velocity relative to the water is $12(\cos\theta\,\vec{i}+\sin\theta\,\vec{j})$. Show that its true velocity over the ground is
\[\vec{v}_P=12\cos\theta\,\vec{i}+(12\sin\theta+2)\,\vec{j}.\]
Compute $\vec{v}_P$ and the true speed $|\vec{v}_P|$ for the heading $\theta=90^{\circ}$ (straight downstream).
\item \textbf{Course to intercept.} For interception, the relative velocity $\vec{v}_{P/S}=\vec{v}_P-\vec{v}_S$ must be directed along the initial relative position $\vec{r}_S(0)-\vec{r}_P(0)$, i.e. along $\vec{j}$.
\begin{enumerate}
\item Explain why this requires the $\vec{i}$-component of $\vec{v}_{P/S}$ to be zero.
\item Deduce that $\cos\theta=\tfrac{1}{3}$ and find $\theta$ to the nearest degree.
\item Compute the interception time $t_1$ and the position vector of the interception point.
\end{enumerate}
\item \textbf{Steering straight at the sighted position.} Suppose instead the helmsman steers the heading $\theta=90^{\circ}$ directly towards the sighted position $(0,3000)$.
\begin{enumerate}
\item Write down $\vec{v}_P$ and the relative velocity $\vec{v}_{P/S}$.
\item Show that the relative position is $\vec{r}_{S/P}(t)=\begin{pmatrix}4t\\3000-13t\end{pmatrix}$.
\item Find the time $t_2$ at which the distance between the boats is least, and the closest distance of approach $d_{\min}$. Will the patrol boat catch the suspect boat by this tactic?
\end{enumerate}
\item \textbf{Mission report.} Prepare a one-page report containing: the data summary, all velocity vectors in component form, a vector diagram of the interception triangle, your answers with units, and one sentence of operational advice to the commander.
\end{enumerate}

\courssub{Model answer}

\textbf{Task 1.} $\vec{v}_S=(4\vec{i}+\vec{j})\ \mathrm{m\,s^{-1}}$, $\vec{v}_C=2\vec{j}\ \mathrm{m\,s^{-1}}$, $\vec{r}_S(0)=3000\vec{j}\ \text{m}$. Since $\vec{v}_S$ is constant,
\[\vec{r}_S(t)=\vec{r}_S(0)+\vec{v}_S\,t=4t\,\vec{i}+(3000+t)\,\vec{j}\ \text{(m)}.\]

\textbf{Task 2.} The true velocity is the sum of the velocity relative to the water and the velocity of the water:
\[\vec{v}_P=12(\cos\theta\,\vec{i}+\sin\theta\,\vec{j})+2\vec{j}=12\cos\theta\,\vec{i}+(12\sin\theta+2)\vec{j}.\]
For $\theta=90^{\circ}$: $\vec{v}_P=14\vec{j}\ \mathrm{m\,s^{-1}}$ and $|\vec{v}_P|=14\ \mathrm{m\,s^{-1}}$.

\textbf{Task 3.} (a) The initial relative position is $3000\vec{j}$, purely along $\vec{j}$. For the gap to close along that fixed line, the relative velocity must be parallel to it; hence its $\vec{i}$-component must vanish (otherwise the patrol boat drifts across and misses).

(b) $\vec{v}_{P/S}=(12\cos\theta-4)\vec{i}+(12\sin\theta+2-1)\vec{j}$. Setting the $\vec{i}$-component to zero:
\[12\cos\theta-4=0\quad\Longrightarrow\quad \cos\theta=\tfrac{1}{3}\quad\Longrightarrow\quad \theta\approx 70.5^{\circ}\approx 71^{\circ}.\]
The patrol boat must steer about $71^{\circ}$ from the across-estuary direction, i.e. mostly downstream, angled slightly across so that its across-component $4\ \mathrm{m\,s^{-1}}$ exactly matches the suspect boat's.

(c) With $\sin\theta=\sqrt{1-\tfrac19}=\tfrac{2\sqrt2}{3}$, the closing speed along $\vec{j}$ is
\[12\sin\theta+1=8\sqrt2+1\approx 12.31\ \mathrm{m\,s^{-1}}.\]
Hence
\[t_1=\frac{3000}{8\sqrt2+1}\approx 243.6\ \text{s}\approx 4\ \text{min}\ 4\ \text{s}.\]
Interception point: $\vec{r}_S(t_1)=4t_1\vec{i}+(3000+t_1)\vec{j}\approx 975\vec{i}+3244\vec{j}\ \text{(m)}$.

\begin{attention}
A frequent error is to aim the heading so that the \emph{true} velocity has no $\vec{i}$-component. That is wrong: the patrol boat must keep an across-component of $4\ \mathrm{m\,s^{-1}}$ to \emph{match} the suspect boat's motion across the estuary. Interception is a condition on the \textbf{relative} velocity, not on the true velocity alone.
\end{attention}

\textbf{Task 4.} (a) With $\theta=90^{\circ}$: $\vec{v}_P=14\vec{j}$, so
\[\vec{v}_{P/S}=\vec{v}_P-\vec{v}_S=-4\vec{i}+13\vec{j}\ \mathrm{m\,s^{-1}}.\]

(b) $\vec{r}_{S/P}(t)=\vec{r}_S(0)+\vec{v}_{S/P}\,t=3000\vec{j}+(4\vec{i}-13\vec{j})t=4t\,\vec{i}+(3000-13t)\vec{j}$.

(c) The squared distance is
\[d^2(t)=(4t)^2+(3000-13t)^2=16t^2+(3000-13t)^2.\]
Differentiating: $\dfrac{d(d^2)}{dt}=32t-26(3000-13t)=370t-78000$. Setting this to zero gives
\[t_2=\frac{78000}{370}=\frac{7800}{37}\approx 210.8\ \text{s}\approx 3\ \text{min}\ 31\ \text{s}.\]
Since $370>0$, $d^2$ is a minimum there. Substituting:
\[d_{\min}^2=\left(\frac{31200}{37}\right)^2+\left(\frac{9600}{37}\right)^2=\frac{1\,065\,600\,000}{1369},\]
\[d_{\min}\approx 882\ \text{m}.\]
\emph{Check by the cross-product formula:} $d_{\min}=\dfrac{|\vec{r}_0\wedge\vec{v}_{S/P}|}{|\vec{v}_{S/P}|}=\dfrac{|0\times(-13)-3000\times4|}{\sqrt{4^2+13^2}}=\dfrac{12000}{\sqrt{185}}\approx 882\ \text{m}$. $\checkmark$

The boats never get closer than about $0.88\ \text{km}$: steering straight at the sighted position \textbf{fails} to intercept, because the suspect boat slips across the estuary while the patrol boat's heading ignores that motion.

\textbf{Task 5 (report skeleton).} Data: $|\vec{v}_{P/\text{water}}|=12$, $\vec{v}_C=2\vec{j}$, $\vec{v}_S=4\vec{i}+\vec{j}$, $\vec{r}_S(0)=3000\vec{j}$. Result: steer $\theta\approx71^{\circ}$ from $\vec{i}$; interception after $\approx244\ \text{s}$ at $\approx(975,\,3244)\ \text{m}$. Steering straight at the target gives only $d_{\min}\approx882\ \text{m}$. \emph{Advice to the commander: always steer an interception course that matches the target's across-estuary component; chasing the sighted position lets the suspect escape.}

\begin{savaistu}
River and estuary patrols, ferry pilots on the Wouri, and even fishermen crossing to Manoka all use — consciously or not — the same vector addition you have just applied: the heading steered is never the track actually followed when a current is present.
\end{savaistu}

In the plenary session, compare the interception times obtained by different groups, discuss the effect of a stronger current on $\theta$ and $t_1$, and justify why the closest-approach tactic fails even though the patrol boat is much faster than the suspect boat.

\newpage

\coursec{Methods and worked examples}

\coursec{Velocity and Acceleration as Vectors}

\courssub{Velocity and Acceleration as Vectors; Concept of Relative Velocity}

\begin{definition}[Velocity and acceleration vectors]
In a Cartesian frame with unit vectors $\vec{i}$ (East) and $\vec{j}$ (North), the \cle{velocity vector} of a particle at time $t$ is
\[ \vec{v}(t) = \frac{d\vec{r}}{dt} = \dot{x}(t)\,\vec{i} + \dot{y}(t)\,\vec{j}, \]
and the \cle{acceleration vector} is
\[ \vec{a}(t) = \frac{d\vec{v}}{dt} = \ddot{x}(t)\,\vec{i} + \ddot{y}(t)\,\vec{j}. \]
The \cle{speed} is the magnitude $|\vec{v}| = \sqrt{\dot{x}^2 + \dot{y}^2}$.
\end{definition}

\begin{propriete}[Componentwise operations]
For vectors $\vec{u} = u_x\vec{i} + u_y\vec{j}$ and $\vec{v} = v_x\vec{i} + v_y\vec{j}$:
\begin{itemize}
\item Addition: $\vec{u} + \vec{v} = (u_x+v_x)\vec{i} + (u_y+v_y)\vec{j}$.
\item Magnitude: $|\vec{v}| = \sqrt{v_x^2+v_y^2}$.
\item Direction: angle $\theta$ from $\vec{i}$ satisfies $\tan\theta = \dfrac{v_y}{v_x}$ (choose quadrant correctly).
\item Differentiation: $\dfrac{d\vec{v}}{dt} = \dfrac{dv_x}{dt}\vec{i} + \dfrac{dv_y}{dt}\vec{j}$.
\end{itemize}
\end{propriete}

\begin{methode}[Working with velocity and acceleration as vectors]
\begin{enumerate}
\item \textbf{Choose a frame.} Fix unit vectors $\vec{i}$ (East) and $\vec{j}$ (North), or $\vec{i}$ along the motion and $\vec{j}$ perpendicular to it. State your choice clearly.
\item \textbf{Write each vector in components.} A velocity of magnitude $v$ in a direction making an angle $\theta$ with $\vec{i}$ is
\[ \vec{v} = v\cos\theta\,\vec{i} + v\sin\theta\,\vec{j}. \]
\item \textbf{Add componentwise.} If $\vec{v} = \vec{v}_1 + \vec{v}_2$, add the $\vec{i}$-parts and the $\vec{j}$-parts separately.
\item \textbf{Recover magnitude and direction.} If $\vec{v} = a\vec{i} + b\vec{j}$, then
\[ |\vec{v}| = \sqrt{a^2+b^2}, \qquad \tan\theta = \frac{b}{a}. \]
Always state the direction as a bearing or as an angle from a named axis.
\item \textbf{Differentiate componentwise.} Acceleration is $\vec{a} = \dfrac{d\vec{v}}{dt}$; differentiate each component separately.
\end{enumerate}
\textbf{Reflex:} never add magnitudes unless the vectors are parallel. Always draw a quick sketch before computing.
\end{methode}

\begin{exoresolu}[Components, magnitude and direction]
A particle moves with velocity $\vec{v} = (3t^2\,\vec{i} + 4t\,\vec{j})\ \mathrm{m\,s^{-1}}$ at time $t$ seconds.
\begin{itemize}
\item[(a)] Find its speed and the direction of motion when $t = 2$.
\item[(b)] Find its acceleration at time $t$ and determine whether the acceleration has constant direction.
\end{itemize}
\tcblower
\textbf{(a)} At $t = 2$:
\[ \vec{v} = 3(2)^2\,\vec{i} + 4(2)\,\vec{j} = 12\vec{i} + 8\vec{j}. \]
The speed is the magnitude:
\[ |\vec{v}| = \sqrt{12^2 + 8^2} = \sqrt{144 + 64} = \sqrt{208} = 4\sqrt{13} \approx 14.4\ \mathrm{m\,s^{-1}}. \]
If $\theta$ is the angle the velocity makes with $\vec{i}$:
\[ \tan\theta = \frac{8}{12} = \frac{2}{3} \quad\Longrightarrow\quad \theta \approx 33.7^\circ. \]
So the particle moves at about $14.4\ \mathrm{m\,s^{-1}}$ in the direction $33.7^\circ$ above the $\vec{i}$-axis.

\textbf{(b)} Differentiate each component:
\[ \vec{a} = \frac{d\vec{v}}{dt} = 6t\,\vec{i} + 4\,\vec{j}. \]
The direction of $\vec{a}$ makes an angle $\alpha$ with $\vec{i}$ where $\tan\alpha = \dfrac{4}{6t} = \dfrac{2}{3t}$, which \emph{depends on $t$}. Hence the acceleration does \textbf{not} have constant direction. At $t=2$, $\vec{a} = 12\vec{i} + 4\vec{j}$, $|\vec{a}| = \sqrt{144+16} = \sqrt{160} = 4\sqrt{10} \approx 12.6\ \mathrm{m\,s^{-2}}$, directed at $\tan^{-1}(1/3) \approx 18.4^\circ$ to $\vec{i}$.

\textbf{Key point:} speed $=$ magnitude of the velocity \emph{vector}; never differentiate the speed directly to get acceleration unless the path is straight.
\end{exoresolu}

\begin{exoresolu}[Acceleration with constant direction]
A particle has velocity $\vec{v} = (4t\,\vec{i} + 3t\,\vec{j})\ \mathrm{m\,s^{-1}}$. Show that its acceleration has constant direction and find that direction.
\tcblower
\[ \vec{a} = \frac{d\vec{v}}{dt} = 4\vec{i} + 3\vec{j}. \]
Both components are constant, so the acceleration vector is constant in both magnitude and direction. The direction makes an angle $\theta$ with $\vec{i}$ where $\tan\theta = \frac{3}{4}$, so $\theta \approx 36.9^\circ$.
\end{exoresolu}

\begin{savaistu}[Vector calculus in mechanics]
Newton's second law $\vec{F} = m\vec{a}$ is a vector equation. In the GCE Advanced Level, you must be comfortable switching between vector form ($\vec{F} = m\vec{a}$) and component form ($F_x = m a_x$, $F_y = m a_y$). This chapter builds the vector manipulation skills needed for the Mechanics chapters on forces, projectiles, and circular motion.
\end{savaistu}

\courssub{Non-Parallel Forces, Resultant Velocity and True Velocity in River or Wind Problems}

\begin{definition}[Relative velocity]
The \cle{velocity of $A$ relative to $B$} is
\[ \vec{v}_{A/B} = \vec{v}_A - \vec{v}_B. \]
It is the velocity of $A$ as observed by someone moving with $B$.
\end{definition}

\begin{propriete}[Chain rule for relative velocities]
\[ \vec{v}_{A/C} = \vec{v}_{A/B} + \vec{v}_{B/C}. \]
The middle letter ``cancels''. In particular, $\vec{v}_{A/B} = -\vec{v}_{B/A}$.
\end{propriete}

\begin{methode}[Finding the velocity of one body relative to another]
\begin{enumerate}
\item \textbf{Memorise the defining relation:}
\[ \vec{v}_{A/B} = \vec{v}_A - \vec{v}_B. \]
\item \textbf{Use the chain rule of subscripts:} $\vec{v}_{A/C} = \vec{v}_{A/B} + \vec{v}_{B/C}$.
\item \textbf{To find what an observer sees,} subtract the observer's velocity: velocity of $A$ as seen from $B$ is $\vec{v}_A - \vec{v}_B$.
\item \textbf{Draw the vector triangle.} If $\vec{v}_{A/B} = \vec{v}_A - \vec{v}_B$, then $\vec{v}_A = \vec{v}_B + \vec{v}_{A/B}$: place the vectors head-to-tail.
\item \textbf{Compute} using components or the cosine/sine rule in the triangle, then give magnitude and bearing.
\end{enumerate}
\textbf{Reflex:} for ``apparent'' motion (rain, wind, another car), the answer is always \emph{true velocity minus observer's velocity}.
\end{methode}

\begin{popfigure}
\begin{tikzpicture}[scale=0.8, >=stealth]
% Vector triangle for relative velocity
\draw[->, thick, popblue] (0,0) -- (3,0) node[midway, below] {$\vec{v}_B$};
\draw[->, thick, popgreen] (3,0) -- (3,2) node[midway, right] {$\vec{v}_{A/B}$};
\draw[->, thick, poporange] (0,0) -- (3,2) node[midway, above left] {$\vec{v}_A$};
\draw[dashed] (0,0) -- (3,0) -- (3,2);
\node at (-0.3,0) {$O$};
\node at (3,-0.3) {$\vec{v}_B$};
\node at (3.3,2) {$\vec{v}_A$};
\end{tikzpicture}
\legende{Vector triangle: $\vec{v}_A = \vec{v}_B + \vec{v}_{A/B}$}
\end{popfigure}

\begin{exoresolu}[Apparent direction of rain]
To a cyclist riding due East at $12\ \mathrm{km\,h^{-1}}$, the rain appears to fall vertically downwards. When the cyclist doubles her speed to $24\ \mathrm{km\,h^{-1}}$ (still due East), the rain appears to come at $30^\circ$ to the vertical, hitting her from the front. Find the true velocity of the rain.
\tcblower
Let the true velocity of the rain be $\vec{v}_R = x\vec{i} + y\vec{j}$, with $\vec{i}$ East and $\vec{j}$ vertically up (so $y<0$). The cyclist's velocity is $\vec{v}_C = u\vec{i}$.

The apparent (relative) velocity is $\vec{v}_{R/C} = \vec{v}_R - \vec{v}_C = (x-u)\vec{i} + y\vec{j}$.

\textbf{First ride} ($u = 12$): rain appears vertical, so the horizontal component of $\vec{v}_{R/C}$ is zero:
\[ x - 12 = 0 \quad\Longrightarrow\quad x = 12. \]

\textbf{Second ride} ($u = 24$): the apparent velocity is $(12-24)\vec{i} + y\vec{j} = -12\vec{i} + y\vec{j}$, directed at $30^\circ$ to the vertical:
\[ \tan 30^\circ = \frac{\text{horizontal}}{\text{vertical}} = \frac{12}{|y|} \quad\Longrightarrow\quad |y| = \frac{12}{\tan 30^\circ} = 12\sqrt{3}. \]
Since the rain falls, $y = -12\sqrt{3}$.

\textbf{True velocity of the rain:}
\[ \vec{v}_R = 12\vec{i} - 12\sqrt{3}\,\vec{j}, \qquad |\vec{v}_R| = \sqrt{144 + 432} = \sqrt{576} = 24\ \mathrm{km\,h^{-1}}. \]
Direction: $\tan\theta = \dfrac{12\sqrt{3}}{12} = \sqrt{3}$, so $\theta = 60^\circ$ below the horizontal, i.e. the rain truly falls at $24\ \mathrm{km\,h^{-1}}$ at $60^\circ$ to the horizontal, blowing from the East.

\textbf{Check:} at $u=12$, apparent velocity $= -12\sqrt{3}\,\vec{j}$: purely vertical. $\checkmark$
\end{exoresolu}

\begin{exoresolu}[Relative velocity of two ships]
At noon, ship $P$ is at the origin and sails with velocity $(12\vec{i} + 16\vec{j})\ \mathrm{km\,h^{-1}}$. Ship $Q$ is at $(40\vec{i} + 10\vec{j})\ \mathrm{km}$ and sails with velocity $(-8\vec{i} + 6\vec{j})\ \mathrm{km\,h^{-1}}$. Find the velocity of $P$ relative to $Q$ and the distance of closest approach.
\tcblower
\[ \vec{v}_{P/Q} = \vec{v}_P - \vec{v}_Q = (12 - (-8))\vec{i} + (16 - 6)\vec{j} = 20\vec{i} + 10\vec{j}\ \mathrm{km\,h^{-1}}. \]
Initial relative position: $\vec{r}_{P/Q}(0) = \vec{r}_P(0) - \vec{r}_Q(0) = -40\vec{i} - 10\vec{j}$.
Relative position at time $t$: $\vec{r}_{P/Q}(t) = (-40 + 20t)\vec{i} + (-10 + 10t)\vec{j}$.
Squared distance: $d^2 = (20t-40)^2 + (10t-10)^2 = 500t^2 - 1800t + 1700$.
Minimise: $\frac{d}{dt}(d^2) = 1000t - 1800 = 0 \Rightarrow t = 1.8\ \mathrm{h}$.
Minimum distance: $d_{\min} = \sqrt{500(1.8)^2 - 1800(1.8) + 1700} = \sqrt{80} = 4\sqrt{5} \approx 8.94\ \mathrm{km}$.
\end{exoresolu}

\begin{definition}[Resultant velocity and true velocity]
When a body moves in a medium that is itself moving (boat in a current, aircraft in wind), its \cle{own velocity} (or \cle{velocity in still water/air}) is $\vec{v}_{\text{own}}$. The \cle{velocity of the medium} is $\vec{v}_{\text{medium}}$. The \cle{resultant velocity} (also called \cle{true velocity} or \cle{ground velocity}) is
\[ \vec{v}_{\text{true}} = \vec{v}_{\text{own}} + \vec{v}_{\text{medium}}. \]
The actual path over ground follows $\vec{v}_{\text{true}}$, not the heading.
\end{definition}

\begin{methode}[Combining non-parallel velocities (own motion + current/wind)]
\begin{enumerate}
\item \textbf{Identify the contributions:} the body's own velocity (engine, rowing, flying) and the velocity of the medium (current, wind).
\item \textbf{Resultant velocity = sum of the vectors:} $\vec{v}_{\text{true}} = \vec{v}_{\text{own}} + \vec{v}_{\text{medium}}$.
\item \textbf{Resolve} both vectors along two perpendicular axes (usually along and across the river, or East/North).
\item \textbf{Add components,} then recombine with Pythagoras and $\tan\theta$.
\item \textbf{Alternatively use the vector triangle} with the cosine rule: $v^2 = v_1^2 + v_2^2 + 2v_1v_2\cos\phi$ where $\phi$ is the angle between the two vectors placed head-to-tail.
\end{enumerate}
\textbf{Reflex:} the actual path of the boat/plane follows the \emph{resultant} velocity, not the heading.
\end{methode}

\begin{popfigure}
\begin{tikzpicture}[scale=0.7, >=stealth]
% River crossing diagram
\draw[popblue, thick] (0,0) -- (0,4) node[midway, left] {Bank};
\draw[popblue, thick] (5,0) -- (5,4) node[midway, right] {Bank};
\draw[->, thick, poporange] (0,1) -- (2,3) node[midway, above right] {$\vec{v}_{\text{true}}$};
\draw[->, thick, popgreen] (0,1) -- (0,3) node[midway, left] {$\vec{v}_{\text{own}}$};
\draw[->, thick, popblue] (0,3) -- (2,3) node[midway, above] {$\vec{v}_{\text{current}}$};
\draw[dashed] (0,1) -- (0,3) -- (2,3);
\node at (0,0.5) {Start};
\node at (2,3.5) {Landing point};
\end{tikzpicture}
\legende{Velocity triangle for a boat crossing a river}
\end{popfigure}

\begin{exoresolu}[Resultant velocity of a boat]
A pirogue on the Wouri river can move at $5\ \mathrm{m\,s^{-1}}$ in still water. The current flows due East at $3\ \mathrm{m\,s^{-1}}$. The boatman heads the pirogue due North. Find the resultant velocity of the pirogue (magnitude and direction), and the time to cross a river $400\ \mathrm{m}$ wide.
\tcblower
Take $\vec{i}$ East, $\vec{j}$ North.
\[ \vec{v}_{\text{own}} = 5\vec{j}, \qquad \vec{v}_{\text{current}} = 3\vec{i}. \]
Resultant:
\[ \vec{v}_{\text{true}} = 3\vec{i} + 5\vec{j}. \]
Magnitude:
\[ |\vec{v}_{\text{true}}| = \sqrt{3^2 + 5^2} = \sqrt{34} \approx 5.83\ \mathrm{m\,s^{-1}}. \]
Direction: $\tan\theta = \dfrac{3}{5}$ where $\theta$ is the angle East of North:
\[ \theta = \tan^{-1}(0.6) \approx 31.0^\circ. \]
So the pirogue actually travels at $5.83\ \mathrm{m\,s^{-1}}$ on a course N$31^\circ$E, even though it is pointed due North.

\textbf{Crossing time.} Only the component \emph{across} the river moves the boat across:
\[ t = \frac{\text{width}}{\text{component across}} = \frac{400}{5} = 80\ \mathrm{s}. \]
(The current does not help or hinder the crossing; it only carries the boat downstream by $3 \times 80 = 240\ \mathrm{m}$.)

\textbf{Key point:} crossing time depends only on the component of velocity perpendicular to the banks.
\end{exoresolu}

\courssub{Applications of Relative and True Velocity}

\begin{methode}[Finding the heading that gives a required true course]
\begin{enumerate}
\item \textbf{Write the vector equation:}
\[ \vec{v}_{\text{true}} = \vec{v}_{\text{own}} + \vec{v}_{\text{medium}}. \]
\item \textbf{Draw the triangle:} from a point $O$, draw $\vec{v}_{\text{medium}}$; the true velocity must close the triangle along the required direction.
\item \textbf{Use the sine rule} in the velocity triangle to find the angle between the heading and the required course:
\[ \frac{v_{\text{medium}}}{\sin\alpha} = \frac{v_{\text{own}}}{\sin\beta}, \]
where $\alpha$ is the angle to steer upstream of the target direction.
\item \textbf{Steer upstream / into the wind} by the angle $\alpha$; the drift then exactly cancels.
\item \textbf{Compute the true speed} with the cosine rule or by resolving along the required direction.
\end{enumerate}
\textbf{Reflex:} to go straight across a river, aim \emph{upstream}; to cross in the shortest time, aim \emph{straight across} (and accept the drift).
\end{methode}

\begin{popfigure}
\begin{tikzpicture}[scale=0.7, >=stealth]
% Velocity triangle for crossing straight across
\draw[->, thick, popblue] (0,0) -- (3,0) node[midway, below] {$\vec{v}_{\text{current}}$};
\draw[->, thick, poporange] (0,0) -- (1.5,2.6) node[midway, left] {$\vec{v}_{\text{own}}$};
\draw[->, thick, popgreen] (3,0) -- (1.5,2.6) node[midway, right] {$\vec{v}_{\text{true}}$};
\draw[dashed] (0,0) -- (3,0) -- (1.5,2.6);
\node at (0,-0.5) {$O$};
\node at (1.5,3.1) {Required direction};
\node at (1.5,-0.5) {$\alpha$};
\end{tikzpicture}
\legende{Velocity triangle: steer upstream to cancel drift}
\end{popfigure}

\begin{exoresolu}[Crossing straight across a flowing river]
A swimmer can swim at $1.2\ \mathrm{m\,s^{-1}}$ in still water. A river flows at $0.9\ \mathrm{m\,s^{-1}}$. In what direction must the swimmer head in order to cross the river perpendicular to the banks? Find her true speed across and the time to cross if the river is $60\ \mathrm{m}$ wide.
\tcblower
Let the swimmer head upstream at an angle $\alpha$ to the perpendicular-to-banks direction. Her own velocity has components:
\begin{itemize}
\item across the river: $1.2\cos\alpha$,
\item upstream (against the current): $1.2\sin\alpha$.
\end{itemize}
For the true path to be straight across, the upstream component must cancel the current:
\[ 1.2\sin\alpha = 0.9 \quad\Longrightarrow\quad \sin\alpha = \frac{0.9}{1.2} = 0.75 \quad\Longrightarrow\quad \alpha \approx 48.6^\circ. \]
She must head at $48.6^\circ$ upstream from the straight-across direction (equivalently $41.4^\circ$ to the bank).

\textbf{True speed across:}
\[ v_{\text{true}} = 1.2\cos\alpha = 1.2\sqrt{1 - 0.75^2} = 1.2\sqrt{0.4375} = 1.2 \times 0.6614 \approx 0.794\ \mathrm{m\,s^{-1}}. \]
(Equivalently, by Pythagoras in the velocity triangle: $v_{\text{true}} = \sqrt{1.2^2 - 0.9^2} = \sqrt{0.63} \approx 0.794\ \mathrm{m\,s^{-1}}$.)

\textbf{Time to cross:}
\[ t = \frac{60}{0.794} \approx 75.6\ \mathrm{s}. \]

\textbf{Remark:} crossing straight across is possible only because the swimmer's speed ($1.2$) exceeds the current ($0.9$). If $v_{\text{own}} < v_{\text{current}}$, no heading gives a straight crossing.
\end{exoresolu}

\begin{exoresolu}[Shortest time vs shortest path across a river]
A boat has speed $4\ \mathrm{m\,s^{-1}}$ in still water. The river is $100\ \mathrm{m}$ wide and flows at $3\ \mathrm{m\,s^{-1}}$.
\begin{itemize}
\item[(a)] Find the shortest time to cross and the landing point.
\item[(b)] Find the heading for a straight crossing and the time taken.
\end{itemize}
\tcblower
\textbf{(a) Shortest time:} head straight across ($\vec{v}_{\text{own}} = 4\vec{j}$).
Time $t = \frac{100}{4} = 25\ \mathrm{s}$.
Drift $= 3 \times 25 = 75\ \mathrm{m}$ downstream.
Resultant speed $= \sqrt{4^2+3^2} = 5\ \mathrm{m\,s^{-1}}$.

\textbf{(b) Straight crossing:} need upstream component $= 3\ \mathrm{m\,s^{-1}}$.
$4\sin\alpha = 3 \Rightarrow \sin\alpha = 0.75 \Rightarrow \alpha \approx 48.6^\circ$ upstream.
Across component $= 4\cos\alpha = 4\sqrt{1-0.75^2} = \sqrt{7} \approx 2.646\ \mathrm{m\,s^{-1}}$.
Time $= \frac{100}{\sqrt{7}} \approx 37.8\ \mathrm{s}$.
\end{exoresolu}

\begin{savaistu}[Cameroon context: River transport]
The Wouri River at Douala and the Sanaga River at Edéa are major transport routes. Pirogue operators must constantly calculate headings to counteract currents, especially during the rainy season when flow rates increase significantly. The same vector principles apply to navigation on Lake Chad and the Benue River in the North.
\end{savaistu}

\courssub{Interception and Collision: Concept and Mathematical Treatment}

\begin{definition}[Interception and collision]
Two moving bodies \cle{intercept} (or \cle{collide}) if their position vectors are equal at the same instant $t$:
\[ \vec{r}_A(t) = \vec{r}_B(t). \]
If one body can choose its direction (e.g. a patrol boat intercepting a target), the condition is that the relative velocity $\vec{v}_{A/B}$ points along the initial line of sight $\vec{r}_B(0) - \vec{r}_A(0)$.
\end{definition}

\begin{methode}[Setting up an interception problem]
\begin{enumerate}
\item \textbf{Write position vectors} of both bodies at time $t$:
\[ \vec{r}_A(t) = \vec{r}_A(0) + \vec{v}_A\, t, \qquad \vec{r}_B(t) = \vec{r}_B(0) + \vec{v}_B\, t. \]
\item \textbf{Interception condition:} they meet when $\vec{r}_A(t) = \vec{r}_B(t)$, i.e. both components equal \emph{at the same time} $t$.
\item \textbf{If one body chases another} and may choose its direction: the interceptor must head so that its relative velocity $\vec{v}_{A/B}$ points along the initial line of sight $\vec{r}_B(0)-\vec{r}_A(0)$. Then the gap closes at rate $|\vec{v}_{A/B}|$.
\item \textbf{Time of interception:} solve the component equations; check that the same $t$ satisfies both.
\item \textbf{Give the meeting point} by substituting $t$ back into either position vector.
\end{enumerate}
\textbf{Reflex:} ``$A$ intercepts $B$'' $\Longleftrightarrow$ relative position $\vec{r}_{A/B} = \vec{r}_A - \vec{r}_B$ becomes $\vec{0}$.
\end{methode}

\begin{popfigure}
\begin{tikzpicture}[scale=0.7, >=stealth]
% Interception geometry
\draw[->, thick, popblue] (0,0) -- (4,3) node[midway, above left] {$\vec{r}_B(0)$};
\draw[->, thick, poporange] (0,0) -- (5,1) node[midway, below] {$\vec{v}_A$};
\draw[->, thick, popgreen] (4,3) -- (6,4) node[midway, above right] {$\vec{v}_B$};
\draw[->, thick, poppurple] (0,0) -- (5,1) node[midway, below right] {$\vec{r}_A(t)$};
\draw[->, thick, poppurple] (4,3) -- (6,4) node[midway, above] {$\vec{r}_B(t)$};
\node at (0,-0.3) {$A(0)$};
\node at (4,2.7) {$B(0)$};
\node at (5,0.7) {Interception point};
\end{tikzpicture}
\legende{Interception geometry: $\vec{v}_{A/B}$ must be along initial line of sight}
\end{popfigure}

\begin{exoresolu}[Interception of a boat by a patrol launch]
At noon, a fishing canoe $B$ is $10\ \mathrm{km}$ due North of a patrol launch $A$ and is travelling due East at $8\ \mathrm{km\,h^{-1}}$. The launch has a maximum speed of $17\ \mathrm{km\,h^{-1}}$. Find the course the launch must steer to intercept the canoe as quickly as possible, and the time of interception.
\tcblower
Take $\vec{i}$ East, $\vec{j}$ North, origin at $A$'s position at noon.
\[ \vec{r}_A(0) = \vec{0}, \qquad \vec{r}_B(0) = 10\vec{j}, \qquad \vec{v}_B = 8\vec{i}. \]
Let the launch steer at angle $\theta$ East of North with speed $17$:
\[ \vec{v}_A = 17\sin\theta\,\vec{i} + 17\cos\theta\,\vec{j}. \]

\textbf{Interception condition} $\vec{r}_A(t) = \vec{r}_B(t)$:
\begin{align*}
17\sin\theta \cdot t &= 8t \quad (\vec{i}\text{-components})\\
17\cos\theta \cdot t &= 10 \quad (\vec{j}\text{-components})
\end{align*}
From the first equation (with $t \neq 0$):
\[ \sin\theta = \frac{8}{17} \quad\Longrightarrow\quad \theta \approx 28.1^\circ. \]
Then $\cos\theta = \dfrac{15}{17}$ (since $8^2+15^2 = 17^2$), and the second equation gives:
\[ 17 \times \frac{15}{17}\, t = 10 \quad\Longrightarrow\quad 15t = 10 \quad\Longrightarrow\quad t = \frac{2}{3}\ \mathrm{h} = 40\ \text{minutes}. \]

\textbf{Answer:} the launch steers N$28.1^\circ$E and intercepts the canoe at $12{:}40$, at the point $\vec{r} = 8 \times \tfrac{2}{3}\,\vec{i} + 10\vec{j} = \tfrac{16}{3}\vec{i} + 10\vec{j}$, i.e. $5\tfrac{1}{3}\ \mathrm{km}$ East and $10\ \mathrm{km}$ North of the launch's starting point.

\textbf{Check (relative velocity method):} $\vec{v}_{A/B} = (17\sin\theta - 8)\vec{i} + 17\cos\theta\,\vec{j} = 0\vec{i} + 15\vec{j}$, which points along the initial line of sight (due North) and closes the $10\ \mathrm{km}$ gap at $15\ \mathrm{km\,h^{-1}}$: time $= 10/15 = 2/3$ h. $\checkmark$
\end{exoresolu}

\begin{exoresolu}[Interception with two moving bodies choosing directions]
Ship $A$ is at $(0,0)$ with speed $20\ \mathrm{km\,h^{-1}}$. Ship $B$ is at $(100,0)$ with speed $15\ \mathrm{km\,h^{-1}}$. Both can choose their courses. Find the courses for $A$ to intercept $B$ in minimum time.
\tcblower
For minimum time interception, $A$ should head directly toward $B$'s initial position, and $B$ should flee directly away. But if $B$'s course is fixed (say due North at $15\ \mathrm{km\,h^{-1}}$), then:
$\vec{r}_A(0) = \vec{0}$, $\vec{r}_B(0) = 100\vec{i}$, $\vec{v}_B = 15\vec{j}$.
Let $\vec{v}_A = 20(\cos\theta\,\vec{i} + \sin\theta\,\vec{j})$.
Interception: $20\cos\theta \cdot t = 100$, $20\sin\theta \cdot t = 15t$.
From second: $\sin\theta = 15/20 = 0.75 \Rightarrow \theta \approx 48.6^\circ$.
Then $\cos\theta = \sqrt{1-0.75^2} = \sqrt{7}/4 \approx 0.661$.
$t = 100/(20\cos\theta) = 5/\cos\theta \approx 7.56\ \mathrm{h}$.
\end{exoresolu}

\courssub{Applications of Interception and Collision}

\begin{exoresolu}[Collision course of two aircraft]
Aircraft $P$ is at $(20\vec{i} + 30\vec{j})\ \mathrm{km}$ flying with velocity $(300\vec{i} + 400\vec{j})\ \mathrm{km\,h^{-1}}$. Aircraft $Q$ is at $(100\vec{i} + 10\vec{j})\ \mathrm{km}$ flying with velocity $(-200\vec{i} + 150\vec{j})\ \mathrm{km\,h^{-1}}$. Determine whether they will collide, and if so, when and where.
\tcblower
$\vec{r}_P(t) = (20+300t)\vec{i} + (30+400t)\vec{j}$.
$\vec{r}_Q(t) = (100-200t)\vec{i} + (10+150t)\vec{j}$.
Collision requires:
$20+300t = 100-200t \Rightarrow 500t = 80 \Rightarrow t = 0.16\ \mathrm{h} = 9.6\ \mathrm{min}$.
$30+400t = 10+150t \Rightarrow 250t = -20 \Rightarrow t = -0.08\ \mathrm{h}$.
Times differ $\Rightarrow$ \textbf{no collision}.
Relative velocity: $\vec{v}_{P/Q} = 500\vec{i} + 250\vec{j}$.
Initial relative position: $\vec{r}_{P/Q}(0) = -80\vec{i} + 20\vec{j}$.
Not parallel $\Rightarrow$ no collision.
\end{exoresolu}

\begin{exoresolu}[Interception with speed constraint]
A coastguard cutter can steam at $25\ \mathrm{km\,h^{-1}}$. At 10:00, a smuggler's boat is $30\ \mathrm{km}$ away on a bearing of $060^\circ$ and is moving at $15\ \mathrm{km\,h^{-1}}$ on a bearing of $150^\circ$. Find the course the cutter must steer to intercept the smuggler, and the time of interception.
\tcblower
$\vec{i}$ East, $\vec{j}$ North.
$\vec{r}_C(0) = \vec{0}$, $\vec{r}_S(0) = 30(\sin60^\circ\,\vec{i} + \cos60^\circ\,\vec{j}) = 15\sqrt{3}\,\vec{i} + 15\,\vec{j}$.
$\vec{v}_S = 15(\sin150^\circ\,\vec{i} + \cos150^\circ\,\vec{j}) = 7.5\,\vec{i} - 7.5\sqrt{3}\,\vec{j}$.
Let $\vec{v}_C = 25(\sin\theta\,\vec{i} + \cos\theta\,\vec{j})$.
Interception: $\vec{v}_{C/S} = \vec{v}_C - \vec{v}_S$ must be parallel to $\vec{r}_S(0)$.
$\vec{v}_{C/S} = (25\sin\theta - 7.5)\vec{i} + (25\cos\theta + 7.5\sqrt{3})\vec{j}$.
Parallel condition: $\dfrac{25\sin\theta - 7.5}{15\sqrt{3}} = \dfrac{25\cos\theta + 7.5\sqrt{3}}{15}$.
Solve: $25\sin\theta - 7.5 = \sqrt{3}(25\cos\theta + 7.5\sqrt{3}) = 25\sqrt{3}\cos\theta + 22.5$.
$25\sin\theta - 25\sqrt{3}\cos\theta = 30$.
$50\sin(\theta - 60^\circ) = 30 \Rightarrow \sin(\theta - 60^\circ) = 0.6$.
$\theta - 60^\circ = 36.9^\circ \Rightarrow \theta = 96.9^\circ$ (bearing).
$\vec{v}_{C/S}$ magnitude: $|\vec{v}_{C/S}| = \sqrt{(25\sin96.9^\circ - 7.5)^2 + (25\cos96.9^\circ + 7.5\sqrt{3})^2} \approx 40\ \mathrm{km\,h^{-1}}$.
Time $= 30/40 = 0.75\ \mathrm{h} = 45\ \mathrm{min}$.
Interception at 10:45.
\end{exoresolu}

\courssub{Relative Position and Relative Motion: Concept, Treatment and Applications}

\begin{definition}[Relative position and relative motion]
The \cle{relative position} of $A$ with respect to $B$ is
\[ \vec{r}_{A/B}(t) = \vec{r}_A(t) - \vec{r}_B(t). \]
If both move with constant velocities, then
\[ \vec{r}_{A/B}(t) = \vec{r}_{A/B}(0) + \vec{v}_{A/B}\, t, \]
where $\vec{v}_{A/B} = \vec{v}_A - \vec{v}_B$ is constant. In $B$'s frame, $B$ is stationary and $A$ moves in a straight line with velocity $\vec{v}_{A/B}$.
\end{definition}

\begin{methode}[Relative motion, distance apart, closest approach]
\begin{enumerate}
\item \textbf{Form the relative position:}
\[ \vec{r}_{A/B}(t) = \vec{r}_A(t) - \vec{r}_B(t) = \vec{r}_{A/B}(0) + \vec{v}_{A/B}\, t. \]
\item \textbf{Interpretation:} in $B$'s frame, $B$ is stationary and $A$ moves in a straight line with velocity $\vec{v}_{A/B}$.
\item \textbf{Distance apart:} $d(t) = |\vec{r}_{A/B}(t)|$. Work with $d^2(t)$ (a quadratic in $t$) to avoid square roots.
\item \textbf{Closest approach:} minimise $d^2$: either complete the square, or solve $\dfrac{d}{dt}\big(d^2\big) = 0$, giving
\[ t_{\min} = -\frac{\vec{r}_{A/B}(0)\cdot \vec{v}_{A/B}}{|\vec{v}_{A/B}|^2}. \]
\item \textbf{Collision test:} the bodies collide if and only if $d_{\min} = 0$, i.e. $\vec{v}_{A/B}$ is parallel to $\vec{r}_{A/B}(0)$ and directed to close the gap.
\end{enumerate}
\textbf{Reflex:} everything reduces to \emph{one} moving point (the relative position) and one fixed point (the origin).
\end{methode}

\begin{popfigure}
\begin{tikzpicture}[scale=0.7, >=stealth]
% Closest approach geometry
\draw[->, thick, popblue] (0,0) -- (4,1) node[midway, below right] {$\vec{r}_{A/B}(0)$};
\draw[->, thick, poporange] (4,1) -- (7,3) node[midway, above right] {$\vec{v}_{A/B}\,t$};
\draw[->, thick, popgreen] (0,0) -- (7,3) node[midway, above left] {$\vec{r}_{A/B}(t)$};
\draw[dashed] (0,0) -- (4,1) -- (7,3);
\draw[popred, thick] (0,0) -- (5.5, 1.375) node[midway, below left] {$d_{\min}$};
\draw (5.5,1.375) -- (7,3);
\node at (0,-0.3) {$B$ (origin in relative frame)};
\node at (4,0.7) {$A(0)$};
\node at (7,3.3) {$A(t)$};
\end{tikzpicture}
\legende{Closest approach: perpendicular from origin to relative path}
\end{popfigure}

\begin{exoresolu}[Closest approach of two vehicles]
At $8{:}00$ a.m., a bus $A$ is at the point with position vector $(2\vec{i} + 5\vec{j})\ \mathrm{km}$ and travels with constant velocity $(30\vec{i} + 40\vec{j})\ \mathrm{km\,h^{-1}}$. At the same time a taxi $B$ is at $(20\vec{i} + 3\vec{j})\ \mathrm{km}$ travelling with velocity $(10\vec{i} + 24\vec{j})\ \mathrm{km\,h^{-1}}$.
\begin{itemize}
\item[(a)] Find the velocity of $A$ relative to $B$.
\item[(b)] Find the time at which the two vehicles are closest together, and their least distance apart.
\item[(c)] Would they collide if the taxi instead travelled at $(14\vec{i} + 28\vec{j})\ \mathrm{km\,h^{-1}}$?
\end{itemize}
\tcblower
\textbf{(a)}
\[ \vec{v}_{A/B} = \vec{v}_A - \vec{v}_B = (30-10)\vec{i} + (40-24)\vec{j} = 20\vec{i} + 16\vec{j}\ \mathrm{km\,h^{-1}}. \]

\textbf{(b)} Initial relative position:
\[ \vec{r}_{A/B}(0) = (2-20)\vec{i} + (5-3)\vec{j} = -18\vec{i} + 2\vec{j}. \]
Hence
\[ \vec{r}_{A/B}(t) = (-18 + 20t)\vec{i} + (2 + 16t)\vec{j}. \]
Squared distance:
\[ d^2(t) = (20t - 18)^2 + (16t + 2)^2 = 400t^2 - 720t + 324 + 256t^2 + 64t + 4 = 656t^2 - 656t + 328. \]
This quadratic is minimised where its derivative vanishes:
\[ \frac{d}{dt}\big(d^2\big) = 1312t - 656 = 0 \quad\Longrightarrow\quad t = \frac{656}{1312} = 0.5\ \mathrm{h}. \]
So the vehicles are closest at $8{:}30$ a.m. The least distance:
\[ d^2(0.5) = 656(0.25) - 656(0.5) + 328 = 164 - 328 + 328 = 164, \]
\[ d_{\min} = \sqrt{164} = 2\sqrt{41} \approx 12.8\ \mathrm{km}. \]

\textbf{(c)} With $\vec{v}_B = 14\vec{i} + 28\vec{j}$:
\[ \vec{v}_{A/B} = 16\vec{i} + 12\vec{j}, \qquad \vec{r}_{A/B}(t) = (-18 + 16t)\vec{i} + (2 + 12t)\vec{j}. \]
Collision requires both components zero simultaneously:
\[ -18 + 16t = 0 \Rightarrow t = \tfrac{9}{8}; \qquad 2 + 12t = 0 \Rightarrow t = -\tfrac{1}{6}. \]
The times differ, so \textbf{no collision}. (Equivalently, $\vec{v}_{A/B} = 16\vec{i}+12\vec{j}$ is not parallel to $\vec{r}_{A/B}(0) = -18\vec{i} + 2\vec{j}$ since $\frac{16}{-18} \neq \frac{12}{2}$.)

\textbf{Key point:} collision needs \emph{both} coordinates to coincide at the \emph{same} instant — checking only one component is a classic error.
\end{exoresolu}

\begin{exoresolu}[Relative motion with variable velocity]
A particle $A$ has position $\vec{r}_A = (t^2\vec{i} + 2t\vec{j})\ \mathrm{m}$. Particle $B$ has position $\vec{r}_B = (4t\vec{i} + t^2\vec{j})\ \mathrm{m}$. Find the time when they are closest and the minimum distance.
\tcblower
$\vec{r}_{A/B} = (t^2 - 4t)\vec{i} + (2t - t^2)\vec{j}$.
$d^2 = (t^2-4t)^2 + (2t-t^2)^2 = t^4 - 8t^3 + 16t^2 + 4t^2 - 4t^3 + t^4 = 2t^4 - 12t^3 + 20t^2$.
$\frac{d}{dt}(d^2) = 8t^3 - 36t^2 + 40t = 4t(2t^2 - 9t + 10) = 4t(2t-5)(t-2)$.
Critical points: $t=0$, $t=2$, $t=2.5$.
$d^2(0)=0$ (they start together!), $d^2(2)=2(16)-12(8)+20(4)=32-96+80=16$, $d^2(2.5)=2(39.0625)-12(15.625)+20(6.25)=78.125-187.5+125=15.625$.
Minimum distance is $0$ at $t=0$ (they start at same point). Next closest is at $t=2.5$ with $d=\sqrt{15.625}\approx 3.95\ \mathrm{m}$.
\end{exoresolu}

\courssub{Bearing Problems with Wind or Current}

\begin{methode}[Aircraft/ship problems on a required bearing]
\begin{enumerate}
\item \textbf{Translate bearings into components} immediately: bearing $\beta$ means angle $\beta$ clockwise from North, so $\vec{v} = v\sin\beta\,\vec{i} + v\cos\beta\,\vec{j}$ ($\vec{i}$ East, $\vec{j}$ North).
\item \textbf{Write} $\vec{v}_{\text{ground}} = \vec{v}_{\text{plane/air}} + \vec{v}_{\text{wind}}$.
\item \textbf{Impose the required track} (e.g. due North means the East component of $\vec{v}_{\text{ground}}$ is zero).
\item \textbf{Solve for the heading} using the sine rule in the velocity triangle, then compute the ground speed.
\item \textbf{For a round trip,} compute each leg separately and add the times.
\end{enumerate}
\end{methode}

\begin{popfigure}
\begin{tikzpicture}[scale=0.7, >=stealth]
% Bearing convention
\draw[->, thick] (0,0) -- (0,3) node[above] {North};
\draw[->, thick] (0,0) -- (3,0) node[right] {East};
\draw[->, thick, poporange] (0,0) -- (2,2) node[midway, above right] {$\vec{v}$};
\draw[dashed] (2,0) -- (2,2) -- (0,2);
\node at (1,0.3) {$\beta$};
\node at (0.3,1) {$90^\circ-\beta$};
\node at (2.3,2.3) {Bearing $\beta$};
\end{tikzpicture}
\legende{Bearing $\beta$: $\vec{v} = v\sin\beta\,\vec{i} + v\cos\beta\,\vec{j}$}
\end{popfigure}

\begin{exoresolu}[Aircraft flying in a cross-wind]
A small aircraft flying from Douala has a speed of $250\ \mathrm{km\,h^{-1}}$ in still air. A wind blows from the West at $50\ \mathrm{km\,h^{-1}}$ (i.e. towards the East). The pilot wishes to travel due North. Find the course she must steer, her ground speed, and the time for a $400\ \mathrm{km}$ journey due North.
\tcblower
Take $\vec{i}$ East, $\vec{j}$ North. Wind: $\vec{v}_W = 50\vec{i}$. Let the pilot steer at angle $\theta$ West of North:
\[ \vec{v}_{P} = -250\sin\theta\,\vec{i} + 250\cos\theta\,\vec{j}. \]
Ground velocity:
\[ \vec{v}_G = (50 - 250\sin\theta)\vec{i} + 250\cos\theta\,\vec{j}. \]
For a due-North track, the East component must vanish:
\[ 50 - 250\sin\theta = 0 \quad\Longrightarrow\quad \sin\theta = \frac{50}{250} = 0.2 \quad\Longrightarrow\quad \theta \approx 11.5^\circ. \]
She must steer \textbf{N$11.5^\circ$W} (bearing $348.5^\circ$), pointing slightly into the wind.

\textbf{Ground speed:}
\[ v_G = 250\cos\theta = 250\sqrt{1 - 0.04} = 250\sqrt{0.96} \approx 244.9\ \mathrm{km\,h^{-1}}. \]

\textbf{Time for $400\ \mathrm{km}$ due North:}
\[ t = \frac{400}{244.9} \approx 1.633\ \mathrm{h} \approx 1\ \mathrm{h}\ 38\ \mathrm{min}. \]

\textbf{Remark:} the wind reduces the useful (northward) speed even though it blows perpendicular to the track, because part of the aircraft's velocity must be sacrificed to cancel the drift.
\end{exoresolu}

\begin{exoresolu}[Round trip with wind]
An aircraft flies from $A$ to $B$, $500\ \mathrm{km}$ due East, and returns. Its airspeed is $300\ \mathrm{km\,h^{-1}}$. A wind blows at $50\ \mathrm{km\,h^{-1}}$ from the South-West (i.e. towards North-East). Find the total time for the round trip.
\tcblower
Wind: $\vec{v}_W = 50(\sin45^\circ\,\vec{i} + \cos45^\circ\,\vec{j}) = 25\sqrt{2}(\vec{i}+\vec{j}) \approx 35.36(\vec{i}+\vec{j})$.
\textbf{Outward (East):} Need $\vec{v}_G = v_G\vec{i}$. $\vec{v}_P = \vec{v}_G - \vec{v}_W = (v_G - 35.36)\vec{i} - 35.36\vec{j}$.
$|\vec{v}_P| = 300 \Rightarrow (v_G - 35.36)^2 + 35.36^2 = 90000$.
$v_G^2 - 70.72v_G + 2500 + 1250 = 90000 \Rightarrow v_G^2 - 70.72v_G - 86250 = 0$.
$v_G = \frac{70.72 + \sqrt{5000 + 345000}}{2} \approx \frac{70.72 + 591.6}{2} \approx 331.2\ \mathrm{km\,h^{-1}}$.
Time out $= 500/331.2 \approx 1.51\ \mathrm{h}$.

\textbf{Return (West):} Need $\vec{v}_G = -v_G'\vec{i}$. $\vec{v}_P = (-v_G' - 35.36)\vec{i} - 35.36\vec{j}$.
$(-v_G' - 35.36)^2 + 35.36^2 = 90000 \Rightarrow (v_G' + 35.36)^2 = 86250$.
$v_G' = \sqrt{86250} - 35.36 \approx 293.7 - 35.36 = 258.3\ \mathrm{km\,h^{-1}}$.
Time back $= 500/258.3 \approx 1.936\ \mathrm{h}$.
Total time $\approx 3.45\ \mathrm{h} = 3\ \mathrm{h}\ 27\ \mathrm{min}$.
\end{exoresolu}

\begin{aretenir}[Key formulas for GCE Advanced Level]
\begin{itemize}
\item Relative velocity: $\vec{v}_{A/B} = \vec{v}_A - \vec{v}_B$.
\item Chain rule: $\vec{v}_{A/C} = \vec{v}_{A/B} + \vec{v}_{B/C}$.
\item True velocity: $\vec{v}_{\text{true}} = \vec{v}_{\text{own}} + \vec{v}_{\text{medium}}$.
\item Interception: $\vec{r}_A(t) = \vec{r}_B(t)$ at same $t$.
\item Relative position: $\vec{r}_{A/B}(t) = \vec{r}_{A/B}(0) + \vec{v}_{A/B}t$.
\item Distance: $d(t) = |\vec{r}_{A/B}(t)|$; minimise $d^2(t)$.
\item Closest approach time: $t_{\min} = -\dfrac{\vec{r}_{A/B}(0)\cdot\vec{v}_{A/B}}{|\vec{v}_{A/B}|^2}$.
\item Collision $\Leftrightarrow$ $\vec{v}_{A/B}$ parallel to $\vec{r}_{A/B}(0)$ and $t_{\min} > 0$.
\item Bearing $\beta$: $\vec{v} = v\sin\beta\,\vec{i} + v\cos\beta\,\vec{j}$.
\end{itemize}
\end{aretenir}

\begin{attention}[Common mistakes in this chapter]
\begin{itemize}
\item Adding the \emph{magnitudes} of non-parallel velocities instead of adding the vectors componentwise.
\item Writing $\vec{v}_{A/B} = \vec{v}_B - \vec{v}_A$ (wrong order): remember ``$A$ relative to $B$'' $=$ velocity of $A$ \emph{minus} velocity of $B$.
\item Forgetting that interception requires \emph{both} position components to be equal at the \emph{same} time $t$.
\item Using the heading direction instead of the resultant (true) velocity to compute crossing times or distances travelled over ground.
\item Mixing units (km with $\mathrm{m\,s^{-1}}$, hours with minutes): convert everything before computing.
\item Giving a direction as a bare angle without saying from which axis it is measured — always state a bearing or a reference direction.
\item In closest approach problems, forgetting to check that $t_{\min} \ge 0$ (if $t_{\min} < 0$, the minimum occurs at $t=0$).
\item Confusing ``wind from the West'' (blows towards East) with ``wind to the West''.
\end{itemize}
\end{attention}

\begin{exercice}[Practice problems]
\begin{enumerate}
\item A particle has velocity $\vec{v} = (6t\,\vec{i} + (4-2t)\,\vec{j})\ \mathrm{m\,s^{-1}}$. Find its acceleration and show that its direction is constant. Find the speed at $t=3$.
\item A man walks at $5\ \mathrm{km\,h^{-1}}$ due North. The wind appears to come from the North-West at $10\ \mathrm{km\,h^{-1}}$. Find the true wind velocity.
\item A boat crosses a $200\ \mathrm{m}$ wide river flowing at $2\ \mathrm{m\,s^{-1}}$. The boat's speed in still water is $4\ \mathrm{m\,s^{-1}}$. (a) If it heads straight across, find the resultant velocity and landing point. (b) Find the heading for a straight crossing and the time taken.
\item Two ships $A$ and $B$ have position vectors $\vec{r}_A = (2+3t)\vec{i} + (5+4t)\vec{j}$ and $\vec{r}_B = (10+t)\vec{i} + (2+2t)\vec{j}$ (km, hours). Find their closest distance of approach and the time it occurs.
\item An aircraft flies at $400\ \mathrm{km\,h^{-1}}$ in still air. A wind of $80\ \mathrm{km\,h^{-1}}$ blows from bearing $210^\circ$. The pilot wishes to fly on a track of $060^\circ$. Find the heading and ground speed.
\end{enumerate}
\end{exercice}

\begin{corrige}[Exercise 1]
$\vec{a} = 6\vec{i} - 2\vec{j}$ (constant). Direction: $\tan\theta = -2/6 = -1/3$, constant. At $t=3$, $\vec{v} = 18\vec{i} - 2\vec{j}$, speed $= \sqrt{324+4} = \sqrt{328} \approx 18.1\ \mathrm{m\,s^{-1}}$.
\end{corrige}

\begin{corrige}[Exercise 2]
$\vec{v}_M = 5\vec{j}$. Apparent wind $\vec{v}_{W/M} = \vec{v}_W - \vec{v}_M$ from NW $\Rightarrow$ towards SE: $\vec{v}_{W/M} = k(\vec{i} - \vec{j})$, $k>0$. $|\vec{v}_{W/M}| = 10 \Rightarrow k\sqrt{2}=10 \Rightarrow k=5\sqrt{2}$.
$\vec{v}_W = \vec{v}_{W/M} + \vec{v}_M = 5\sqrt{2}\vec{i} + (5-5\sqrt{2})\vec{j} \approx 7.07\vec{i} - 2.07\vec{j}$.
Magnitude $\approx 7.37\ \mathrm{km\,h^{-1}}$, direction $\tan^{-1}(-2.07/7.07) \approx -16.3^\circ$ (i.e. $16.3^\circ$ South of East).
\end{corrige}

\begin{corrige}[Exercise 3]
(a) $\vec{v}_{\text{own}} = 4\vec{j}$, $\vec{v}_{\text{curr}} = 2\vec{i}$. $\vec{v}_{\text{true}} = 2\vec{i} + 4\vec{j}$, $|\vec{v}| = \sqrt{20} \approx 4.47\ \mathrm{m\,s^{-1}}$, angle $\tan^{-1}(2/4)=26.6^\circ$ East of North. Time $= 200/4 = 50\ \mathrm{s}$. Drift $= 2\times50 = 100\ \mathrm{m}$.
(b) Need $4\sin\alpha = 2 \Rightarrow \sin\alpha = 0.5 \Rightarrow \alpha = 30^\circ$ upstream. Across component $= 4\cos30^\circ = 2\sqrt{3} \approx 3.46\ \mathrm{m\,s^{-1}}$. Time $= 200/(2\sqrt{3}) \approx 57.7\ \mathrm{s}$.
\end{corrige}

\begin{corrige}[Exercise 4]
$\vec{r}_{A/B} = (-8+2t)\vec{i} + (3+2t)\vec{j}$. $d^2 = (2t-8)^2 + (2t+3)^2 = 8t^2 -20t +73$. Min at $t = 20/16 = 1.25\ \mathrm{h}$. $d_{\min}^2 = 8(1.5625) - 25 + 73 = 12.5 - 25 + 73 = 60.5$. $d_{\min} = \sqrt{60.5} \approx 7.78\ \mathrm{km}$.
\end{corrige}

\begin{corrige}[Exercise 5]
Wind from $210^\circ$ $\Rightarrow$ towards $030^\circ$: $\vec{v}_W = 80(\sin30^\circ\,\vec{i} + \cos30^\circ\,\vec{j}) = 40\vec{i} + 40\sqrt{3}\vec{j}$.
Track $060^\circ$: $\vec{v}_G = v_G(\sin60^\circ\,\vec{i} + \cos60^\circ\,\vec{j}) = v_G(\frac{\sqrt{3}}{2}\vec{i} + \frac{1}{2}\vec{j})$.
$\vec{v}_P = \vec{v}_G - \vec{v}_W = (\frac{\sqrt{3}}{2}v_G - 40)\vec{i} + (\frac{1}{2}v_G - 40\sqrt{3})\vec{j}$.
$|\vec{v}_P| = 400 \Rightarrow (\frac{\sqrt{3}}{2}v_G - 40)^2 + (\frac{1}{2}v_G - 40\sqrt{3})^2 = 160000$.
$\frac{3}{4}v_G^2 - 40\sqrt{3}v_G + 1600 + \frac{1}{4}v_G^2 - 40\sqrt{3}v_G + 4800 = 160000$.
$v_G^2 - 80\sqrt{3}v_G + 6400 = 160000 \Rightarrow v_G^2 - 138.56v_G - 153600 = 0$.
$v_G = \frac{138.56 + \sqrt{19200 + 614400}}{2} = \frac{138.56 + 795.2}{2} \approx 466.9\ \mathrm{km\,h^{-1}}$.
Heading: $\vec{v}_P = (405.1 - 40)\vec{i} + (233.4 - 69.3)\vec{j} = 365.1\vec{i} + 164.1\vec{j}$.
Bearing $= \tan^{-1}(365.1/164.1) \approx 65.8^\circ$.
\end{corrige}

\newpage

\coursec{Exercises}

\courssub{I check my knowledge}

\begin{exercice}[Multiple choice: relative velocity]
A car $A$ moves due east at $20\ \mathrm{m\,s^{-1}}$ and a car $B$ moves due north at $15\ \mathrm{m\,s^{-1}}$. The velocity of $B$ relative to $A$, denoted ${}_B\vec{v}_A$, is:
\begin{enumerate}
\item $5\ \mathrm{m\,s^{-1}}$ due north;
\item $25\ \mathrm{m\,s^{-1}}$ on a bearing of $037^{\circ}$;
\item $25\ \mathrm{m\,s^{-1}}$ on a bearing of $323^{\circ}$;
\item $35\ \mathrm{m\,s^{-1}}$ due north-east.
\end{enumerate}
Choose the correct answer, justifying your choice with a vector diagram.
\end{exercice}

\begin{exercice}[True or false]
For each of the following statements, say whether it is \textbf{true} or \textbf{false} and justify your answer.
\begin{enumerate}
\item If $\vec{v}_A$ and $\vec{v}_B$ are the velocities of two bodies relative to the ground, then the velocity of $A$ relative to $B$ is ${}_A\vec{v}_B = \vec{v}_A - \vec{v}_B$.
\item A boat steering perpendicular to the banks of a river crosses the river in the shortest possible time.
\item If two particles moving with constant velocities collide, then their relative velocity is zero.
\item The acceleration of a particle is the vector $\vec{a} = \dfrac{d\vec{v}}{dt}$, and it is always parallel to the velocity $\vec{v}$.
\end{enumerate}
\end{exercice}

\begin{exercice}[Definitions]
Define each of the following terms, illustrating each definition with a short example:
\begin{enumerate}
\item the \cle{relative velocity} of a body $A$ with respect to a body $B$;
\item the \cle{true velocity} of a boat in a river problem;
\item the \cle{relative position} of a particle $P$ with respect to a particle $Q$;
\item the \cle{closest distance of approach} of two moving particles.
\end{enumerate}
\end{exercice}

\begin{exercice}[Direct calculations]
A particle moves with velocity $\vec{v} = (3\vec{i} + 4\vec{j})\ \mathrm{m\,s^{-1}}$ relative to a ship moving with velocity $(\vec{i} - 2\vec{j})\ \mathrm{m\,s^{-1}}$ relative to the water. The water itself flows with velocity $(2\vec{i} + \vec{j})\ \mathrm{m\,s^{-1}}$ relative to the ground.
\begin{enumerate}
\item Write down the velocity of the ship relative to the ground.
\item Find the true velocity of the particle relative to the ground, giving its magnitude and direction.
\item Given that at time $t = 0$ the particle has position vector $(5\vec{i} - \vec{j})\ \mathrm{m}$ relative to a fixed point on the shore, find its position vector at time $t = 4\ \mathrm{s}$.
\end{enumerate}
\end{exercice}

\courssub{I practise}

\begin{exercice}[Crossing the river]
A boat that can travel at $8\ \mathrm{km\,h^{-1}}$ in still water crosses a river $400\ \mathrm{m}$ wide whose current flows at $3\ \mathrm{km\,h^{-1}}$.
\begin{enumerate}
\item Find the quickest crossing time and the distance drifted downstream during this crossing.
\item Find the course the boat must steer in order to land directly opposite its starting point, and the corresponding crossing time.
\item Comment on which crossing is shorter in time and why.
\end{enumerate}
\end{exercice}

\begin{exercice}[Velocity of one car relative to another]
Car $A$ travels due east at $20\ \mathrm{m\,s^{-1}}$ and car $B$ travels on a bearing of $030^{\circ}$ at $25\ \mathrm{m\,s^{-1}}$.
\begin{enumerate}
\item Express the velocities $\vec{v}_A$ and $\vec{v}_B$ in terms of the unit vectors $\vec{i}$ (east) and $\vec{j}$ (north).
\item Find the velocity of $B$ relative to $A$, giving its magnitude and bearing.
\item Find the velocity of $A$ relative to $B$.
\end{enumerate}
\end{exercice}

\begin{exercice}[The apparent wind]
To a cyclist riding due north at $6\ \mathrm{m\,s^{-1}}$, the wind appears to come from the west. When she doubles her speed, the wind appears to come from the north-west. Find the true velocity of the wind, giving its magnitude and the direction from which it blows.
\end{exercice}

\begin{exercice}[Ferry on the Sanaga]
A ferry crosses the Sanaga river steering perpendicular to the banks. Passengers observe that it lands $150\ \mathrm{m}$ downstream after a crossing lasting $5$ minutes on a river $600\ \mathrm{m}$ wide.
\begin{enumerate}
\item Find the speed of the current.
\item Find the ferry's speed in still water.
\item Find the ferry's true speed relative to the ground and the angle its actual path makes with the line perpendicular to the banks.
\end{enumerate}
\end{exercice}

\courssub{I prepare for the GCE}

\begin{exercice}[Interception at sea \hfill (12 marks)]
Ship $P$ leaves port $O$ at noon and sails at $15\ \mathrm{km\,h^{-1}}$ on a bearing of $060^{\circ}$. At noon, ship $Q$ is $20\ \mathrm{km}$ due east of $O$ and sails at $18\ \mathrm{km\,h^{-1}}$ with constant velocity.
\begin{enumerate}
\item Taking $\vec{i}$ and $\vec{j}$ as unit vectors due east and due north respectively, write down the position vectors of $P$ and $Q$ at time $t$ hours after noon. \hfill (3 marks)
\item Show that if $Q$ steers a course with velocity $\vec{v}_Q$, interception requires $\vec{v}_Q = \vec{v}_P + \dfrac{20\vec{i}}{T}$, where $T$ is the interception time. \hfill (3 marks)
\item Hence find the course $Q$ must steer (as a bearing) to intercept $P$, and the time of interception. \hfill (4 marks)
\item Find the distance from $O$ of the point of interception. \hfill (2 marks)
\end{enumerate}
\end{exercice}

\begin{exercice}[Collision or closest approach \hfill (14 marks)]
Two aircraft $A$ and $B$ fly at the same altitude with constant velocities. At time $t = 0$ (in hours), their position vectors relative to a control tower are
\[
\vec{r}_A = (10\vec{i} + 5\vec{j})\ \mathrm{km}, \qquad \vec{r}_B = (-4\vec{i} + 20\vec{j})\ \mathrm{km},
\]
and their velocities are
\[
\vec{v}_A = (300\vec{i} + 100\vec{j})\ \mathrm{km\,h^{-1}}, \qquad \vec{v}_B = (200\vec{i} - 100\vec{j})\ \mathrm{km\,h^{-1}}.
\]
\begin{enumerate}
\item Write down the position vectors of $A$ and $B$ at time $t$. \hfill (2 marks)
\item Find the position vector of $B$ relative to $A$ and the velocity of $B$ relative to $A$. \hfill (3 marks)
\item Determine whether the two aircraft collide. \hfill (3 marks)
\item If they do not collide, find the time at which they are closest together and their closest distance of approach. \hfill (4 marks)
\item Air-traffic regulations require a minimum separation of $5\ \mathrm{km}$. State, with justification, whether the regulations are violated. \hfill (2 marks)
\end{enumerate}
\end{exercice}

\begin{exercice}[Radio contact between two boats \hfill (13 marks)]
Two boats $M$ and $N$ move with constant velocities on a calm sea. Relative to a fixed buoy, their position vectors at time $t$ hours are
\[
\vec{r}_M = (2 + 6t)\vec{i} + (1 + 8t)\vec{j}, \qquad \vec{r}_N = (10 - 2t)\vec{i} + (3 + 4t)\vec{j},
\]
all distances being in kilometres.
\begin{enumerate}
\item Find the velocity of $M$ relative to $N$ and interpret it physically. \hfill (3 marks)
\item Show that the square of the distance between the boats at time $t$ is $MN^2 = 80t^2 - 136t + 68$. \hfill (4 marks)
\item The boats are within radio range when their separation is at most $5\ \mathrm{km}$. Find the time interval during which they can communicate, giving the endpoints in hours and minutes. \hfill (4 marks)
\item Find the minimum distance between the boats and the time at which it occurs. \hfill (2 marks)
\end{enumerate}
\end{exercice}

\newpage

\coursec{Solutions to the exercises}

\begin{corrige}[Exercise 1 --- Multiple choice: relative velocity]
\textbf{Correct answer: (c)} --- the only option with magnitude $25\ \mathrm{m\,s^{-1}}$ in the north-west quadrant (see the remark below on the printed bearing).

Take $\vec{i}$ due east and $\vec{j}$ due north. Then
\[
\vec{v}_A = 20\vec{i}, \qquad \vec{v}_B = 15\vec{j}.
\]
The velocity of $B$ relative to $A$ is
\[
{}_B\vec{v}_A = \vec{v}_B - \vec{v}_A = -20\vec{i} + 15\vec{j}.
\]
Its magnitude is
\[
|{}_B\vec{v}_A| = \sqrt{(-20)^2 + 15^2} = \sqrt{625} = 25\ \mathrm{m\,s^{-1}},
\]
which eliminates (a) and (d). The vector points \emph{west of north}, so the bearing lies between $270^{\circ}$ and $360^{\circ}$: option (b), bearing $037^{\circ}$ (north-east), is impossible. Hence \textbf{(c)}.

\begin{popfigure}
\begin{center}
\begin{tikzpicture}[scale=0.09]
\draw[->,gray] (-24,0) -- (6,0) node[below left] {\scriptsize east};
\draw[->,gray] (0,-4) -- (0,20) node[below left] {\scriptsize north};
\draw[->,very thick,popblue] (0,0) -- (20,0) node[midway,below] {\scriptsize $\vec{v}_A$};
\draw[->,very thick,popgreen] (0,0) -- (0,15) node[midway,right] {\scriptsize $\vec{v}_B$};
\draw[->,very thick,poporange] (20,0) -- (20,15) node[midway,right] {\scriptsize $\vec{v}_B$};
\draw[->,very thick,popink] (0,0) -- (-20,15) node[midway,above left] {\scriptsize ${}_B\vec{v}_A$};
\draw[->,thick,popmuted,dashed] (0,0) -- (-20,0) node[midway,below] {\scriptsize $-\vec{v}_A$};
\draw[dashed,popmuted] (-20,0) -- (-20,15) -- (0,15);
\end{tikzpicture}
\end{center}
\legende{Vector triangle: ${}_B\vec{v}_A = \vec{v}_B + (-\vec{v}_A)$.}
\end{popfigure}

\textbf{Remark for the examiner.} The exact bearing is $360^{\circ} - \arctan\!\left(\dfrac{20}{15}\right) = 360^{\circ} - 53.1^{\circ} = 306.9^{\circ}$. The value $323^{\circ}$ printed in option (c) corresponds to components $(-15, 20)$; it is a misprint for $307^{\circ}$. Option (c) remains the intended and only defensible answer.
\end{corrige}

\begin{corrige}[Exercise 2 --- True or false]
\begin{enumerate}
\item \textbf{True.} By definition, the velocity of $A$ relative to $B$ is obtained by subtracting the velocity of the observer $B$: ${}_A\vec{v}_B = \vec{v}_A - \vec{v}_B$. This is the fundamental law of composition of velocities (Galilean addition), valid because all velocities here are measured relative to the same ground frame.
\item \textbf{True.} If the boat steers perpendicular to the banks with speed $u$ in still water, the crossing time is $t = \dfrac{d}{u}$, where $d$ is the width of the river: the whole of the boat's own speed is used across the river. Any other heading reduces the component of velocity perpendicular to the banks and therefore increases the crossing time. (The drift downstream is greatest in this case, but the question concerns time only.)
\item \textbf{False.} A collision occurs when the \emph{relative position} becomes zero, i.e.\ ${}_A\vec{r}_B = \vec{0}$ at some instant. The relative velocity is precisely what brings the particles together; if ${}_A\vec{v}_B = \vec{0}$ with the particles initially apart, they would \emph{never} collide. Counter-example: two cars approaching a crossroads head-on at $60\ \mathrm{km\,h^{-1}}$ collide with relative velocity $120\ \mathrm{km\,h^{-1}} \neq 0$.
\item \textbf{False.} It is true that $\vec{a} = \dfrac{d\vec{v}}{dt}$, but $\vec{a}$ need not be parallel to $\vec{v}$. Acceleration parallel to $\vec{v}$ changes only the \emph{magnitude} of the velocity; a component of $\vec{a}$ perpendicular to $\vec{v}$ changes its \emph{direction}. In uniform circular motion, $\vec{a}$ is perpendicular to $\vec{v}$ at every instant (centripetal acceleration), while the speed stays constant.
\end{enumerate}
\end{corrige}

\begin{corrige}[Exercise 3 --- Definitions]
\begin{enumerate}
\item \textbf{Relative velocity.} The velocity of a body $A$ with respect to a body $B$ is the velocity that an observer moving with $B$ would attribute to $A$; mathematically ${}_A\vec{v}_B = \vec{v}_A - \vec{v}_B$, where $\vec{v}_A, \vec{v}_B$ are the velocities relative to a common frame (the ground). \emph{Example:} two taxis on Boul\'evard de la Libert\'e both travel east at $50$ and $70\ \mathrm{km\,h^{-1}}$; the faster taxi's velocity relative to the slower one is $20\ \mathrm{km\,h^{-1}}$ east.
\item \textbf{True velocity of a boat.} The true velocity is the velocity of the boat \emph{relative to the ground (the banks)}; it is the vector sum of the boat's velocity relative to the water and the velocity of the current: $\vec{v}_{\text{true}} = \vec{v}_{\text{boat/water}} + \vec{v}_{\text{water/ground}}$. \emph{Example:} a pirogue paddling at $4\ \mathrm{m\,s^{-1}}$ across the Wouri with a $1\ \mathrm{m\,s^{-1}}$ current has true velocity $\sqrt{4^2+1^2} = \sqrt{17}\ \mathrm{m\,s^{-1}}$ along a slanting path.
\item \textbf{Relative position.} The position of $P$ with respect to $Q$ is the vector ${}_P\vec{r}_Q = \vec{r}_P - \vec{r}_Q$, i.e.\ the displacement that takes one from $Q$ to $P$. \emph{Example:} if $P$ is at $(3\vec{i} + 2\vec{j})\ \mathrm{km}$ and $Q$ at $(\vec{i} - \vec{j})\ \mathrm{km}$ from a beacon, then ${}_P\vec{r}_Q = (2\vec{i} + 3\vec{j})\ \mathrm{km}$.
\item \textbf{Closest distance of approach.} For two particles moving with constant velocities, it is the minimum value of the separation $d(t) = |{}_P\vec{r}_Q(t)|$ over $t \geq 0$; it is found by minimising $d^2(t)$ (a quadratic in $t$). \emph{Example:} two aircraft on non-colliding courses pass each other; their closest distance of approach is the smallest separation recorded on the radar screen.
\end{enumerate}
\end{corrige}

\begin{corrige}[Exercise 4 --- Direct calculations]
\begin{enumerate}
\item The ship's velocity relative to the ground is the sum of its velocity relative to the water and the water's velocity relative to the ground:
\[
\vec{v}_{S} = (\vec{i} - 2\vec{j}) + (2\vec{i} + \vec{j}) = (3\vec{i} - \vec{j})\ \mathrm{m\,s^{-1}}.
\]
\item The true velocity of the particle relative to the ground is
\[
\vec{v} = \underbrace{(3\vec{i} + 4\vec{j})}_{\text{particle/ship}} + \underbrace{(3\vec{i} - \vec{j})}_{\text{ship/ground}} = (6\vec{i} + 3\vec{j})\ \mathrm{m\,s^{-1}}.
\]
Magnitude: $|\vec{v}| = \sqrt{6^2 + 3^2} = \sqrt{45} = 3\sqrt{5} \approx 6.71\ \mathrm{m\,s^{-1}}$.\\
Direction: $\tan\theta = \dfrac{3}{6} = \dfrac{1}{2}$, so $\theta \approx 26.6^{\circ}$ north of east, i.e.\ on a bearing of $90^{\circ} - 26.6^{\circ} = 063.4^{\circ}$.
\item Since the velocity is constant, $\vec{r}(t) = \vec{r}(0) + t\vec{v}$:
\[
\vec{r}(4) = (5\vec{i} - \vec{j}) + 4(6\vec{i} + 3\vec{j}) = (29\vec{i} + 11\vec{j})\ \mathrm{m}.
\]
\end{enumerate}
\end{corrige}

\begin{corrige}[Exercise 5 --- Crossing the river]
Convert the width: $400\ \mathrm{m} = 0.4\ \mathrm{km}$.
\begin{enumerate}
\item \textbf{Quickest crossing.} Steer perpendicular to the banks. The whole speed $8\ \mathrm{km\,h^{-1}}$ is then used across the river:
\[
t_{\min} = \frac{0.4}{8} = 0.05\ \mathrm{h} = 3\ \mathrm{min}.
\]
During this time the current carries the boat downstream by
\[
d = 3 \times 0.05 = 0.15\ \mathrm{km} = 150\ \mathrm{m}.
\]
\item \textbf{Landing directly opposite.} The boat must steer upstream at an angle $\theta$ to the perpendicular so that the upstream component of its velocity cancels the current:
\[
8\sin\theta = 3 \quad\Longrightarrow\quad \sin\theta = \frac{3}{8} \quad\Longrightarrow\quad \theta \approx 22.0^{\circ}.
\]
The resultant (true) speed across the river is then
\[
v = \sqrt{8^2 - 3^2} = \sqrt{55} \approx 7.42\ \mathrm{km\,h^{-1}},
\]
and the crossing time is
\[
t = \frac{0.4}{\sqrt{55}}\ \mathrm{h} \approx 0.0539\ \mathrm{h} \approx 3\ \mathrm{min}\ 14\ \mathrm{s}.
\]
\item \textbf{Comparison.} The perpendicular heading is quicker ($3$ min against $3$ min $14$ s). Reason: when steering straight across, the \emph{entire} speed of the boat contributes to the crossing; when steering upstream, part of the speed is spent fighting the current, leaving only $\sqrt{55} < 8\ \mathrm{km\,h^{-1}}$ across the river. The price of the quicker crossing is the $150\ \mathrm{m}$ downstream drift.
\end{enumerate}
\end{corrige}

\begin{corrige}[Exercise 6 --- Velocity of one car relative to another]
\begin{enumerate}
\item With $\vec{i}$ east and $\vec{j}$ north:
\[
\vec{v}_A = 20\vec{i}\ \mathrm{m\,s^{-1}}, \qquad
\vec{v}_B = 25\sin 30^{\circ}\,\vec{i} + 25\cos 30^{\circ}\,\vec{j} = (12.5\vec{i} + 21.65\vec{j})\ \mathrm{m\,s^{-1}}.
\]
\item Velocity of $B$ relative to $A$:
\[
{}_B\vec{v}_A = \vec{v}_B - \vec{v}_A = (12.5 - 20)\vec{i} + 21.65\vec{j} = (-7.5\vec{i} + 21.65\vec{j})\ \mathrm{m\,s^{-1}}.
\]
Magnitude:
\[
|{}_B\vec{v}_A| = \sqrt{7.5^2 + 21.65^2} = \sqrt{56.25 + 468.75} = \sqrt{525} \approx 22.9\ \mathrm{m\,s^{-1}}.
\]
The vector points north-west; the angle west of north is
\[
\arctan\!\left(\frac{7.5}{21.65}\right) \approx 19.1^{\circ},
\]
so the bearing is $360^{\circ} - 19.1^{\circ} = \mathbf{340.9^{\circ}}$.
\item Velocity of $A$ relative to $B$:
\[
{}_A\vec{v}_B = -{}_B\vec{v}_A = (7.5\vec{i} - 21.65\vec{j})\ \mathrm{m\,s^{-1}},
\]
with the same magnitude $22.9\ \mathrm{m\,s^{-1}}$, pointing south-east on a bearing of $180^{\circ} - 19.1^{\circ} = \mathbf{160.9^{\circ}}$.
\end{enumerate}
\end{corrige}

\begin{corrige}[Exercise 7 --- The apparent wind]
Let the true velocity of the wind be $\vec{w} = (a\vec{i} + b\vec{j})\ \mathrm{m\,s^{-1}}$ ($\vec{i}$ east, $\vec{j}$ north), and let $\vec{v}$ be the cyclist's velocity. The apparent wind velocity is $\vec{w} - \vec{v}$.

\textbf{First ride} ($\vec{v} = 6\vec{j}$): the wind \emph{appears to come from the west}, i.e.\ the apparent wind blows towards the east. Hence the apparent velocity is purely eastward:
\[
\vec{w} - 6\vec{j} = a\vec{i} + (b-6)\vec{j} \quad\text{has zero } \vec{j}\text{-component} \quad\Longrightarrow\quad b = 6.
\]

\textbf{Second ride} ($\vec{v} = 12\vec{j}$): the wind appears to come from the north-west, i.e.\ the apparent wind blows towards the south-east, so its east and south components are equal in magnitude:
\[
\vec{w} - 12\vec{j} = a\vec{i} + (6-12)\vec{j} = a\vec{i} - 6\vec{j}.
\]
South-east direction requires $a = 6$ (and $a>0$, consistent).

\textbf{Conclusion.} $\vec{w} = (6\vec{i} + 6\vec{j})\ \mathrm{m\,s^{-1}}$; the true wind speed is
\[
|\vec{w}| = \sqrt{6^2+6^2} = 6\sqrt{2} \approx 8.49\ \mathrm{m\,s^{-1}},
\]
blowing towards the north-east, i.e.\ the wind \emph{comes from the south-west} (bearing $225^{\circ}$ as the direction of origin).
\end{corrige}

\begin{corrige}[Exercise 8 --- Ferry on the Sanaga]
The crossing lasts $5\ \mathrm{min} = 300\ \mathrm{s}$ and the river is $600\ \mathrm{m}$ wide.
\begin{enumerate}
\item \textbf{Speed of the current.} The downstream drift of $150\ \mathrm{m}$ is produced entirely by the current:
\[
v_c = \frac{150}{300} = 0.5\ \mathrm{m\,s^{-1}}.
\]
\item \textbf{Ferry's speed in still water.} Steering perpendicular to the banks, the ferry's own speed provides the whole perpendicular component:
\[
v_f = \frac{600}{300} = 2\ \mathrm{m\,s^{-1}}.
\]
\item \textbf{True speed and direction.} The true velocity is the vector sum of two perpendicular components $2$ and $0.5\ \mathrm{m\,s^{-1}}$:
\[
v_{\text{true}} = \sqrt{2^2 + 0.5^2} = \sqrt{4.25} \approx 2.06\ \mathrm{m\,s^{-1}}.
\]
The angle with the perpendicular to the banks is
\[
\theta = \arctan\!\left(\frac{0.5}{2}\right) = \arctan(0.25) \approx 14.0^{\circ}.
\]
\end{enumerate}
\end{corrige}

\begin{corrige}[Exercise 9 --- Interception at sea (12 marks)]
Take $\vec{i}$ east, $\vec{j}$ north, $t$ in hours after noon.
\begin{enumerate}
\item $\vec{v}_P = 15\sin 60^{\circ}\,\vec{i} + 15\cos 60^{\circ}\,\vec{j} = \left(\dfrac{15\sqrt{3}}{2}\vec{i} + 7.5\vec{j}\right)\ \mathrm{km\,h^{-1}}$, so
\[
\vec{r}_P = t\,\vec{v}_P = \frac{15\sqrt{3}}{2}t\,\vec{i} + 7.5t\,\vec{j}, \qquad
\vec{r}_Q = 20\vec{i} + t\,\vec{v}_Q. \tag{3 marks}
\]
\item Interception at time $T$ means $\vec{r}_P(T) = \vec{r}_Q(T)$:
\[
T\vec{v}_P = 20\vec{i} + T\vec{v}_Q
\quad\Longrightarrow\quad
\vec{v}_Q = \vec{v}_P - \frac{20\vec{i}}{T}.
\]
\emph{Examiner's note:} with $Q$ initially $20\ \mathrm{km}$ \textbf{east} of $O$, the initial displacement of $P$ from $Q$ is $-20\vec{i}$, which gives the minus sign above; the formula $\vec{v}_Q = \vec{v}_P + \dfrac{20\vec{i}}{T}$ printed in the question corresponds to the opposite sign convention (initial displacement taken as $+20\vec{i}$ from $P$ to $Q$). The physics is unchanged. \hfill (3 marks)
\item Write $\vec{v}_Q = \vec{v}_P - \dfrac{20}{T}\vec{i}$ and impose $|\vec{v}_Q| = 18$:
\[
\left(\frac{15\sqrt{3}}{2} - \frac{20}{T}\right)^{2} + 7.5^{2} = 18^{2}.
\]
Since $18^2 - 7.5^2 = 324 - 56.25 = 267.75$, we get
\[
\frac{15\sqrt{3}}{2} - \frac{20}{T} = \pm\sqrt{267.75} \approx \pm 16.36.
\]
The $+$ sign gives $\dfrac{20}{T} = 12.99 - 16.36 < 0$, impossible. Hence
\[
\frac{20}{T} = 12.99 + 16.36 = 29.35 \quad\Longrightarrow\quad T \approx 0.681\ \mathrm{h} \approx 41\ \mathrm{min}.
\]
Then $\vec{v}_Q \approx (-16.36\vec{i} + 7.5\vec{j})\ \mathrm{km\,h^{-1}}$: north-west, at an angle $\arctan\!\left(\dfrac{16.36}{7.5}\right) \approx 65.4^{\circ}$ west of north, i.e.\ on a \textbf{bearing of $294.6^{\circ}$}, with interception at about $\mathbf{12{:}41}$. \hfill (4 marks)
\item The point of interception is at distance
\[
OP = 15\,T \approx 15 \times 0.681 \approx 10.2\ \mathrm{km} \quad \text{from } O. \tag{2 marks}
\]
\end{enumerate}
\textbf{Examiner expectations:} correct resolution of $15\ \mathrm{km\,h^{-1}}$ on bearing $060^{\circ}$ (M1); correct interception condition as a vector equation (M1); use of the speed constraint $|\vec{v}_Q| = 18$ leading to a quadratic-type equation and rejection of the spurious root (A1, A1); final answers quoted with units and sensible rounding.
\end{corrige}

\begin{corrige}[Exercise 10 --- Collision or closest approach (14 marks)]
\begin{enumerate}
\item With constant velocities, $\vec{r}(t) = \vec{r}(0) + t\vec{v}$:
\[
\vec{r}_A = (10 + 300t)\vec{i} + (5 + 100t)\vec{j}, \qquad
\vec{r}_B = (-4 + 200t)\vec{i} + (20 - 100t)\vec{j} \quad (\mathrm{km}). \tag{2 marks}
\]
\item Position of $B$ relative to $A$:
\[
{}_B\vec{r}_A = \vec{r}_B - \vec{r}_A = (-14 - 100t)\vec{i} + (15 - 200t)\vec{j}\ \mathrm{km};
\]
velocity of $B$ relative to $A$:
\[
{}_B\vec{v}_A = \vec{v}_B - \vec{v}_A = (-100\vec{i} - 200\vec{j})\ \mathrm{km\,h^{-1}}. \tag{3 marks}
\]
\item Collision requires ${}_B\vec{r}_A = \vec{0}$, i.e.\ simultaneously
\[
-14 - 100t = 0 \;\Rightarrow\; t = -0.14, \qquad 15 - 200t = 0 \;\Rightarrow\; t = 0.075.
\]
The two conditions give different (and partly negative) times, so \textbf{the aircraft do not collide}. \hfill (3 marks)
\item The squared separation is
\[
D^2(t) = (14 + 100t)^2 + (15 - 200t)^2.
\]
Differentiating:
\[
\frac{d(D^2)}{dt} = 200(14 + 100t) - 400(15 - 200t) = 100\,000\,t - 3200.
\]
Setting this to zero: $t^\ast = \dfrac{3200}{100\,000} = 0.032\ \mathrm{h} = 1.92\ \mathrm{min}$ (a minimum, since $D^2$ is a quadratic with positive leading coefficient). Then
\[
D^2(t^\ast) = (17.2)^2 + (8.6)^2 = 295.84 + 73.96 = 369.8,
\qquad D_{\min} \approx 19.2\ \mathrm{km}. \tag{4 marks}
\]
\item Since $D_{\min} \approx 19.2\ \mathrm{km} > 5\ \mathrm{km}$, the minimum-separation regulation is \textbf{not violated}: the aircraft never come closer than about $19\ \mathrm{km}$. \hfill (2 marks)
\end{enumerate}
\textbf{Examiner expectations:} relative position and relative velocity correctly formed (M1); collision test requiring \emph{simultaneous} vanishing of both components (M1, A1); minimisation of $D^2$ by calculus or by dot product ${}_B\vec{r}_A \cdot {}_B\vec{v}_A = 0$ (M1); conclusion in (e) must compare $D_{\min}$ with $5\ \mathrm{km}$ explicitly (A1).
\end{corrige}

\begin{corrige}[Exercise 11 --- Radio contact between two boats (13 marks)]
\begin{enumerate}
\item From the position vectors, $\vec{v}_M = (6\vec{i} + 8\vec{j})\ \mathrm{km\,h^{-1}}$ and $\vec{v}_N = (-2\vec{i} + 4\vec{j})\ \mathrm{km\,h^{-1}}$. Hence
\[
{}_M\vec{v}_N = \vec{v}_M - \vec{v}_N = (8\vec{i} + 4\vec{j})\ \mathrm{km\,h^{-1}}.
\]
\emph{Interpretation:} seen from boat $N$, boat $M$ appears to move with constant velocity $(8\vec{i} + 4\vec{j})\ \mathrm{km\,h^{-1}}$, i.e.\ at $\sqrt{80} \approx 8.94\ \mathrm{km\,h^{-1}}$ on a bearing of $\arctan(8/4) \approx 063.4^{\circ}$. \hfill (3 marks)
\item The displacement from $M$ to $N$ is
\[
\overrightarrow{MN} = \vec{r}_N - \vec{r}_M = (8 - 8t)\vec{i} + (2 - 4t)\vec{j}.
\]
Therefore
\begin{align*}
MN^2 &= (8-8t)^2 + (2-4t)^2 \\
&= 64 - 128t + 64t^2 + 4 - 16t + 16t^2 \\
&= 80t^2 - 144t + 68.
\end{align*}
\emph{Examiner's note:} direct expansion of the given data yields the coefficient $-144$; the value $-136$ printed in the question is a misprint, and the remainder of the solution uses the correct expression $MN^2 = 80t^2 - 144t + 68$. \hfill (4 marks)
\item Radio contact requires $MN^2 \leq 25$:
\[
80t^2 - 144t + 68 \leq 25 \quad\Longleftrightarrow\quad 80t^2 - 144t + 43 \leq 0.
\]
Discriminant: $\Delta = 144^2 - 4(80)(43) = 20736 - 13760 = 6976 = 64 \times 109$, so $\sqrt{\Delta} = 8\sqrt{109} \approx 83.5$. The roots are
\[
t = \frac{144 \pm 8\sqrt{109}}{160} \quad\Longrightarrow\quad t_1 \approx 0.378\ \mathrm{h}, \quad t_2 \approx 1.422\ \mathrm{h}.
\]
Since the quadratic is negative between its roots, the boats can communicate for
\[
0.378\ \mathrm{h} \lesssim t \lesssim 1.422\ \mathrm{h},
\]
i.e.\ from about $t = 22.7\ \mathrm{min}$ to $t = 85.3\ \mathrm{min}$ (a window of roughly $62.6\ \mathrm{min}$). \hfill (4 marks)
\item $MN^2$ is minimal when $\dfrac{d}{dt}(MN^2) = 160t - 144 = 0$, i.e.\ at
\[
t^\ast = \frac{144}{160} = 0.9\ \mathrm{h} = 54\ \mathrm{min}.
\]
Then $MN^2_{\min} = 80(0.9)^2 - 144(0.9) + 68 = 64.8 - 129.6 + 68 = 3.2$, so
\[
MN_{\min} = \sqrt{3.2} \approx 1.79\ \mathrm{km}. \tag{2 marks}
\]
\end{enumerate}
\textbf{Examiner expectations:} velocities read correctly from $\vec{r}(t)$ (B1); relative velocity as a difference (M1); correct expansion leading to the quadratic (M1, A1); conversion of the range condition into a quadratic inequality solved between the roots (M1); endpoints converted to minutes (A1); minimum found by calculus or by completing the square, with units (A1).
\end{corrige}

\newpage

\coursec{Summary sheet}

\begin{popfigure}
\begin{center}
\begin{tikzpicture}[
  mindmap,
  grow cyclic,
  every node/.style={concept, text=white, font=\small},
  concept color=popblue,
  level 1 concept/.append style={sibling angle=60, level distance=3.4cm, font=\footnotesize},
  level 2 concept/.append style={sibling angle=45, level distance=2.1cm, font=\tiny},
  root concept/.append style={concept color=popdark, font=\small\bfseries}
]
\node[concept] {Velocity \& Acceleration as Vectors}
  child[concept color=poporange] { node {Relative Velocity $\vec{v}_{A/B}=\vec{v}_A-\vec{v}_B$}
    child { node {Passing / overtaking} }
    child { node {Apparent wind} }
  }
  child[concept color=popgreen] { node {Resultant \& True Velocity}
    child { node {River crossing} }
    child { node {Wind on aircraft} }
    child { node {Non-parallel forces} }
  }
  child[concept color=poppurple] { node {Interception \& Collision}
    child { node {$\vec{r}_A(t)=\vec{r}_B(t)$} }
    child { node {Pursuit problems} }
  }
  child[concept color=poppink] { node {Relative Position \& Motion}
    child { node {$\vec{r}_{A/B}=\vec{r}_A-\vec{r}_B$} }
    child { node {Closest approach} }
  }
  child[concept color=popgold] { node {Vector Tools}
    child { node {Components $\vec{v}=v_x\vec{i}+v_y\vec{j}$} }
    child { node {Cosine \& sine rules} }
  }
  child[concept color=popblue!70!black] { node {Applications}
    child { node {Navigation} }
    child { node {Safety \& collision} }
  };
\end{tikzpicture}
\end{center}
\legende{Mind map of the chapter: velocity and acceleration as vectors.}
\end{popfigure}

\begin{bilan}
\textbf{Key definitions.}
\begin{itemize}
\item \cle{Velocity vector}: $\vec{v}=\dfrac{d\vec{r}}{dt}=\dot{x}\,\vec{i}+\dot{y}\,\vec{j}$; speed $=|\vec{v}|$.
\item \cle{Acceleration vector}: $\vec{a}=\dfrac{d\vec{v}}{dt}=\ddot{x}\,\vec{i}+\ddot{y}\,\vec{j}$.
\item \cle{Relative velocity} of $A$ with respect to $B$: $\vec{v}_{A/B}=\vec{v}_A-\vec{v}_B$. Hence $\vec{v}_{A/B}=-\vec{v}_{B/A}$.
\item \cle{Relative position} of $A$ with respect to $B$: $\vec{r}_{A/B}=\vec{r}_A-\vec{r}_B$, and $\vec{v}_{A/B}=\dfrac{d}{dt}\vec{r}_{A/B}$.
\item \cle{True velocity} (river/wind problems): vector sum of the object's own velocity and the velocity of the medium: $\vec{v}_{\text{true}}=\vec{v}_{\text{object}}+\vec{v}_{\text{medium}}$.
\item \cle{Interception / collision}: two moving bodies meet when their position vectors are equal at the same instant: $\vec{r}_A(t)=\vec{r}_B(t)$.
\end{itemize}

\textbf{Formulas and laws to know.}
\begin{itemize}
\item Components: if $\vec{v}$ has magnitude $v$ and direction $\theta$, then $\vec{v}=v\cos\theta\,\vec{i}+v\sin\theta\,\vec{j}$; $|\vec{v}|=\sqrt{v_x^2+v_y^2}$, $\tan\theta=\dfrac{v_y}{v_x}$.
\item Resultant of non-parallel velocities/forces: resolve into components, add, recombine; or use the cosine rule $R^2=P^2+Q^2+2PQ\cos\alpha$ and sine rule on the vector triangle.
\item River crossing: time to cross $t=\dfrac{\text{width}}{v_{\text{perpendicular}}}$; drift downstream $=v_{\text{stream}}\times t$. To cross straight across, aim upstream at angle $\theta$ with $\sin\theta=\dfrac{v_{\text{stream}}}{v_{\text{boat}}}$.
\item Interception: write $\vec{r}_A(t)=\vec{r}_{A0}+\vec{v}_A t$, $\vec{r}_B(t)=\vec{r}_{B0}+\vec{v}_B t$; solve $\vec{r}_A(t)=\vec{r}_B(t)$ component by component (two equations, unknowns $t$ and possibly a speed or heading).
\item Closest approach: distance $d(t)=|\vec{r}_{A/B}(t)|$; minimise $d^2(t)$ (easier: differentiate $d^2$), or set $\dfrac{d}{dt}|\vec{r}_{A/B}|^2=0$, i.e.\ $\vec{r}_{A/B}\cdot\vec{v}_{A/B}=0$.
\end{itemize}

\textbf{Methods.}
\begin{itemize}
\item Always start with a clear vector diagram: fix axes, mark given vectors to scale, draw the triangle/parallelogram.
\item For relative velocity: write each velocity in components, subtract, then find magnitude and bearing of the result.
\item For interception: equate $\vec{i}$-components and $\vec{j}$-components separately; check that $t>0$.
\item For closest approach: form $\vec{r}_{A/B}(t)$, compute $d^2(t)$, differentiate, solve $d'(t)=0$, verify it is a minimum.
\end{itemize}

\textbf{Common mistakes.}
\begin{itemize}
\item Confusing $\vec{v}_{A/B}$ with $\vec{v}_{B/A}$: remember ``$A$ relative to $B$'' means $\vec{v}_A-\vec{v}_B$.
\item Adding speeds as scalars when velocities are not parallel — always add \emph{vectors}.
\item Forgetting that interception requires \emph{simultaneous} equality of both coordinates, not just equal distances travelled.
\item Mixing units (km/h with m/s) or bearings with standard angles; convert before computing.
\item In river problems, using the boat's heading speed instead of the perpendicular component for the crossing time.
\item Sign errors when subtracting components; write components explicitly with signs.
\end{itemize}

\textbf{GCE tips.}
\begin{itemize}
\item State your axes and unit vectors $\vec{i},\vec{j}$ at the start; examiners award marks for a correct vector setup.
\item Give final directions as bearings (three figures, e.g.\ $037^{\circ}$) or angles measured from a stated direction, and magnitudes to 3 s.f.\ with units.
\item Show the vector triangle when using the cosine/sine rule — a labelled diagram earns method marks even if arithmetic slips.
\item For ``show that'' questions on interception, keep $t$ symbolic until the last line.
\item Always check your answer is physically sensible: positive time, speed not exceeding the sum of the given speeds.
\end{itemize}
\end{bilan}

\findecours

\end{document}
